\documentclass[12pt]{book}
\usepackage[T1]{fontenc}
\usepackage[utf8]{inputenc}
\usepackage[brazil]{babel}
\usepackage{times}
\usepackage{listings}
\usepackage{graphicx}
\usepackage{xcolor}
\usepackage{url}
% \usepackage{showidx}
\usepackage{imakeidx}
\usepackage{booktabs}
\usepackage{url}
\usepackage{etoolbox}  % for showidx
\usepackage{fullpage}
\usepackage{soul}

\lstset{
    literate=  % Define substituições para caracteres especiais
    {á}{{\'a}}1 {é}{{\'e}}1 {í}{{\'i}}1 {ó}{{\'o}}1 {ú}{{\'u}}1
    {Á}{{\'A}}1 {É}{{\'E}}1 {Í}{{\'I}}1 {Ó}{{\'O}}1 {Ú}{{\'U}}1
    {ã}{{\~a}}1 {õ}{{\~o}}1 {Ã}{{\~A}}1 {Õ}{{\~O}}1
    {ç}{{\c c}}1 {Ç}{{\c C}}1,
    basicstyle=\ttfamily\small, % Estilo básico do código
    commentstyle=\color{gray},   % Estilo dos comentários
}

\usepackage{hyperref}  % should be last

% One space after periods
\frenchspacing

\hypersetup{pdfauthor={Russ Cox, Frans Kaashoek, Robert Morris},
            pdftitle={xv6: a simple, Unix-like teaching operating system},}
  
\lstset{basicstyle=\small\ttfamily}
\lstset{morecomment=[is]{[[[}{]]]}}
\lstset{escapeinside={(*@}{@*)}}
\lstset{xleftmargin=5.0ex}

\newcommand{\github}{https://github.com/mit-pdos/xv6-riscv/blob/riscv/}

\newcommand{\fileref}[1]{\href{\github/#1}{\small{(#1)}}}
\newcommand{\lineref}[2]{\href{\github/#1\#L#2}{\small{(#1:#2)}}}
\newcommand{\linerefs}[3]{\href{\github/#1\#L#2-L#3}{\small(#1:#2-#3)}}

\newcommand{\indextext}[1]{\textit{#1}\index{#1}}
\newcommand{\indextextx}[1]{{#1}\index{#1}}
\newcommand{\indexcode}[1]{\lstinline{#1}\index{#1@\lstinline{#1}}}

%% editing markup
\newcommand{\insertnote}[3]{\noindent\textcolor{#1}{\textbf{#2:} #3}}
\newcommand{\note}[1]{\insertnote{blue}{NOTE}{#1}}
\newcommand{\rtm}[1]{\insertnote{red}{RTM}{#1}}
\newcommand{\mfk}[1]{\insertnote{red}{MFK}{#1}}
%% for publishing book without notes
%\renewcommand{\insertnote}[3]{}

\title{\textbf{xv6: um sistema operacional simples, tipo Unix, para ensino}}
\author{Russ Cox \and Frans Kaashoek \and Robert Morris \and Alan Veloso (Tradutor)}

\makeindex

\begin{document}

\maketitle

\tableofcontents

\chapter*{Prefácio e Agradecimentos}

Este é um texto preliminar destinado a uma disciplina sobre sistemas operacionais.  
Ele explica os principais conceitos de sistemas operacionais por meio do estudo de um kernel de exemplo, chamado xv6.  
O xv6 é inspirado no Unix Version 6 (v6) de Dennis Ritchie e Ken Thompson~\cite{unix}.  
O xv6 segue de forma livre a estrutura e o estilo do v6, mas é implementado em ANSI C~\cite{kernighan} para uma arquitetura RISC-V multicore~\cite{riscv}.

Este texto deve ser lido em conjunto com o código-fonte do xv6,  
uma abordagem inspirada no comentário de John Lions sobre o UNIX 6ª Edição~\cite{lions};  
o texto possui hiperlinks para o código-fonte em  
\url{https://github.com/mit-pdos/xv6-riscv}.  
Consulte também  
\url{https://pdos.csail.mit.edu/6.1810}  
para recursos online adicionais sobre o v6 e o xv6, incluindo diversos exercícios práticos utilizando o xv6.

Utilizamos este texto nos cursos 6.828 e 6.1810, disciplinas de sistemas operacionais no MIT.  
Agradecemos ao corpo docente, assistentes de ensino e estudantes dessas turmas, que direta ou indiretamente contribuíram com o xv6.  
Em particular, gostaríamos de agradecer a Adam Belay, Austin Clements e Nickolai Zeldovich.  
Finalmente, agradecemos às pessoas que nos enviaram relatórios de erros no texto ou sugestões de melhorias:  
Abutalib Aghayev, Sebastian Boehm, brandb97, Anton Burtsev, Raphael Carvalho, Tej Chajed, Brendan Davidson, Rasit Eskicioglu, Color Fuzzy, Wojciech Gac, Giuseppe, Tao Guo, Haibo Hao, Naoki Hayama, Chris Henderson, Robert Hilderman, Eden Hochbaum, Wolfgang Keller, Paweł Kraszewski, Henry Laih, Jin Li, Austin Liew, lyazj@github.com, Pavan Maddamsetti, Jacek Masiulaniec, Michael McConville, m3hm00d, miguelgvieira, Mark Morrissey, Muhammed Mourad, Harry Pan, Harry Porter, Siyuan Qian, Zhefeng Qiao, Askar Safin, Salman Shah, Huang Sha, Vikram Shenoy, Adeodato Simó, Ruslan Savchenko, Pawel Szczurko, Warren Toomey, tyfkda, tzerbib, Vanush Vaswani, Xi Wang, Zou Chang Wei, Sam Whitlock, Qiongsi Wu, LucyShawYang, ykf1114@gmail.com e Meng Zhou.

Se você encontrar erros ou tiver sugestões de melhorias, por favor envie um e-mail para  
Frans Kaashoek e Robert Morris (kaashoek, rtm@csail.mit.edu).

\chapter{Interfaces de sistemas operacionais}
\label{CH:UNIX}

A função de um sistema operacional é compartilhar um computador entre
múltiplos programas e fornecer um conjunto de serviços mais úteis
do que o suporte oferecido apenas pelo hardware.
Um sistema operacional gerencia e abstrai
o hardware de baixo nível, de modo que, por exemplo,
um processador de texto não precise se preocupar com qual tipo
de disco está sendo utilizado.
Um sistema operacional compartilha o hardware entre vários programas para que
eles executem (ou aparentem executar) ao mesmo tempo.
Por fim, sistemas operacionais fornecem maneiras controladas para que os programas
interajam, de modo que possam compartilhar dados ou trabalhar em conjunto.

Um sistema operacional fornece serviços aos programas de usuário por meio de uma interface.
\index{projeto de interface}
Projetar uma boa interface acaba sendo
uma tarefa difícil. Por um lado, gostaríamos que a interface fosse
simples e enxuta, pois isso facilita a implementação correta.
Por outro lado,
podemos ser tentados a oferecer muitos recursos sofisticados às aplicações.
O truque para resolver essa tensão é projetar interfaces que se baseiem em poucos
mecanismos que possam ser combinados para oferecer grande generalidade.

Este livro utiliza um único sistema operacional como exemplo concreto para
ilustrar conceitos de sistemas operacionais. Esse sistema operacional,
o xv6, fornece as interfaces básicas introduzidas por Ken Thompson e
Dennis Ritchie no sistema Unix~\cite{unix}, além de imitar o projeto interno do Unix.
O Unix fornece uma
interface enxuta cujos mecanismos se combinam bem, oferecendo um surpreendente
grau de generalidade. Essa interface foi tão bem-sucedida que
sistemas operacionais modernos—BSD, Linux, macOS, Solaris e até, em menor grau,
o Microsoft Windows—possuem interfaces semelhantes ao Unix.
Compreender o xv6 é um bom começo para entender qualquer um desses
sistemas e muitos outros.

Como mostra a 
Figura~\ref{fig:os},
o xv6 assume a forma tradicional de um
\indextext{kernel} (núcleo),
um programa especial que fornece
serviços aos programas em execução.
Cada programa em execução, chamado de
\indextext{processo},
possui memória contendo instruções, dados e uma pilha. As
instruções implementam o
cálculo do programa. Os dados são as variáveis sobre as quais
o cálculo atua. A pilha organiza as chamadas de procedimentos do programa.
Um computador geralmente possui muitos processos, mas apenas um único
kernel.

Quando um
processo precisa invocar um serviço do kernel, ele chama
uma \indextext{chamada de sistema} (system call),
uma das chamadas
da interface do sistema operacional.
A chamada de sistema entra no kernel;
o kernel realiza o serviço e retorna.
Assim, um processo alterna entre a execução no
\indextext{espaço do usuário}
e no
\indextext{espaço do kernel}.

Como será descrito em detalhes nos capítulos seguintes, o kernel utiliza os mecanismos de proteção de hardware fornecidos por uma
CPU\footnote{
Este texto geralmente se refere ao elemento de hardware que executa uma
computação com o termo \indextext{CPU}, sigla para unidade central
de processamento (central processing unit). Outras documentações (por exemplo, a especificação RISC-V)
também utilizam os termos processador, núcleo (core) e hart em vez de CPU.
}
para
garantir que cada processo executando no espaço do usuário possa acessar apenas
sua própria memória.
O kernel executa com os privilégios de hardware necessários para
implementar essas proteções; os programas de usuário executam sem
esses privilégios.
Quando um programa de usuário invoca uma chamada de sistema, o hardware
eleva o nível de privilégio e começa a executar uma função pré-configurada
no kernel.

\begin{figure}[t]
\center
\includegraphics[scale=0.5]{fig/os.pdf}
\caption{Um kernel e dois processos de usuário.}
\label{fig:os}
\end{figure}

O conjunto de chamadas de sistema que um kernel fornece
é a interface que os programas de usuário enxergam.
O kernel do xv6 fornece um subconjunto dos serviços e chamadas de sistema
que os kernels Unix tradicionalmente oferecem.  
A Figura~\ref{fig:api} 
lista todas as chamadas de sistema do xv6.

O restante deste capítulo descreve os serviços do xv6—processos, memória,
descritores de arquivos, pipes e sistema de arquivos—e os ilustra com
trechos de código e discussões sobre como o \indextext{shell},
a interface de linha de comando do Unix, os utiliza.
O uso das chamadas de sistema pelo shell ilustra o quanto essas interfaces
foram cuidadosamente projetadas.

O shell é um programa comum que lê comandos do usuário
e os executa.
O fato de o shell ser um programa de usuário, e não parte do kernel,
ilustra o poder da interface de chamadas de sistema: não há nada de
especial sobre o shell.
Isso também significa que o shell pode ser facilmente substituído; como resultado,
sistemas Unix modernos possuem uma variedade de
shells disponíveis, cada um com sua própria interface
e recursos de script.
O shell do xv6 é uma implementação simples da essência do
shell Bourne do Unix. Sua implementação pode ser encontrada em 
\lineref{user/sh.c}{1}.

%% 
%% 	Processes and memory
%% 
\section{Processos e memória}

Um processo no xv6 consiste em memória no espaço do usuário (instruções, dados e pilha)
e estado por processo, privado do kernel.
O xv6
\indextext{compartilha o tempo} (time-share)
entre processos: ele alterna transparentemente as CPUs disponíveis
entre o conjunto de processos que aguardam execução.
Quando um processo não está em execução, o xv6 salva os registradores da CPU do processo,
restaurando-os quando for executá-lo novamente.
O kernel associa um identificador de processo, ou
\indexcode{PID},
a cada processo.

\begin{figure}[t]
\center
\begin{tabular}{ll}
\textbf{Chamada de sistema}            & \textbf{Descrição}                                                                                                                           \\ \hline
int fork()                             & Cria um processo, retorna o PID do filho.                                                                                                    \\
int exit(int status)                   & Termina o processo atual; status retornado para wait(). Sem retorno.                                                                         \\
int wait(int *status)                  & \begin{tabular}[c]{@{}l@{}}Aguarda a finalização de um processo filho; status é colocado \\ em *status; retorna o PID do filho.\end{tabular} \\
int kill(int pid)                      & Encerra o processo PID. Retorna 0, ou -1 em caso de erro.                                                                                    \\
int getpid()                           & Retorna o PID do processo atual.                                                                                                             \\
int sleep(int n)                       & Pausa por n ciclos de clock.                                                                                                                 \\
int exec(char *file, char *argv{[}{]}) & \begin{tabular}[c]{@{}l@{}}Carrega um arquivo e executa com argumentos; só retorna em caso \\ de erro.\end{tabular}                          \\
char *sbrk(int n)                      & \begin{tabular}[c]{@{}l@{}}Expande a memória do processo em n bytes zero. Retorna o início da \\ nova memória.\end{tabular}                  \\
int open(char *file, int flags)        & \begin{tabular}[c]{@{}l@{}}Abre um arquivo; flags indicam leitura/escrita; retorna um fd \\ (descritor de arquivo).\end{tabular}             \\
int write(int fd, char *buf, int n)    & Escreve n bytes de buf no descritor fd; retorna n.                                                                                           \\
int read(int fd, char *buf, int n)     & Lê n bytes em buf; retorna quantidade lida, ou 0 se fim de arquivo.                                                                          \\
int close(int fd)                      & Libera o arquivo aberto fd.                                                                                                                  \\
int dup(int fd)                        & Retorna um novo descritor que referencia o mesmo arquivo que fd.                                                                             \\
int pipe(int p{[}{]})                  & Cria um pipe, insere descritores de leitura/escrita em p{[}0{]} e p{[}1{]}.                                                                  \\
int chdir(char *dir)                   & Altera o diretório atual.                                                                                                                    \\
int mkdir(char *dir)                   & Cria um novo diretório.                                                                                                                      \\
int mknod(char *file, int, int)        & Cria um arquivo de dispositivo.                                                                                                              \\
int fstat(int fd, struct stat *st)     & Insere informações de um arquivo aberto em *st.                                                                                              \\
int link(char *file1, char *file2)     & Cria outro nome (file2) para o arquivo file1.                                                                                                \\
int unlink(char *file)                 & Remove um arquivo.                                                                                                                          
\end{tabular}
\caption{Chamadas de sistema do xv6. Salvo indicação contrária, essas chamadas retornam 0 se não houver erro, e -1 em caso de erro.}
\label{fig:api}
\end{figure}

Um processo pode criar um novo processo usando a chamada de sistema
\indexcode{fork}.
\lstinline{fork}
cria no novo processo uma cópia exata da memória do processo que fez a chamada: ele copia as instruções, dados e
a pilha do processo original para o novo processo.
\lstinline{fork}
retorna em ambos os processos: o original e o novo.
No processo original, \lstinline{fork} retorna o PID do novo processo.
No processo novo, \lstinline{fork} retorna zero.
Os processos original e novo são frequentemente chamados de
\indextext{pai}
e
\indextext{filho}.

Por exemplo, considere o seguinte trecho de programa escrito na linguagem C~\cite{kernighan}:
\begin{lstlisting}[]
int pid = fork();
if(pid > 0){
  printf("pai: filho=%d\n", pid);
  pid = wait((int *) 0);
  printf("filho %d terminou\n", pid);
} else if(pid == 0){
  printf("filho: saindo\n");
  exit(0);
} else {
  printf("erro no fork\n");
}
\end{lstlisting}
A chamada de sistema
\indexcode{exit}
faz com que o processo que a chamou pare de executar e
libere recursos como memória e arquivos abertos.
Exit recebe um argumento inteiro de status,
convencionalmente 0 para indicar sucesso e 1 para indicar falha.
A chamada
\indexcode{wait}
retorna o PID de um filho finalizado (ou morto) do processo atual
e copia o status de saída do filho para o endereço passado para wait;
se nenhum dos filhos do chamador tiver finalizado,
\indexcode{wait}
aguarda que um finalize.
Se o processo pai não tiver filhos, \lstinline{wait} retorna imediatamente -1.
Se o pai não se importa com o status de saída do filho, pode
passar o endereço nulo (0) para
\lstinline{wait}.

No exemplo, as linhas de saída
\begin{lstlisting}[]
pai: filho=1234
filho: saindo
\end{lstlisting}
podem aparecer em qualquer ordem (ou até intercaladas), dependendo de quem
executar o \indexcode{printf} primeiro: o pai ou o filho.
Após o filho finalizar, o \indexcode{wait} do pai retorna, fazendo com que o pai imprima:
\begin{lstlisting}[]
pai: filho 1234 terminou
\end{lstlisting}
Embora o filho comece com o mesmo conteúdo de memória do pai,
pai e filho executam com memórias e registradores separados:
alterar uma variável em um não afeta o outro. Por exemplo, quando o valor de retorno de
\lstinline{wait}
é armazenado em
\lstinline{pid}
no processo pai,
isso não altera a variável
\lstinline{pid}
no filho. O valor de
\lstinline{pid}
no filho continuará sendo zero.

A chamada de sistema
\indexcode{exec}
substitui a memória do processo chamador por uma nova imagem
de memória carregada de um arquivo no sistema de arquivos.
O arquivo deve ter um formato específico, que indica quais partes contêm instruções,
dados, onde iniciar a execução, etc.
O xv6 usa o formato ELF, que será discutido com mais detalhes no Capítulo~\ref{CH:MEM}.
Normalmente, o arquivo é o resultado da compilação de um código-fonte.
Quando
\indexcode{exec}
é bem-sucedido, ele não retorna ao programa chamador;
em vez disso, as instruções carregadas começam a ser executadas a partir do ponto de entrada
indicado no cabeçalho ELF.
\lstinline{exec}
recebe dois argumentos: o nome do arquivo executável e um vetor de argumentos em formato string.
Por exemplo:
\begin{lstlisting}[]
char *argv[3];

argv[0] = "echo";
argv[1] = "hello";
argv[2] = 0;
exec("/bin/echo", argv);
printf("erro no exec\n");
\end{lstlisting}
Esse trecho substitui o programa chamador por uma instância
do programa 
\lstinline{/bin/echo}
executando com a lista de argumentos
\lstinline{echo hello}.
A maioria dos programas ignora o primeiro elemento do vetor de argumentos, que
convencionalmente é o nome do programa.

O shell do xv6 usa essas chamadas para executar programas em nome dos
usuários. A estrutura principal do shell é simples; veja
\lstinline{main}.
O loop principal lê uma linha de entrada do usuário com
\indexcode{getcmd}.
Depois, chama
\lstinline{fork},
que cria uma cópia do processo do shell.
O processo pai chama
\lstinline{wait},
enquanto o filho executa o comando. Por exemplo, se o usuário digitar
``\lstinline{echo hello}''
no shell,
\lstinline{runcmd}
será chamado com
``\lstinline{echo hello}''
como argumento.
\lstinline{runcmd}
executa o comando propriamente dito. Para
``\lstinline{echo hello}'',
ele chamaria
\lstinline{exec}.
Se
\lstinline{exec}
for bem-sucedido, o filho passará a executar as instruções de
\lstinline{echo}
em vez de
\lstinline{runcmd}.
Em algum momento,
\lstinline{echo}
chamará
\lstinline{exit},
fazendo com que o pai retorne de
\lstinline{wait}
em
\lstinline{main}.

Você pode se perguntar por que
\indexcode{fork}
e
\indexcode{exec}
não são combinadas em uma única chamada; veremos mais adiante que
o shell se aproveita dessa separação para implementar
redirecionamento de entrada/saída.
Para evitar o desperdício de criar um processo duplicado apenas para substituí-lo
imediatamente (com \lstinline{exec}),
os kernels modernos otimizam a implementação de
\lstinline{fork}
usando técnicas de memória virtual, como
copy-on-write (ver Seção~\ref{sec:pagefaults}).

O xv6 aloca a maior parte da memória do espaço do usuário
implicitamente:
\indexcode{fork}
aloca a memória necessária para a cópia da memória do pai no filho, e
\indexcode{exec}
aloca memória suficiente para carregar o executável.
Um processo que precisa de mais memória em tempo de execução (por exemplo, para
\indexcode{malloc})
pode chamar
\lstinline{sbrk(n)}
para expandir sua memória de dados em
\lstinline{n} bytes zero;
\indexcode{sbrk}
retorna o endereço da nova memória.


%% 
%% 	I/O and File descriptors
%% 
\section{E/S e Descritores de Arquivo}

Um 
\indextext{descritor de arquivo}
é um pequeno número inteiro que representa um objeto gerenciado pelo kernel
do qual um processo pode ler ou para o qual pode escrever.
Um processo pode obter um descritor de arquivo ao abrir um arquivo, diretório
ou dispositivo, ou ao criar um pipe, ou duplicando um descritor
existente.
Para simplificar, frequentemente nos referiremos ao objeto associado ao descritor de arquivo
simplesmente como um ``arquivo'';  
a interface de descritor de arquivo abstrai as diferenças entre
arquivos, pipes e dispositivos, fazendo com que todos pareçam fluxos de bytes.
Chamaremos leitura e escrita de \indextext{E/S} (entrada/saída).

Internamente, o kernel do xv6 usa o descritor de arquivo
como um índice em uma tabela por processo,
de modo que cada processo possui um espaço privado de descritores de arquivo,
começando pelo zero.
Por convenção, um processo lê do descritor 0 (entrada padrão),
escreve na saída padrão pelo descritor 1 e
escreve mensagens de erro no descritor 2 (erro padrão).
Como veremos, o shell aproveita essa convenção para implementar redirecionamento de E/S
e pipelines. O shell garante que sempre tenha três descritores de arquivo
abertos,
que, por padrão, estão conectados ao console.

As chamadas de sistema
\lstinline{read}
e
\lstinline{write}
leem bytes de e escrevem bytes em
arquivos abertos identificados por descritores de arquivo.
A chamada
\lstinline{read(fd, buf, n)}
lê no máximo
\lstinline{n}
bytes do descritor de arquivo
\lstinline{fd},
copia-os para
\lstinline{buf},
e retorna o número de bytes lidos.
Cada descritor de arquivo associado a um arquivo
possui um deslocamento associado.
\lstinline{read}
lê os dados a partir do deslocamento atual do arquivo e então avança
esse deslocamento pelo número de bytes lidos:
uma chamada subsequente a
\lstinline{read}
retornará os bytes seguintes aos da primeira chamada.
Quando não há mais bytes a serem lidos,
\lstinline{read}
retorna zero, indicando fim de arquivo.

A chamada
\lstinline{write(fd, buf, n)}
escreve
\lstinline{n}
bytes de
\lstinline{buf}
no descritor de arquivo
\lstinline{fd}
e retorna o número de bytes escritos.
Menos de
\lstinline{n}
bytes serão escritos apenas em caso de erro.
Assim como
\lstinline{read},
\lstinline{write}
escreve os dados no deslocamento atual e depois o avança
de acordo com o número de bytes escritos:
cada
\lstinline{write}
continua a partir de onde a anterior parou.

O trecho de programa a seguir (que representa a essência do programa
\lstinline{cat})
copia dados da entrada padrão
para a saída padrão. Se ocorrer algum erro, ele escreve uma mensagem
no erro padrão.
\begin{lstlisting}[]
char buf[512];
int n;

for(;;){
  n = read(0, buf, sizeof buf);
  if(n == 0)
    break;
  if(n < 0){
    fprintf(2, "erro de leitura\n");
    exit(1);
  }
  if(write(1, buf, n) != n){
    fprintf(2, "erro de escrita\n");
    exit(1);
  }
}
\end{lstlisting}
O importante a se notar nesse trecho é que o programa
\lstinline{cat}
não sabe se está lendo de um arquivo, do console ou de um pipe.
Da mesma forma,
\lstinline{cat}
não sabe se está escrevendo no console, em um arquivo, ou em outro lugar.
O uso de descritores de arquivo e a convenção de que o descritor 0
é a entrada e o descritor 1 é a saída permite uma implementação
simples do
\lstinline{cat}.

A chamada de sistema
\lstinline{close}
libera um descritor de arquivo, tornando-o disponível para reutilização por uma futura
chamada
\lstinline{open},
\lstinline{pipe}
ou
\lstinline{dup}
(ver abaixo).
Um novo descritor de arquivo alocado
é sempre o menor descritor não utilizado
pelo processo atual.

Descritores de arquivo e a chamada
\indexcode{fork}
interagem para tornar fácil a implementação de redirecionamento de E/S.
\lstinline{fork}
copia a tabela de descritores de arquivo do processo pai junto com sua memória,
de modo que o filho inicia com exatamente os mesmos arquivos abertos que o pai.
A chamada de sistema
\indexcode{exec}
substitui a memória do processo chamador, mas preserva sua tabela de arquivos.
Esse comportamento permite que o shell
implemente \indextext{redirecionamento de E/S} fazendo fork,
reabrindo os descritores de arquivo desejados no filho
e, então, chamando \lstinline{exec} para executar o novo programa.
Aqui está uma versão simplificada do código que um shell executa para o comando
\lstinline{cat < input.txt}:
\begin{lstlisting}[]
char *argv[2];

argv[0] = "cat";
argv[1] = 0;
if(fork() == 0) {
  close(0);
  open("input.txt", O_RDONLY);
  exec("cat", argv);
}
\end{lstlisting}
Após o processo filho fechar o descritor de arquivo 0,
\lstinline{open}
garantidamente usará esse descritor
para abrir
\lstinline{input.txt}:
0 será o menor descritor disponível.
O programa
\lstinline{cat}
executará com o descritor 0 (entrada padrão) referindo-se a
\lstinline{input.txt}.
Os descritores do processo pai não são afetados por essa
sequência, pois ela modifica apenas os descritores do filho.

O código de redirecionamento de E/S no shell do xv6 funciona exatamente dessa forma.
Lembre-se de que, nesse ponto do código, o shell já fez fork do
processo filho, e
\lstinline{runcmd}
chamará
\lstinline{exec}
para carregar o novo programa.

O segundo argumento de \lstinline{open} consiste em um conjunto de
flags, expressos como bits, que controlam o comportamento da chamada \lstinline{open}.
Os valores possíveis são definidos no cabeçalho de controle de arquivos (fcntl):
\lstinline{O_RDONLY},
\lstinline{O_WRONLY},
\lstinline{O_RDWR},
\lstinline{O_CREATE}, e
\lstinline{O_TRUNC},
que instruem \lstinline{open} a
abrir o arquivo para leitura,
ou para escrita,
ou para leitura e escrita,
a criar o arquivo se ele não existir,
e a truncar (zerar) o arquivo.

Agora fica claro por que é útil que
\lstinline{fork}
e
\lstinline{exec}
sejam chamadas separadas: entre elas, o shell tem a chance
de redirecionar a E/S do processo filho sem afetar a configuração de E/S do shell principal.
Pode-se imaginar uma hipotética chamada combinada
\lstinline{forkexec},
mas as opções para redirecionamento de E/S com essa chamada
seriam complicadas.
O shell teria que modificar sua própria E/S antes de chamar \lstinline{forkexec} (e então
desfazer as modificações); ou
\lstinline{forkexec} teria que aceitar instruções de redirecionamento como argumentos;
ou (menos desejável) cada programa como \lstinline{cat}
teria que ser modificado para fazer seu próprio redirecionamento.

Embora
\lstinline{fork}
copie a tabela de descritores de arquivos, o deslocamento interno de cada arquivo é compartilhado
entre pai e filho.
Considere este exemplo:
\begin{lstlisting}[]
if(fork() == 0) {
  write(1, "hello ", 6);
  exit(0);
} else {
  wait(0);
  write(1, "world\n", 6);
}
\end{lstlisting}
Ao final desse trecho, o arquivo associado ao descritor 1
conterá os dados
\lstinline{hello}
\lstinline{world}.
A chamada
\lstinline{write}
no processo pai
(que, graças ao
\lstinline{wait},
executa apenas após o filho terminar)
continua a partir do ponto onde a chamada
\lstinline{write}
do filho parou.
Esse comportamento ajuda a produzir saídas sequenciais de comandos do shell como:
\lstinline{(echo hello; echo world) >output.txt}.

A chamada
\lstinline{dup}
duplica um descritor de arquivo existente,
retornando um novo que se refere ao mesmo objeto de E/S.
Ambos os descritores compartilham o mesmo deslocamento, assim como os descritores
duplicados por
\lstinline{fork}.
Esta é outra maneira de escrever
\lstinline{hello world}
em um arquivo:
\begin{lstlisting}[]
fd = dup(1);
write(1, "hello ", 6);
write(fd, "world\n", 6);
\end{lstlisting}

Dois descritores compartilham o deslocamento se derivarem do
mesmo descritor original por meio de chamadas de
\lstinline{fork}
e
\lstinline{dup}.
Caso contrário, os descritores não compartilham deslocamento, mesmo que resultem de
chamadas
\lstinline{open}
para o mesmo arquivo.
\lstinline{dup}
permite que shells implementem comandos como:
\lstinline{ls existing-file non-existing-file >tmp1 2>&1}.
O trecho
\lstinline{2>&1}
informa ao shell que o descritor 2 deve ser uma duplicata do descritor 1.
Tanto o nome do arquivo existente quanto a mensagem de erro do arquivo
inexistente aparecerão no arquivo
\lstinline{tmp1}.
O shell do xv6 não suporta redirecionamento de E/S para o descritor de erro,
mas agora você sabe como implementá-lo.

Descritores de arquivos são uma abstração poderosa,
pois escondem os detalhes do que está conectado a eles:
um processo escrevendo no descritor 1 pode estar escrevendo em um
arquivo, em um dispositivo como o console ou em um pipe.

%% 
%% 	Pipes
%% 
\section{Pipes}

Um 
\indextext{pipe}
é um pequeno buffer no kernel exposto aos processos como um par de
descritores de arquivo: um para leitura e outro para escrita.
Escrever dados em uma extremidade do pipe
torna esses dados disponíveis para leitura na outra extremidade.
Pipes fornecem um meio de comunicação entre processos.

O código de exemplo a seguir executa o programa
\lstinline{wc}
com a entrada padrão conectada à extremidade de leitura de um pipe.
\begin{lstlisting}[]
int p[2];
char *argv[2];

argv[0] = "wc";
argv[1] = 0;

pipe(p);
if(fork() == 0) {
  close(0);
  dup(p[0]);
  close(p[0]);
  close(p[1]);
  exec("/bin/wc", argv);
} else {
  close(p[0]);
  write(p[1], "hello world\n", 12);
  close(p[1]);
}
\end{lstlisting}

O programa chama
\lstinline{pipe},
que cria um novo pipe e armazena os descritores de leitura e escrita
no array
\lstinline{p}.
Após o
\lstinline{fork},
tanto o processo pai quanto o processo filho possuem descritores de arquivo
referentes ao pipe.
O processo filho chama \lstinline{close} e \lstinline{dup} para fazer com que o descritor 0
passe a se referir à extremidade de leitura do pipe,
fecha os descritores do array
\lstinline{p},
e chama \lstinline{exec} para executar
\lstinline{wc}.
Quando
\lstinline{wc}
lê da entrada padrão, ele estará lendo do pipe.
O processo pai fecha o lado de leitura do pipe,
escreve no pipe,
e então fecha o lado de escrita.

Se não houver dados disponíveis, uma chamada
\lstinline{read}
em um pipe espera até que dados sejam escritos ou até que todos os
descritores referentes à extremidade de escrita sejam fechados;
neste último caso,
\lstinline{read}
retorna zero, assim como quando se chega ao final de um arquivo.
O fato de que
\lstinline{read}
bloqueia até que seja impossível novos dados chegarem
é uma das razões pelas quais é importante que o processo filho
feche a extremidade de escrita do pipe
antes de executar
\lstinline{wc}
no exemplo acima: se algum dos descritores do
\lstinline{wc}
ainda referisse a extremidade de escrita,
\lstinline{wc}
nunca veria o fim de arquivo.

O shell do xv6 implementa pipelines como
\lstinline{grep fork sh.c | wc -l}
de maneira semelhante ao código acima.
O processo filho cria um pipe para conectar a extremidade esquerda do pipeline
com a extremidade direita. Em seguida, ele chama
\lstinline{fork}
e
\lstinline{runcmd}
para a extremidade esquerda do pipeline,
e depois chama
\lstinline{fork}
e
\lstinline{runcmd}
para a extremidade direita, e espera que ambos terminem.
A extremidade direita do pipeline pode ser um comando que por sua vez inclua outro
pipe (por exemplo,
\lstinline{a | b | c}),
o que leva à criação de dois novos processos filhos (um para
\lstinline{b}
e outro para
\lstinline{c}).
Assim, o shell pode criar uma árvore de processos. As folhas dessa árvore são os comandos,
e os nós internos são processos que esperam até que seus filhos à esquerda e à direita terminem.

Pipes podem parecer não ser mais poderosos do que arquivos temporários:
o pipeline
\begin{lstlisting}[]
echo hello world | wc
\end{lstlisting}
poderia ser implementado sem pipes como:
\begin{lstlisting}[]
echo hello world >/tmp/xyz; wc </tmp/xyz
\end{lstlisting}
Mas os pipes possuem ao menos três vantagens sobre arquivos temporários
nesta situação.

Primeiro, pipes se limpam automaticamente;
com redirecionamento de arquivos, o shell precisaria
se preocupar em remover
\lstinline{/tmp/xyz}
ao final da execução.

Segundo, pipes podem transmitir fluxos de dados arbitrariamente longos,
enquanto o redirecionamento para arquivos exige espaço livre em disco
suficiente para armazenar todos os dados.

Terceiro, pipes permitem a execução paralela das etapas do pipeline,
enquanto a abordagem com arquivos exige que o primeiro programa finalize
antes do segundo começar.

% Quarto, se você estiver implementando comunicação entre processos,
% pipes com leituras e escritas bloqueantes são mais eficientes
% do que arquivos com semântica não bloqueante.

%% 
%% 	File system
%% 
\section{Sistema de arquivos}

O sistema de arquivos do xv6 fornece arquivos de dados,
que contêm arrays de bytes não interpretados,
e diretórios, que
contêm referências nomeadas para arquivos de dados e outros diretórios.
Os diretórios formam uma árvore, começando
em um diretório especial chamado
\indextext{raiz} (root).
Um
\indextext{caminho} (path)
como
\lstinline{/a/b/c}
refere-se ao arquivo ou diretório chamado
\lstinline{c}
dentro do diretório chamado
\lstinline{b},
dentro do diretório chamado
\lstinline{a},
no diretório raiz
\lstinline{/}.
Caminhos que não começam com
\lstinline{/}
são avaliados em relação ao
\indextext{diretório atual}
do processo chamador,
que pode ser alterado com a chamada de sistema
\lstinline{chdir}.
Ambos os trechos de código a seguir abrem o mesmo arquivo
(assumindo que todos os diretórios existam):

\begin{lstlisting}[]
chdir("/a");
chdir("b");
open("c", O_RDONLY);

open("/a/b/c", O_RDONLY);
\end{lstlisting}

O primeiro trecho altera o diretório atual do processo para
\lstinline{/a/b};
o segundo não altera nem depende do diretório atual do processo.

Existem chamadas de sistema para criar novos arquivos e diretórios:
\lstinline{mkdir}
cria um novo diretório,
\lstinline{open}
com a flag
\lstinline{O_CREATE}
cria um novo arquivo de dados,
e
\lstinline{mknod}
cria um novo arquivo de dispositivo.
Este exemplo ilustra os três casos:
\begin{lstlisting}[]
mkdir("/dir");
fd = open("/dir/file", O_CREATE|O_WRONLY);
close(fd);
mknod("/console", 1, 1);
\end{lstlisting}

\lstinline{mknod}
cria um arquivo especial que se refere a um dispositivo.
Associados a um arquivo de dispositivo estão
o número major e o número minor do dispositivo
(os dois argumentos de
\lstinline{mknod}),
que identificam de forma única um dispositivo do kernel.
Quando um processo posteriormente abre um arquivo de dispositivo, o kernel
desvia as chamadas de sistema
\lstinline{read}
e
\lstinline{write}
para a implementação do dispositivo no kernel,
em vez de passá-las ao sistema de arquivos.

O nome de um arquivo é distinto do próprio arquivo;
o mesmo arquivo subjacente, chamado
\indextext{inode},
pode ter vários nomes,
chamados
\indextext{links}.
Cada link consiste em uma entrada de diretório;
essa entrada contém um nome de arquivo e uma referência
a um inode.
Um inode armazena
\indextext{metadados}
sobre um arquivo, incluindo
seu tipo (arquivo, diretório ou dispositivo),
seu tamanho,
a localização do conteúdo do arquivo no disco
e o número de links associados ao arquivo.

A chamada de sistema
\lstinline{fstat}
recupera informações do inode ao qual um
descritor de arquivo se refere.
Ela preenche uma
\lstinline{struct stat},
definida em
\lstinline{stat.h} \fileref{kernel/stat.h}
como:
\begin{lstlisting}[]
#define T_DIR     1   // Diretório
#define T_FILE    2   // Arquivo
#define T_DEVICE  3   // Dispositivo

struct stat {
  int dev;     // Dispositivo de disco do sistema de arquivos
  uint ino;    // Número do inode
  short type;  // Tipo de arquivo
  short nlink; // Número de links para o arquivo
  uint64 size; // Tamanho do arquivo em bytes
};
\end{lstlisting}

A chamada de sistema
\lstinline{link}
cria outro nome no sistema de arquivos
referente ao mesmo inode de um arquivo existente.
Este trecho cria um novo arquivo com os nomes
\lstinline{a}
e
\lstinline{b}:
\begin{lstlisting}[]
open("a", O_CREATE|O_WRONLY);
link("a", "b");
\end{lstlisting}
Ler ou escrever em
\lstinline{a}
é o mesmo que ler ou escrever em
\lstinline{b}.
Cada inode é identificado por um número único de
\textit{inode}.
Após a execução do código acima, é possível
verificar que
\lstinline{a}
e
\lstinline{b}
referem-se ao mesmo conteúdo subjacente, inspecionando o resultado de
\lstinline{fstat}:
ambos retornarão o mesmo número de inode
(\lstinline{ino}),
e a contagem de
\lstinline{nlink}
será 2.

A chamada de sistema
\lstinline{unlink}
remove um nome do sistema de arquivos.
O inode e o espaço de disco associado ao conteúdo do arquivo
só são liberados quando a contagem de links for zero e
nenhum descritor de arquivo estiver referenciando o arquivo.
Assim, ao adicionar:
\begin{lstlisting}[]
unlink("a");
\end{lstlisting}
ao código anterior, o inode
e o conteúdo do arquivo continuam acessíveis por meio de
\lstinline{b}.
Além disso, o trecho:
\begin{lstlisting}[]
fd = open("/tmp/xyz", O_CREATE|O_RDWR);
unlink("/tmp/xyz");
\end{lstlisting}
é uma forma idiomática de criar um inode temporário
sem nome,
que será automaticamente removido quando o processo fechar
\lstinline{fd}
ou terminar sua execução.

O Unix fornece
utilitários de arquivos que podem ser chamados a partir do shell
como programas de usuário, por exemplo
\lstinline{mkdir},
\lstinline{ln}
e
\lstinline{rm}.
Esse design permite que qualquer pessoa estenda a interface de linha de comando
adicionando novos programas de usuário.
Em retrospecto, esse plano parece óbvio,
mas outros sistemas projetados na mesma época que o Unix muitas vezes embutiam tais comandos diretamente
no shell (e até mesmo o shell no kernel).

Uma exceção é o comando
\lstinline{cd},
que é incorporado ao próprio shell.
\lstinline{cd}
precisa alterar o diretório de trabalho atual do
próprio shell.
Se
\lstinline{cd}
fosse executado como um comando comum, o shell faria um
\textit{fork} de um processo filho, e o processo filho executaria
\lstinline{cd},
mudando o diretório de trabalho do
\textit{filho},
mas o diretório de trabalho do processo pai (ou seja, o shell)
permaneceria inalterado.

%% 
%% 	Real world
%% 
\section{Mundo real}

A combinação do Unix de descritores de arquivos “padrão”,
pipes e uma sintaxe de shell conveniente para operações sobre eles
foi um grande avanço na escrita de programas reutilizáveis de propósito geral.
Essa ideia deu origem a uma cultura de “ferramentas de software” que foi
responsável por grande parte do poder e da popularidade do Unix,
e o shell foi a primeira chamada “linguagem de script”.
A interface de chamadas de sistema do Unix persiste até hoje em sistemas como
BSD, Linux e macOS.

A interface de chamadas de sistema do Unix foi padronizada através do padrão
Portable Operating System Interface (POSIX).
O xv6
\textit{não}
é compatível com POSIX: ele carece de muitas chamadas de sistema (incluindo algumas básicas como
\lstinline{lseek}),
e muitas das chamadas que ele fornece diferem do padrão.
Nossos principais objetivos com o xv6 são
simplicidade e clareza, ao mesmo tempo em que fornece uma interface de chamadas de sistema simples e semelhante ao UNIX.
Várias pessoas estenderam o xv6 com algumas chamadas de sistema adicionais e uma biblioteca C simples para executar programas Unix básicos.
No entanto, kernels modernos fornecem muitas mais chamadas de sistema e tipos de serviços do kernel do que o xv6.
Por exemplo, eles suportam redes, sistemas de janelas, threads em nível de usuário,
drivers para muitos dispositivos e assim por diante.
Kernels modernos evoluem continuamente e rapidamente,
e oferecem muitos recursos além do POSIX.

O Unix unificou o acesso a múltiplos tipos de recursos (arquivos,
diretórios e dispositivos) com um único conjunto de interfaces
baseadas em nomes de arquivos e descritores de arquivos.
Essa ideia pode ser estendida a mais tipos de recursos;
um bom exemplo é o sistema Plan 9~\cite{Presotto91plan9},
que aplicou o conceito de “recursos são arquivos”
a redes, gráficos e mais.
No entanto, a maioria dos sistemas operacionais derivados do Unix
não seguiu esse caminho.

O sistema de arquivos e os descritores de arquivos têm sido abstrações poderosas.
Ainda assim, existem outros modelos para interfaces de sistemas operacionais.
O Multics, predecessor do Unix,
abstraía o armazenamento de arquivos de forma que o tornava semelhante à memória,
resultando em uma interface com uma abordagem bastante diferente.
A complexidade do projeto do Multics teve influência direta
nos projetistas do Unix, que buscaram criar algo mais simples.

O xv6 não fornece a noção de usuários ou proteção
entre usuários; em termos de Unix, todos os processos no xv6
são executados como root.

Este livro examina como o xv6 implementa sua interface semelhante ao Unix,
mas as ideias e conceitos se aplicam a mais do que apenas ao Unix.
Qualquer sistema operacional deve multiplexar processos sobre o
hardware subjacente, isolar processos entre si
e fornecer mecanismos para comunicação interprocessos controlada.
Após estudar o xv6, você deverá ser capaz de
analisar outros sistemas operacionais mais complexos
e reconhecer neles os conceitos fundamentais do xv6.

%%
\section{Exercícios}
%%

\begin{enumerate}

\item Escreva um programa que utilize chamadas de sistema UNIX para
fazer um “ping-pong” de um byte entre dois processos usando um par
de pipes, um para cada direção. Meça o
desempenho do programa, em trocas por segundo.

\end{enumerate}


 Capítulo  1
 Interfaces  do  sistema  operacional
 A  função  de  um  sistema  operacional  é  compartilhar  um  computador  entre  vários  programas  e  fornecer  um  conjunto  de  serviços  
mais  útil  do  que  o  hardware  suporta  sozinho.  Um  sistema  operacional  gerencia  e  abstrai  o  hardware  de  baixo  nível,  de  modo  que,  
por  exemplo,  um  processador  de  texto  não  precise  se  preocupar  com  o  tipo  de  hardware  de  disco  que  está  sendo  usado.  Um  
sistema  operacional  compartilha  o  hardware  entre  vários  programas  para  que  eles  sejam  executados  (ou  pareçam  ser  executados)  
ao  mesmo  tempo.  Por  fim,  os  sistemas  operacionais  fornecem  maneiras  controladas  para  os  programas  interagirem,  para  que  
possam  compartilhar  dados  ou  trabalhar  em  conjunto.
 Um  sistema  operacional  fornece  serviços  aos  programas  do  usuário  por  meio  de  uma  interface.  Projetar  uma  boa  interface  
acaba  sendo  difícil.  Por  um  lado,  gostaríamos  que  a  interface  fosse  simples  e  concisa,  pois  isso  facilita  a  implementação  correta.  
Por  outro  lado,  podemos  ser  tentados  a  oferecer  muitos  recursos  sofisticados  aos  aplicativos.  O  segredo  para  resolver  essa  
tensão  é  projetar  interfaces  que  dependam  de  alguns  mecanismos  que  possam  ser  combinados  para  fornecer  muita  generalidade.
 Este  livro  utiliza  um  único  sistema  operacional  como  exemplo  concreto  para  ilustrar  conceitos  de  sistema  operacional.  Esse  
sistema  operacional,  o  xv6,  fornece  as  interfaces  básicas  introduzidas  pelo  sistema  operacional  Unix  de  Ken  Thompson  e  Dennis  
Ritchie  [17],  além  de  imitar  o  design  interno  do  Unix.
 O  Unix  oferece  uma  interface  estreita  cujos  mecanismos  se  combinam  bem,  oferecendo  um  grau  surpreendente  de  generalidade.  
Essa  interface  tem  sido  tão  bem-sucedida  que  os  sistemas  operacionais  modernos  —  BSD,  Linux,  macOS,  Solaris  e  até  mesmo,  
em  menor  grau,  o  Microsoft  Windows  —  possuem  interfaces  semelhantes  às  do  Unix.  Entender  o  xv6  é  um  bom  começo  para  
entender  qualquer  um  desses  sistemas  e  muitos  outros.
 Como  mostra  a  Figura  1.1,  o  xv6  assume  a  forma  tradicional  de  um  kernel,  um  programa  especial  que  fornece  serviços  aos  
programas  em  execução.  Cada  programa  em  execução,  chamado  de  processo,  possui  memória  contendo  instruções,  dados  e  
uma  pilha.  As  instruções  implementam  a  computação  do  programa.  Os  dados  são  as  variáveis  sobre  as  quais  a  computação  atua.  
A  pilha  organiza  as  chamadas  de  procedimento  do  programa.
 Um  determinado  computador  normalmente  tem  muitos  processos,  mas  apenas  um  único  kernel.
 Quando  um  processo  precisa  invocar  um  serviço  do  kernel,  ele  invoca  uma  chamada  de  sistema,  uma  das  chamadas  na  
interface  do  sistema  operacional.  A  chamada  de  sistema  entra  no  kernel;  o  kernel  executa  o  serviço  e  retorna.  Assim,  um  processo  
alterna  entre  a  execução  no  espaço  do  usuário  e  no  espaço  do  kernel.
 Conforme  descrito  em  detalhes  nos  capítulos  subsequentes,  o  kernel  usa  os  mecanismos  de  proteção  de  hardware  
fornecidos  por  uma  CPU1  para  garantir  que  cada  processo  em  execução  no  espaço  do  usuário  possa  acessar  apenas
 1Este  texto  geralmente  se  refere  ao  elemento  de  hardware  que  executa  uma  computação  com  o  termo  CPU,  uma  sigla
 9
Machine Translated by Google
 usuário
 espaço
 núcleo
 espaço
 concha
 gato
 chamada  de  
sistema
 Núcleo
 Figura  1.1:  Um  kernel  e  dois  processos  de  usuário.
 sua  própria  memória.  O  kernel  executa  com  os  privilégios  de  hardware  necessários  para  implementar  essas  proteções;  os  programas  do  
usuário  são  executados  sem  esses  privilégios.  Quando  um  programa  do  usuário  invoca  uma  chamada  de  sistema,  o  hardware  aumenta  
o  nível  de  privilégio  e  começa  a  executar  uma  função  pré-definida  no  kernel.
 O  conjunto  de  chamadas  de  sistema  que  um  kernel  fornece  é  a  interface  que  os  programas  do  usuário  veem.  O  kernel  xv6  fornece  
um  subconjunto  dos  serviços  e  chamadas  de  sistema  que  os  kernels  Unix  tradicionalmente  oferecem.
 A  Figura  1.2  lista  todas  as  chamadas  de  sistema  do  xv6.
 O  restante  deste  capítulo  descreve  os  serviços  do  xv6  —  processos,  memória,  descritores  de  arquivo,  pipes  e  um  sistema  de  
arquivos  —  e  os  ilustra  com  trechos  de  código  e  discussões  sobre  como  o  shell,  a  interface  de  linha  de  comando  do  Unix,  os  utiliza.  O  
uso  de  chamadas  de  sistema  pelo  shell  ilustra  o  cuidado  com  que  foram  projetados.
 O  shell  é  um  programa  comum  que  lê  comandos  do  usuário  e  os  executa.  O  fato  de  o  shell  ser  um  programa  do  usuário,  e  não  parte  do  kernel,  ilustra  o  poder  da  
interface  de  chamada  de  sistema:  não  há  nada  de  especial  no  shell.  Isso  também  significa  que  o  shell  é  fácil  de  substituir;  como  resultado,  os  sistemas  Unix  modernos  
têm  uma  variedade  de  shells  para  escolher,  cada  um  com  sua  própria  interface  de  usuário  e  recursos  de  script.  O  shell  xv6  é  uma  implementação  simples  da  essência  do  
shell  Bourne  do  Unix.  Sua  implementação  pode  ser  encontrada  em  (user/sh.c:1).
 1.1  Processos  e  memória
 Um  processo  xv6  consiste  em  memória  de  espaço  do  usuário  (instruções,  dados  e  pilha)  e  estado  por  processo,  privado  do  kernel.  O  Xv6  
compartilha  o  tempo  dos  processos:  ele  alterna  de  forma  transparente  as  CPUs  disponíveis  entre  o  conjunto  de  processos  aguardando  
execução.  Quando  um  processo  não  está  em  execução,  o  xv6  salva  os  registradores  de  CPU  do  processo,  restaurando-os  na  próxima  
execução.  O  kernel  associa  um  identificador  de  processo,  ou  PID,  a  cada  processo.
 Um  processo  pode  criar  um  novo  processo  usando  a  chamada  de  sistema  fork .  Fork  fornece  ao  novo  processo  uma  cópia  exata  da  
memória  do  processo  que  fez  a  chamada,  tanto  instruções  quanto  dados.  Fork  retorna  tanto  no  processo  original  quanto  no  novo.  No  
processo  original,  fork  retorna  o  PID  do  novo  processo.  No  novo  processo,  fork  retorna  zero.  Os  processos  original  e  novo  são  
frequentemente  chamados  de  pai  e  filho.
 para  unidade  central  de  processamento.  Outra  documentação  (por  exemplo,  a  especificação  RISC-V)  também  utiliza  as  palavras  processador,  
núcleo  e  hart  em  vez  de  CPU.
 10
Machine Translated by Google
 Chamada  de  sistema
 int  fork()  int  
exit(int  status)  int  
wait(int  *status)  int  kill(int  
pid)  int  getpid()  int  
sleep(int  n)  int  
exec(char  *file,  char  
Descrição
 Crie  um  processo  e  retorne  o  PID  do  filho.
 Encerra  o  processo  atual;  status  reportado  para  wait().  Sem  retorno.
 Aguarde  a  saída  de  uma  criança;  saia  do  status  em  *status;  retorna  o  PID  da  criança.
 Encerra  o  PID  do  processo.  Retorna  0  ou  -1  em  caso  de  erro.
 Retorna  o  PID  do  processo  atual.
 Pausa  por  n  tiques  de  relógio.
 *argv[])  Carrega  um  arquivo  e  o  executa  com  argumentos;  retorna  somente  se  houver  erro.
 char  *sbrk(int  n)  int  
open(char  *file,  int  flags)  int  write(int  
fd,  char  *buf,  int  n)  Grava  n  bytes  de  buf  no  descritor  de  arquivo  fd;  retorna  n.
 Aumenta  a  memória  do  processo  em  n  bytes.  Retorna  o  início  da  nova  memória.
 Abre  um  arquivo;  sinalizadores  indicam  leitura/gravação;  retorna  um  fd  (descritor  de  arquivo).
 int  read(int  fd,  char  *buf,  int  n)  Lê  n  bytes  em  buf;  retorna  o  número  lido;  ou  0  se  for  o  fim  do  arquivo.
 int  fechar(int  fd)
 Libere  o  arquivo  aberto  fd.
 int  dup(int  fd)
 int  pipe(int  p[])
 int  chdir(char  *dir)
 int  mkdir(char  *dir)
 int  mknod(char  *arquivo,  int,  int)
 int  fstat(int  fd,  struct  stat  *st)
 Retorna  um  novo  descritor  de  arquivo  referindo-se  ao  mesmo  arquivo  que  fd.
 Crie  um  pipe,  coloque  descritores  de  arquivo  de  leitura/gravação  em  p[0]  e  p[1].
 Altera  o  diretório  atual.
 Crie  um  novo  diretório.
 Crie  um  arquivo  de  dispositivo.
 Coloque  informações  sobre  um  arquivo  aberto  em  *st.
 int  stat(char  *file,  struct  stat  *st)  Coloque  informações  sobre  um  arquivo  nomeado  em  *st.
 int  link(char  *file1,  char  *file2)  Cria  outro  nome  (file2)  para  o  arquivo  file1.
 int  unlink(char  *arquivo)
 Remover  um  arquivo.
 Figura  1.2:  Chamadas  de  sistema  Xv6.  Se  não  indicado  de  outra  forma,  essas  chamadas  retornam  0  para  nenhum  erro  e  -1  se
 há  um  erro.
 Por  exemplo,  considere  o  seguinte  fragmento  de  programa  escrito  na  linguagem  de  programação  C  [7]:
 int  pid  =  garfo();
 se(pid  >  0){
 printf("pai:  filho=%d\n",  pid);
 pid  =  esperar((int  *)  0);
 printf("criança  %d  terminou\n",  pid);
 }  senão  se(pid  ==  0){
 printf("filho:  saindo\n");
 saída(0);
 }  outro  {
 printf("erro  de  bifurcação\n");
 }
 A  chamada  de  sistema  exit  faz  com  que  o  processo  de  chamada  pare  de  executar  e  libere  recursos  como
 memória  e  arquivos  abertos.  Exit  recebe  um  argumento  de  status  inteiro,  convencionalmente  0  para  indicar  sucesso
 e  1  para  indicar  falha.  A  chamada  de  sistema  wait  retorna  o  PID  de  um  filho  saído  (ou  eliminado)  de
 11
Machine Translated by Google
 o  processo  atual  e  copia  o  status  de  saída  do  filho  para  o  endereço  passado  para  wait;  se  nenhum  dos  filhos  do  chamador  
tiver  saído,  wait  aguarda  que  um  o  faça.  Se  o  chamador  não  tiver  filhos,  wait  retorna  imediatamente  -1.  Se  o  pai  não  se  
importar  com  o  status  de  saída  de  um  filho,  ele  pode  passar  um  endereço  0  para  wait.
 No  exemplo,  as  linhas  de  saída
 pai:  filho=1234  filho:  saindo
 pode  sair  em  qualquer  ordem  (ou  até  mesmo  misturados),  dependendo  se  o  pai  ou  o  filho  chega  primeiro  à  sua  
chamada  printf .  Após  a  saída  do  filho,  a  espera  do  pai  retorna,  fazendo  com  que  o  pai  imprima
 pai:  filho  1234  está  pronto
 Embora  o  processo  filho  tenha  inicialmente  o  mesmo  conteúdo  de  memória  que  o  processo  pai,  ambos  estão  
executando  com  memória  e  registradores  separados:  alterar  uma  variável  em  um  não  afeta  o  outro.  Por  exemplo,  
quando  o  valor  de  retorno  de  wait  é  armazenado  em  pid  no  processo  pai,  isso  não  altera  a  variável  pid  no  processo  
filho.  O  valor  de  pid  no  processo  filho  ainda  será  zero.
 A  chamada  de  sistema  exec  substitui  a  memória  do  processo  que  fez  a  chamada  por  uma  nova  imagem  de  
memória  carregada  de  um  arquivo  armazenado  no  sistema  de  arquivos.  O  arquivo  deve  ter  um  formato  específico,  
que  especifica  qual  parte  do  arquivo  contém  instruções,  qual  parte  são  dados,  em  qual  instrução  iniciar,  etc.  O  Xv6  
usa  o  formato  ELF,  que  o  Capítulo  3  discute  em  mais  detalhes.  Normalmente,  o  arquivo  é  o  resultado  da  compilação  
do  código-fonte  de  um  programa.  Quando  exec  é  bem-sucedido,  ele  não  retorna  ao  programa  que  fez  a  chamada;  
em  vez  disso,  as  instruções  carregadas  do  arquivo  começam  a  ser  executadas  no  ponto  de  entrada  declarado  no  
cabeçalho  ELF.  exec  recebe  dois  argumentos:  o  nome  do  arquivo  que  contém  o  executável  e  uma  matriz  de  
argumentos  de  string.  Por  exemplo:
 char  *argv[3];
 argv[0]  =  "eco";  argv[1]  =  
"olá";  argv[2]  =  0;  exec("/bin/
 eco",  argv);  
printf("erro  de  execução\n");
 Este  fragmento  substitui  o  programa  chamador  por  uma  instância  do  programa /bin/echo  em  execução  com  a  lista  
de  argumentos  echo  hello.  A  maioria  dos  programas  ignora  o  primeiro  elemento  da  matriz  de  argumentos,  que  é  
convencionalmente  o  nome  do  programa.
 O  shell  xv6  usa  as  chamadas  acima  para  executar  programas  em  nome  dos  usuários.  A  estrutura  principal  do  
shell  é  simples;  veja  main  (user/sh.c:146).  O  loop  principal  lê  uma  linha  de  entrada  do  usuário  com  getcmd.  Em  
seguida,  chama  fork,  que  cria  uma  cópia  do  processo  do  shell.  O  processo  pai  chama  wait,  enquanto  o  processo  
filho  executa  o  comando.  Por  exemplo,  se  o  usuário  tivesse  digitado  "echo  hello"  no  shell,  runcmd  teria  sido  
chamado  com  "echo  hello"  como  argumento.  runcmd  (user/sh.c:55)  executa  o  comando  real.  Para  "echo  hello",  ele  
chamaria  exec  (user/sh.c:79).  Se  exec  for  bem-sucedido,  o  filho  executará  instruções  de  echo  em  vez  de  runcmd.  
Em  algum  momento,  echo  chamará  exit,  o  que  fará  com  que  o  pai  retorne  de  wait  em  main  (user/sh.c:146).
 12
Machine Translated by Google
 Você  pode  se  perguntar  por  que  fork  e  exec  não  são  combinados  em  uma  única  chamada;  veremos  mais  
adiante  que  o  shell  explora  essa  separação  em  sua  implementação  de  redirecionamento  de  E/S.  Para  evitar  o  
desperdício  de  criar  um  processo  duplicado  e  substituí-lo  imediatamente  (por  exec),  os  kernels  operacionais  
otimizam  a  implementação  de  fork  para  esse  caso  de  uso  usando  técnicas  de  memória  virtual,  como  cópia  na  
gravação  (consulte  a  Seção  4.6).
 O  Xv6  aloca  a  maior  parte  da  memória  do  espaço  do  usuário  implicitamente:  fork  aloca  a  memória  necessária  
para  a  cópia  filha  da  memória  do  pai,  e  exec  aloca  memória  suficiente  para  armazenar  o  arquivo  executável.  Um  
processo  que  precisa  de  mais  memória  em  tempo  de  execução  (talvez  para  malloc)  pode  chamar  sbrk(n)  para  
aumentar  sua  memória  de  dados  em  n  bytes;  sbrk  retorna  a  localização  da  nova  memória.
 1.2  E/S  e  descritores  de  arquivo
 Um  descritor  de  arquivo  é  um  pequeno  inteiro  que  representa  um  objeto  gerenciado  pelo  kernel  do  qual  um  
processo  pode  ler  ou  gravar.  Um  processo  pode  obter  um  descritor  de  arquivo  abrindo  um  arquivo,  diretório  ou  
dispositivo,  criando  um  pipe  ou  duplicando  um  descritor  existente.  Para  simplificar,  frequentemente  nos  referimos  
ao  objeto  ao  qual  um  descritor  de  arquivo  se  refere  como  "arquivo";  a  interface  do  descritor  de  arquivo  abstrai  
as  diferenças  entre  arquivos,  pipes  e  dispositivos,  fazendo  com  que  todos  pareçam  fluxos  de  bytes.  Nos  
referiremos  à  entrada  e  à  saída  como  E/S.
 Internamente,  o  kernel  xv6  usa  o  descritor  de  arquivo  como  um  índice  em  uma  tabela  por  processo,  de  
modo  que  cada  processo  tenha  um  espaço  privado  de  descritores  de  arquivo  começando  em  zero.  Por  
convenção,  um  processo  lê  do  descritor  de  arquivo  0  (entrada  padrão),  grava  a  saída  no  descritor  de  arquivo  1  
(saída  padrão)  e  grava  mensagens  de  erro  no  descritor  de  arquivo  2  (erro  padrão).  Como  veremos,  o  shell  
explora  a  convenção  para  implementar  redirecionamentos  de  E/S  e  pipelines.  O  shell  garante  que  sempre  tenha  
três  descritores  de  arquivo  abertos  (user/sh.c:152).  que  são,  por  padrão,  descritores  de  arquivo  para  o  console.
 O  sistema  de  leitura  e  gravação  chama  bytes  de  leitura  e  bytes  de  gravação  para  abrir  arquivos  nomeados  por  
descritores  de  arquivo.  A  chamada  read(fd,  buf,  n)  lê  no  máximo  n  bytes  do  descritor  de  arquivo  fd,  copia-os  para  buf  e  
retorna  o  número  de  bytes  lidos.  Cada  descritor  de  arquivo  que  se  refere  a  um  arquivo  possui  um  deslocamento  associado.  
read  lê  os  dados  do  deslocamento  do  arquivo  atual  e,  em  seguida,  avança  esse  deslocamento  pelo  número  de  bytes  lidos:  
uma  leitura  subsequente  retornará  os  bytes  seguintes  aos  retornados  pela  primeira  leitura.  Quando  não  há  mais  bytes  
para  ler,  read  retorna  zero  para  indicar  o  fim  do  arquivo.
 A  chamada  write(fd,  buf,  n)  grava  n  bytes  de  buf  no  descritor  de  arquivo  fd  e  retorna  o  número  de  bytes  gravados.  
Menos  de  n  bytes  são  gravados  somente  quando  ocorre  um  erro.  Assim  como  read,  write  grava  dados  no  deslocamento  
do  arquivo  atual  e,  em  seguida,  avança  esse  deslocamento  pelo  número  de  bytes  gravados:  cada  gravação  continua  de  
onde  a  anterior  parou.
 O  seguinte  fragmento  de  programa  (que  constitui  a  essência  do  programa  cat)  copia  dados  de  sua  entrada  
padrão  para  sua  saída  padrão.  Se  ocorrer  um  erro,  ele  grava  uma  mensagem  na  entrada  padrão.
 erro.
 char  buf[512];  int  n;
 para(;;){
 13
Machine Translated by Google
 n  =  ler(0,  buf,  tamanho  de  buf);  se(n  ==  0)  quebrar;  
se(n  <  0)
 { fprintf(2,  
"erro  de  
leitura\n");  sair(1);
 }  se(escrever(1,  buf,  n) !=  n){
 fprintf(2,  "erro  de  gravação\n");  exit(1);
 }
 }
 O  importante  a  ser  observado  no  fragmento  de  código  é  que  cat  não  sabe  se  está  lendo  de  um  arquivo,  console  ou  pipe.  Da  
mesma  forma,  cat  não  sabe  se  está  imprimindo  para  um  console,  um  arquivo  ou  qualquer  outro  lugar.  O  uso  de  descritores  
de  arquivo  e  a  convenção  de  que  o  descritor  de  arquivo  0  é  a  entrada  e  o  descritor  de  arquivo  1  é  a  saída  permitem  uma  
implementação  simples  de  cat.
 A  chamada  de  sistema  close  libera  um  descritor  de  arquivo,  tornando-o  livre  para  reutilização  em  futuras  chamadas  de  
sistema  open,  pipe  ou  dup  (veja  abaixo).  Um  descritor  de  arquivo  recém-alocado  é  sempre  o  descritor  não  utilizado  de  
menor  numeração  do  processo  atual.
 Descritores  de  arquivo  e  fork  interagem  para  facilitar  a  implementação  do  redirecionamento  de  E/S.  O  fork  copia  a  tabela  
de  descritores  de  arquivo  do  processo  pai  juntamente  com  sua  memória,  de  modo  que  o  processo  filho  comece  com  exatamente  
os  mesmos  arquivos  abertos  que  o  processo  pai.  A  chamada  de  sistema  exec  substitui  a  memória  do  processo  que  fez  a  
chamada,  mas  preserva  sua  tabela  de  arquivos.  Esse  comportamento  permite  que  o  shell  implemente  o  redirecionamento  de  E/
 S  por  meio  de  fork,  reabrindo  os  descritores  de  arquivo  escolhidos  no  processo  filho  e,  em  seguida,  chamando  exec  para  
executar  o  novo  programa.  Aqui  está  uma  versão  simplificada  do  código  que  um  shell  executa  para  o  comando  cat  <  input.txt:
 char  *argv[2];
 argv[0]  =  "gato";  argv[1]  =  
0;  if(fork()  ==  0)  
{ fechar(0);  abrir("input.txt",  
O_RDONLY);  
exec("gato",  argv);
 }
 Após  o  processo  filho  fechar  o  descritor  de  arquivo  0,  open  garante  que  usará  esse  descritor  de  arquivo  para  o  
arquivo  input.txt  recém-aberto:  0  será  o  menor  descritor  de  arquivo  disponível.  cat  então  executa  com  o  descritor  de  
arquivo  0  (entrada  padrão)  referindo-se  ao  arquivo  input.txt.  Os  descritores  de  arquivo  do  processo  pai  não  são  
alterados  por  esta  sequência,  pois  ela  modifica  apenas  os  descritores  do  processo  filho.
 O  código  para  redirecionamento  de  E/S  no  shell  xv6  funciona  exatamente  dessa  maneira  (user/sh.c:83).  Lembre-se  de  
que  neste  ponto  do  código  o  shell  já  bifurcou  o  shell  filho  e  que  runcmd  chamará  exec  para  carregar  o  novo  programa.
 O  segundo  argumento  para  open  consiste  em  um  conjunto  de  sinalizadores,  expressos  como  bits,  que  controlam  
o  que  open  faz.  Os  valores  possíveis  são  definidos  no  cabeçalho  de  controle  de  arquivo  (fcntl)  (kernel/fcntl.h:1-5):
 14
Machine Translated by Google
 O_RDONLY,  O_WRONLY,  O_RDWR,  O_CREATE  e  O_TRUNC,  que  instruem  open  a  abrir  o  arquivo  para  leitura,  ou  para  
gravação,  ou  para  leitura  e  gravação,  para  criar  o  arquivo  se  ele  não  existir  e  para  truncá-lo  para  comprimento  zero.
 Agora  deve  ficar  claro  por  que  é  útil  que  fork  e  exec  sejam  chamadas  separadas:  entre  as  duas,  o  shell  tem  a  chance  de  
redirecionar  a  E/S  do  shell  filho  sem  perturbar  a  configuração  de  E/S  do  shell  principal.  Pode-se,  em  vez  disso,  imaginar  uma  
hipotética  chamada  de  sistema  combinada  forkexec ,  mas  as  opções  para  realizar  o  redirecionamento  de  E/S  com  tal  
chamada  parecem  estranhas.  O  shell  poderia  modificar  sua  própria  configuração  de  E/S  antes  de  chamar  forkexec  (e  então  
desfazer  essas  modificações);  ou  forkexec  poderia  receber  instruções  de  redirecionamento  de  E/S  como  argumentos;  ou  (o  
que  é  menos  atraente)  todo  programa,  como  cat,  poderia  ser  ensinado  a  fazer  seu  próprio  redirecionamento  de  E/S.
 Embora  o  fork  copie  a  tabela  do  descritor  de  arquivo,  cada  deslocamento  de  arquivo  subjacente  é  compartilhado  entre
 Pais  e  filhos.  Considere  este  exemplo:
 if(fork()  ==  0)  { write(1,  "olá  
",  6);  exit(0); }  else  { wait(0);  write(1,  
"mundo\n",  
6);
 }
 No  final  deste  fragmento,  o  arquivo  anexado  ao  descritor  de  arquivo  1  conterá  os  dados  hello  world.
 A  gravação  no  pai  (que,  graças  à  espera,  só  é  executada  após  a  conclusão  do  filho)  continua  de  onde  a  gravação  do  filho  
parou.  Esse  comportamento  ajuda  a  produzir  uma  saída  sequencial  a  partir  de  sequências  de  comandos  de  shell,  como  (echo  
hello;  echo  world)  >output.txt.
 A  chamada  de  sistema  dup  duplica  um  descritor  de  arquivo  existente,  retornando  um  novo  que  se  refere  ao  mesmo  
objeto  de  E/S  subjacente.  Ambos  os  descritores  de  arquivo  compartilham  um  deslocamento,  assim  como  os  descritores  de  
arquivo  duplicados  por  fork .  Esta  é  outra  maneira  de  escrever  hello  world  em  um  arquivo:
 fd  =  dup(1);  
escrever(1,  "olá  ",  6);  escrever(fd,  
"mundo\n",  6);
 Dois  descritores  de  arquivo  compartilham  um  deslocamento  se  forem  derivados  do  mesmo  descritor  de  arquivo  original  
por  uma  sequência  de  chamadas  fork  e  dup .  Caso  contrário,  os  descritores  de  arquivo  não  compartilham  deslocamentos,  
mesmo  que  tenham  resultado  de  chamadas  open  para  o  mesmo  arquivo.  O  dup  permite  que  shells  implementem  comandos  
como  este:  ls  existing-file  non-existing-file  >  tmp1  2>&1.  O  2>&1  informa  ao  shell  para  atribuir  ao  comando  um  descritor  de  
arquivo  2,  que  é  uma  duplicata  do  descritor  1.  Tanto  o  nome  do  arquivo  existente  quanto  a  mensagem  de  erro  para  o  arquivo  
inexistente  serão  exibidos  no  arquivo  tmp1.  O  shell  xv6  não  suporta  redirecionamento  de  E/S  para  o  descritor  de  arquivo  de  
erro,  mas  agora  você  sabe  como  implementá-lo.
 Os  descritores  de  arquivo  são  uma  abstração  poderosa,  porque  eles  escondem  os  detalhes  do  que  estão  conectados:  
um  processo  escrevendo  no  descritor  de  arquivo  1  pode  estar  escrevendo  em  um  arquivo,  em  um  dispositivo  como  o  console  
ou  em  um  pipe.
 15
Machine Translated by Google
 1.3  Tubos
 Um  pipe  é  um  pequeno  buffer  de  kernel  exposto  aos  processos  como  um  par  de  descritores  de  arquivo,  um  para  leitura  e  
outro  para  escrita.  A  gravação  de  dados  em  uma  extremidade  do  pipe  torna  esses  dados  disponíveis  para  leitura  na  outra  
extremidade.  Pipes  fornecem  uma  maneira  para  os  processos  se  comunicarem.
 O  código  de  exemplo  a  seguir  executa  o  programa  wc  com  a  entrada  padrão  conectada  à  extremidade  de  leitura  de  um  
pipe.
 int  p[2];  char  
*argv[2];
 argv[0]  =  "wc";  argv[1]  =  
0;
 pipe(p);  
if(fork()  ==  0)  { fechar(0);  
dup(p[0]);  
fechar(p[0]);  
fechar(p[1]);  exec("/
 bin/wc",  argv); }  else  
{ fechar(p[0]);  escrever(p[1],  "olá  
mundo\n",  
12);  fechar(p[1]);
 }
 O  programa  chama  pipe,  que  cria  um  novo  pipe  e  registra  os  descritores  de  arquivo  de  leitura  e  gravação  no  array  p.  Após  a  
bifurcação,  tanto  o  pai  quanto  o  filho  possuem  descritores  de  arquivo  referentes  ao  pipe.  O  filho  chama  close  e  dup  para  fazer  
com  que  o  descritor  de  arquivo  zero  se  refira  à  extremidade  de  leitura  do  pipe,  fecha  os  descritores  de  arquivo  em  p  e  chama  
exec  para  executar  wc.  Quando  wc  lê  de  sua  entrada  padrão,  ele  lê  do  pipe.  O  pai  fecha  o  lado  de  leitura  do  pipe,  escreve  no  
pipe  e,  em  seguida,  fecha  o  lado  de  gravação.
 Se  não  houver  dados  disponíveis,  uma  leitura  em  um  pipe  aguarda  que  os  dados  sejam  gravados  ou  que  todos  os  
descritores  de  arquivo  referentes  à  extremidade  de  gravação  sejam  fechados;  neste  último  caso,  a  leitura  retornará  0,  como  
se  o  fim  de  um  arquivo  de  dados  tivesse  sido  alcançado.  O  fato  de  a  leitura  bloquear  até  que  seja  impossível  a  chegada  de  
novos  dados  é  um  dos  motivos  pelos  quais  é  importante  que  o  filho  feche  a  extremidade  de  gravação  do  pipe  antes  de  
executar  o  comando  wc  acima:  se  um  dos  descritores  de  arquivo  de  wc  se  referisse  à  extremidade  de  gravação  do  pipe,  wc  
nunca  veria  o  fim  do  arquivo.
 O  shell  xv6  implementa  pipelines  como  grep  fork  sh.c  |  wc  -l  de  maneira  semelhante  ao  código  acima  (usuário/sh.c:101).  
O  processo  filho  cria  um  pipe  para  conectar  a  extremidade  esquerda  do  pipeline  com  a  extremidade  direita.  Em  seguida,  ele  
chama  fork  e  runcmd  para  a  extremidade  esquerda  do  pipeline  e  fork  e  runcmd  para  a  extremidade  direita,  e  aguarda  a  
conclusão  de  ambos.  A  extremidade  direita  do  pipeline  pode  ser  um  comando  que  inclui  um  pipe  (por  exemplo,  a  |  b  |  c),  que  
bifurca  dois  novos  processos  filhos  (um  para  b  e  outro  para  c).  Assim,  o  shell  pode  criar  uma  árvore  de  processos.  As  folhas
 16
Machine Translated by Google
 dessa  árvore  são  comandos  e  os  nós  internos  são  processos  que  esperam  até  que  os  filhos  esquerdo  e  direito  sejam  concluídos.
 Os  pipes  podem  não  parecer  mais  poderosos  do  que  os  arquivos  temporários:  o  pipeline
 eco  olá  mundo  |  wc
 poderia  ser  implementado  sem  tubos  como
 eco  olá  mundo  >/tmp/xyz;  wc  </tmp/xyz
 Pipes  têm  pelo  menos  três  vantagens  sobre  arquivos  temporários  nessa  situação.  Primeiro,  pipes  se  autolimpam;  com  o  
redirecionamento  de  arquivos,  um  shell  teria  que  ter  o  cuidado  de  remover /tmp/xyz  ao  terminar.  Segundo,  pipes  podem  passar  
fluxos  de  dados  arbitrariamente  longos,  enquanto  o  redirecionamento  de  arquivos  requer  espaço  livre  suficiente  em  disco  para  
armazenar  todos  os  dados.  Terceiro,  pipes  permitem  a  execução  paralela  de  etapas  do  pipeline,  enquanto  a  abordagem  de  
arquivo  exige  que  o  primeiro  programa  termine  antes  do  segundo  iniciar.
 1.4  Sistema  de  arquivos
 O  sistema  de  arquivos  xv6  fornece  arquivos  de  dados,  que  contêm  matrizes  de  bytes  não  interpretadas,  e  diretórios,  que  contêm  
referências  nomeadas  a  arquivos  de  dados  e  outros  diretórios.  Os  diretórios  formam  uma  árvore,  começando  em  um  diretório  
especial  chamado  raiz.  Um  caminho  como /a/b/c  refere-se  ao  arquivo  ou  diretório  chamado  c  dentro  do  diretório  chamado  b  
dentro  do  diretório  chamado  a  no  diretório  raiz /.  Caminhos  que  não  começam  com /  são  avaliados  em  relação  ao  diretório  atual  
do  processo  chamador,  que  pode  ser  alterado  com  a  chamada  de  sistema  chdir .  Ambos  os  fragmentos  de  código  abrem  o  mesmo  
arquivo  (assumindo  que  todos  os  diretórios  envolvidos  existam):
 chdir("/a");  chdir("b");  
open("c",  
O_RDONLY);
 aberto("/a/b/c",  O_RDONLY);
 O  primeiro  fragmento  altera  o  diretório  atual  do  processo  para /a/b;  o  segundo  não  faz  referência  nem  altera  o  diretório  atual  do  
processo.
 Existem  chamadas  de  sistema  para  criar  novos  arquivos  e  diretórios:  mkdir  cria  um  novo  diretório,  open  com  o  
sinalizador  O_CREATE  cria  um  novo  arquivo  de  dados  e  mknod  cria  um  novo  arquivo  de  dispositivo.  Este  exemplo  
ilustra  todos  os  três:
 mkdir("/dir");  fd  =  open("/
 dir/arquivo",  O_CREATE|O_WRONLY);  close(fd);  mknod("/console",  1,  
1);
 O  mknod  cria  um  arquivo  especial  que  se  refere  a  um  dispositivo.  Associados  a  um  arquivo  de  dispositivo  estão  os  números  de  
dispositivo  principal  e  secundário  (os  dois  argumentos  do  mknod),  que  identificam  exclusivamente  um  dispositivo  do  kernel.
 Quando  um  processo  abre  posteriormente  um  arquivo  de  dispositivo,  o  kernel  desvia  chamadas  de  leitura  e  gravação  do  sistema  
para  a  implementação  do  dispositivo  do  kernel  em  vez  de  passá-las  para  o  sistema  de  arquivos.
 17
Machine Translated by Google
 O  nome  de  um  arquivo  é  diferente  do  próprio  arquivo;  o  mesmo  arquivo  subjacente,  chamado  inode,  pode  ter  vários  
nomes,  chamados  links.  Cada  link  consiste  em  uma  entrada  em  um  diretório;  a  entrada  contém  um  nome  de  arquivo  e  
uma  referência  a  um  inode.  Um  inode  contém  metadados  sobre  um  arquivo,  incluindo  seu  tipo  (arquivo,  diretório  ou  
dispositivo),  seu  comprimento,  a  localização  do  conteúdo  do  arquivo  no  disco  e  o  número  de  links  para  um  arquivo.
 A  chamada  de  sistema  fstat  recupera  informações  do  inode  ao  qual  um  descritor  de  arquivo  se  refere.
 preenche  uma  estrutura  stat,  definida  em  stat.h  (kernel/stat.h)  como:
 #define  T_DIR
 T_FILE  #define  
T_DEVICE  3 //  Dispositivo
 estrutura  stat  {
 int  dev;
 1 //  Diretório
 2 //  Arquivo  #define  
//  Dispositivo  de  disco  do  sistema  de  arquivos  
//  Número  do  inode  uint  
ino;  short  
type; //  Tipo  de  arquivo  short  nlink; //  Número  
de  links  para  o  arquivo  uint64  size; //  Tamanho  do  arquivo  em  
bytes };
 A  chamada  do  sistema  de  link  cria  outro  nome  de  sistema  de  arquivo  que  se  refere  ao  mesmo  inode  que  um  existente
 arquivo  ing.  Este  fragmento  cria  um  novo  arquivo  chamado  a  e  b.
 abrir("a",  O_CREATE|O_WRONLY);  link("a",  "b");
 Ler  ou  escrever  em  a  é  o  mesmo  que  ler  ou  escrever  em  b.  Cada  inode  é  identificado  por  um  número  de  inode  exclusivo.  
Após  a  sequência  de  código  acima,  é  possível  determinar  que  a  e  b  se  referem  ao  mesmo  conteúdo  subjacente  inspecionando  
o  resultado  de  fstat:  ambos  retornarão  o  mesmo  número  de  inode  (ino)  e  a  contagem  de  nlinks  será  definida  como  2.
 A  chamada  de  sistema  unlink  remove  um  nome  do  sistema  de  arquivos.  O  inode  do  arquivo  e  o  espaço  em  disco  que  
contém  seu  conteúdo  só  são  liberados  quando  a  contagem  de  links  do  arquivo  é  zero  e  nenhum  descritor  de  arquivo  faz  
referência  a  ele.  Adicionando  assim
 desvincular("a");
 para  a  última  sequência  de  código  deixa  o  inode  e  o  conteúdo  do  arquivo  acessíveis  como  b.  Além  disso,
 fd  =  abrir("/tmp/xyz",  O_CREATE|O_RDWR);  desvincular("/tmp/
 xyz");
 é  uma  maneira  idiomática  de  criar  um  inode  temporário  sem  nome  que  será  limpo  quando  o  processo  fechar  o  fd  ou  sair.
 O  Unix  fornece  utilitários  de  arquivo  que  podem  ser  chamados  a  partir  do  shell  como  programas  de  nível  de  usuário,  por  
exemplo ,  mkdir,  ln  e  rm.  Esse  design  permite  que  qualquer  pessoa  estenda  a  interface  de  linha  de  comando  adicionando  
novos  programas  de  nível  de  usuário.  Em  retrospectiva,  esse  plano  parece  óbvio,  mas  outros  sistemas  projetados  na  época  
do  Unix  frequentemente  incorporavam  esses  comandos  ao  shell  (e  o  shell  ao  kernel).
 Uma  exceção  é  o  cd,  que  é  incorporado  ao  shell  (user/sh.c:161).  cd  deve  alterar  o  diretório  de  trabalho  atual  do  próprio  
shell.  Se  cd  fosse  executado  como  um  comando  regular,  o  shell
 18
Machine Translated by Google
 bifurcar  um  processo  filho,  o  processo  filho  executaria  cd,  e  cd  alteraria  o  diretório  de  trabalho  do  filho.  O  diretório  
de  trabalho  do  pai  (ou  seja,  do  shell)  não  mudaria.
 1.5  Mundo  real
 A  combinação  de  descritores  de  arquivo  "padrão",  pipes  e  sintaxe  shell  conveniente  do  Unix  para  operações  sobre  
eles  foi  um  grande  avanço  na  escrita  de  programas  reutilizáveis  de  uso  geral.  A  ideia  deu  início  a  uma  cultura  de  
"ferramentas  de  software"  que  foi  responsável  por  grande  parte  do  poder  e  popularidade  do  Unix,  e  o  shell  foi  a  
primeira  chamada  "linguagem  de  script".  A  interface  de  chamada  de  sistema  do  Unix  persiste  até  hoje  em  sistemas  
como  BSD,  Linux  e  macOS.
 A  interface  de  chamada  de  sistema  do  Unix  foi  padronizada  pelo  padrão  POSIX  (Portable  Operating  System  
Interface).  O  Xv6  não  é  compatível  com  POSIX:  faltam  muitas  chamadas  de  sistema  (incluindo  as  básicas,  como  
lseek),  e  muitas  das  chamadas  de  sistema  que  ele  fornece  diferem  do  padrão.  Nossos  principais  objetivos  para  o  
xv6  são  simplicidade  e  clareza,  ao  mesmo  tempo  em  que  fornecemos  uma  interface  de  chamada  de  sistema  
simples,  semelhante  ao  UNIX.  Várias  pessoas  estenderam  o  xv6  com  algumas  chamadas  de  sistema  a  mais  e  
uma  biblioteca  C  simples  para  executar  programas  Unix  básicos.  Os  kernels  modernos,  no  entanto,  fornecem  
muito  mais  chamadas  de  sistema  e  muitos  mais  tipos  de  serviços  de  kernel  do  que  o  xv6.  Por  exemplo,  eles  
oferecem  suporte  a  redes,  sistemas  de  janelas,  threads  em  nível  de  usuário,  drivers  para  muitos  dispositivos  e  
assim  por  diante.  Os  kernels  modernos  evoluem  contínua  e  rapidamente  e  oferecem  muitos  recursos  além  do  POSIX.
 O  Unix  unificou  o  acesso  a  múltiplos  tipos  de  recursos  (arquivos,  diretórios  e  dispositivos)  com  um  único  
conjunto  de  interfaces  de  nome  e  descritor  de  arquivo.  Essa  ideia  pode  ser  estendida  a  mais  tipos  de  recursos;  
um  bom  exemplo  é  o  Plan  9  [16],  que  aplicou  o  conceito  de  “recursos  são  arquivos”  a  redes,  gráficos  e  outros.  No  
entanto,  a  maioria  dos  sistemas  operacionais  derivados  do  Unix  não  seguiu  esse  caminho.
 O  sistema  de  arquivos  e  os  descritores  de  arquivos  têm  sido  abstrações  poderosas.  Mesmo  assim,  existem  
outros  modelos  para  interfaces  de  sistemas  operacionais.  O  Multics,  um  predecessor  do  Unix,  abstraiu  o  
armazenamento  de  arquivos  de  uma  forma  que  o  fez  parecer  memória,  produzindo  uma  interface  muito  diferente.  
A  complexidade  do  design  do  Multics  teve  uma  influência  direta  nos  designers  do  Unix,  que  buscavam  construir  
algo  mais  simples.
 O  Xv6  não  fornece  uma  noção  de  usuários  ou  de  proteção  de  um  usuário  de  outro;  em  termos  Unix,
 todos  os  processos  xv6  são  executados  como  root.
 Este  livro  examina  como  o  xv6  implementa  sua  interface  semelhante  à  do  Unix,  mas  as  ideias  e  conceitos  se  
aplicam  a  mais  do  que  apenas  o  Unix.  Qualquer  sistema  operacional  deve  multiplexar  processos  no  hardware  
subjacente,  isolar  processos  uns  dos  outros  e  fornecer  mecanismos  para  comunicação  controlada  entre  processos.  
Depois  de  estudar  o  xv6,  você  deverá  ser  capaz  de  analisar  outros  sistemas  operacionais  mais  complexos  e  
entender  os  conceitos  subjacentes  ao  xv6  nesses  sistemas  também

 Capítulo  6
 Bloqueio
 A  maioria  dos  kernels,  incluindo  o  xv6,  intercala  a  execução  de  múltiplas  atividades.  Uma  fonte  de  intercalação  é  o  
hardware  multiprocessador:  computadores  com  múltiplas  CPUs  executando  independentemente,  como  o  RISC-V  do  
xv6.  Essas  múltiplas  CPUs  compartilham  RAM  física,  e  o  xv6  explora  esse  compartilhamento  para  manter  estruturas  
de  dados  que  todas  as  CPUs  leem  e  gravam.  Esse  compartilhamento  aumenta  a  possibilidade  de  uma  CPU  ler  uma  
estrutura  de  dados  enquanto  outra  CPU  está  no  meio  da  atualização,  ou  mesmo  múltiplas  CPUs  atualizando  os  mesmos  
dados  simultaneamente;  sem  um  projeto  cuidadoso,  esse  acesso  paralelo  provavelmente  produzirá  resultados  incorretos  
ou  uma  estrutura  de  dados  quebrada.  Mesmo  em  um  uniprocessador,  o  kernel  pode  alternar  a  CPU  entre  várias  threads,  
fazendo  com  que  sua  execução  seja  intercalada.  Por  fim,  um  manipulador  de  interrupção  de  dispositivo  que  modifica  os  
mesmos  dados  como  algum  código  interrompível  pode  danificar  os  dados  se  a  interrupção  ocorrer  no  momento  errado.  
O  termo  simultaneidade  se  refere  a  situações  em  que  múltiplos  fluxos  de  instruções  são  intercalados,  devido  ao  
paralelismo  do  multiprocessador,  troca  de  threads  ou  interrupções.
 Os  kernels  estão  repletos  de  dados  acessados  concorrentemente.  Por  exemplo,  duas  CPUs  poderiam  chamar  
kalloc  simultaneamente,  saindo  simultaneamente  do  topo  da  lista  de  acessos  livres.  Os  projetistas  de  kernel  gostam  de  
permitir  muita  concorrência,  pois  isso  pode  gerar  maior  desempenho  por  meio  do  paralelismo  e  maior  capacidade  de  
resposta.  No  entanto,  como  resultado,  os  projetistas  de  kernel  precisam  se  convencer  da  correção,  apesar  dessa  
concorrência.  Existem  muitas  maneiras  de  se  chegar  a  um  código  correto,  algumas  mais  fáceis  de  raciocinar  do  que  
outras.  Estratégias  que  visam  à  correção  sob  concorrência,  e  abstrações  que  as  suportam,  são  chamadas  de  técnicas  
de  controle  de  concorrência.
 O  Xv6  utiliza  diversas  técnicas  de  controle  de  concorrência,  dependendo  da  situação;  muitas  outras  são  possíveis.  
Este  capítulo  se  concentra  em  uma  técnica  amplamente  utilizada:  o  bloqueio.  Um  bloqueio  proporciona  exclusão  mútua,  
garantindo  que  apenas  uma  CPU  por  vez  possa  manter  o  bloqueio.  Se  o  programador  associar  um  bloqueio  a  cada  item  
de  dados  compartilhado,  e  o  código  sempre  mantiver  o  bloqueio  associado  ao  usar  um  item,  então  o  item  será  usado  
por  apenas  uma  CPU  por  vez.  Nessa  situação,  dizemos  que  o  bloqueio  protege  o  item  de  dados.  Embora  os  bloqueios  
sejam  um  mecanismo  de  controle  de  concorrência  fácil  de  entender,  a  desvantagem  dos  bloqueios  é  que  eles  podem  
limitar  o  desempenho,  pois  serializam  operações  simultâneas.
 O  restante  deste  capítulo  explica  por  que  o  xv6  precisa  de  bloqueios,  como  o  xv6  os  implementa  e  como  os  utiliza.
 59
Machine Translated by Google
 6.1  Corridas
 Memória
 lista
 l->próximo  =  lista
 CPU
 Figura  6.1:  Arquitetura  SMP  simplificada
 ÔNIBUS
 l->próximo  =  lista
 CPU
 Como  exemplo  de  por  que  precisamos  de  bloqueios,  considere  dois  processos  com  filhos  que  saíram  chamando  wait  em
 duas  CPUs  diferentes.  wait  libera  a  memória  da  criança.  Assim,  em  cada  CPU,  o  kernel  chamará  kfree
 para  liberar  as  páginas  de  memória  das  crianças.  O  alocador  do  kernel  mantém  uma  lista  encadeada:  kalloc()  
(kernel /kalloc.c:69)  extrai  uma  página  de  memória  de  uma  lista  de  páginas  livres  e  kfree()  (kernel/kalloc.c:47)
 coloca  uma  página  na  lista  de  livres.  Para  melhor  desempenho,  podemos  esperar  que  os  kfrees  da
 dois  processos  pais  seriam  executados  em  paralelo  sem  que  nenhum  deles  tivesse  que  esperar  pelo  outro,  mas  isso
 não  seria  correto  dada  a  implementação  kfree  do  xv6.
 A  Figura  6.1  ilustra  a  configuração  com  mais  detalhes:  a  lista  vinculada  de  páginas  livres  está  na  memória  que  é
 compartilhado  pelas  duas  CPUs,  que  manipulam  a  lista  usando  instruções  de  carga  e  armazenamento.  (Na  realidade,  o
 os  processadores  têm  caches,  mas  conceitualmente  os  sistemas  multiprocessadores  se  comportam  como  se  houvesse  um  único,
 memória  compartilhada.)  Se  não  houver  solicitações  simultâneas,  você  pode  implementar  uma  operação  de  envio  de  lista
 do  seguinte  modo:
 1
 2
 3
 4
 5
 6
 elemento  struct  {
 };
 dados  int;
 struct  elemento  *próximo;
 struct  elemento  *lista  =  0;
 7
 8
 9
 10
 11
 12
 13
 14
 15
 16
 vazio
 push(int  dados)
 {
 struct  elemento  *l;
 l  
=  malloc(tamanho  de  *l);
 l->dados  =  dados;
 l->next  =  lista;
 lista  =  l;
 60
Machine Translated by Google
 CPU  1
 Memória
 CPU2
 17
 }
 15
 15
 16
 l->próximo
 l->próximo
 16
 Figura  6.2:  Exemplo  de  corrida
 lista
 lista
 Tempo
 Esta  implementação  está  correta  se  executada  isoladamente.  No  entanto,  o  código  não  está  correto  se  mais
 mais  de  uma  cópia  é  executada  simultaneamente.  Se  duas  CPUs  executarem  push  ao  mesmo  tempo,  ambas  podem
 execute  a  linha  15  conforme  mostrado  na  Fig.  6.1,  antes  de  executar  a  linha  16,  o  que  resulta  em  um  erro
 resultado  conforme  ilustrado  pela  Figura  6.2.  Haveria  então  dois  elementos  de  lista  com  o  próximo  conjunto  para  o
 valor  anterior  da  lista.  Quando  as  duas  atribuições  à  lista  acontecem  na  linha  16,  a  segunda  será
 substituir  o  primeiro;  o  elemento  envolvido  na  primeira  atribuição  será  perdido.
 A  atualização  perdida  na  linha  16  é  um  exemplo  de  corrida.  Uma  corrida  é  uma  situação  em  que  uma  memória
 o  local  é  acessado  simultaneamente  e  pelo  menos  um  acesso  é  de  gravação.  Uma  corrida  geralmente  é  um  sinal  de  bug,
 uma  atualização  perdida  (se  os  acessos  forem  gravações)  ou  uma  leitura  de  uma  estrutura  de  dados  atualizada  de  forma  incompleta.
 O  resultado  de  uma  corrida  depende  do  código  de  máquina  gerado  pelo  compilador,  do  tempo  de
 as  duas  CPUs  envolvidas  e  como  suas  operações  de  memória  são  ordenadas  pelo  sistema  de  memória,
 o  que  pode  dificultar  a  reprodução  e  a  depuração  de  erros  induzidos  por  corrida.  Por  exemplo,  adicionar  print
 declarações  durante  a  depuração  do  push  podem  alterar  o  tempo  de  execução  o  suficiente  para  fazer  o
 raça  desaparece.
 A  maneira  usual  de  evitar  corridas  é  usar  um  cadeado.  Os  cadeados  garantem  a  exclusão  mútua,  de  modo  que  apenas  um
 A  CPU  pode  executar  as  linhas  sensíveis  do  push  por  vez ;  isso  torna  o  cenário  acima  impossível.
 A  versão  corretamente  bloqueada  do  código  acima  adiciona  apenas  algumas  linhas  (destacadas  em  amarelo):
 6
 7
 8
 9  10
 11
 12
 13
 14
 struct  elemento  *lista  =  0;
 struct  lock  lista  de  bloqueio;
 vazio
 push(int  dados)
 {
 struct  elemento  *l;
 l  
=  malloc(tamanho  de  *l);
 l->dados  =  dados;
 61
Machine Translated by Google
 15
 16
 17
 18
 19
 20
 adquirir(&listlock);  l->próximo  =  
lista;  lista  =  l;  
liberar(&listlock);
 }
 A  sequência  de  instruções  entre  aquisição  e  liberação  é  frequentemente  chamada  de  seção  crítica.
 Geralmente,  diz-se  que  o  bloqueio  está  protegendo  a  lista.
 Quando  dizemos  que  um  bloqueio  protege  dados,  na  verdade  queremos  dizer  que  o  bloqueio  protege  alguma  coleção  
de  invariantes  que  se  aplicam  aos  dados.  Invariantes  são  propriedades  de  estruturas  de  dados  que  são  mantidas  entre  
operações.  Normalmente,  o  comportamento  correto  de  uma  operação  depende  de  as  invariantes  serem  verdadeiras  
quando  a  operação  começa.  A  operação  pode  violar  temporariamente  as  invariantes,  mas  deve  restabelecê-las  antes  de  
terminar.  Por  exemplo,  no  caso  da  lista  encadeada,  a  invariante  é  que  a  lista  aponta  para  o  primeiro  elemento  da  lista  e  
que  o  próximo  campo  de  cada  elemento  aponta  para  o  próximo  elemento.  A  implementação  de  push  viola  essa  invariante  
temporariamente:  na  linha  17,  l  aponta  para  o  próximo  elemento  da  lista,  mas  list  ainda  não  aponta  para  l  (restabelecido  
na  linha  18).  A  corrida  que  examinamos  acima  ocorreu  porque  uma  segunda  CPU  executou  código  que  dependia  das  
invariantes  da  lista  enquanto  elas  estavam  (temporariamente)  violadas.  O  uso  adequado  de  um  bloqueio  garante  que  
apenas  uma  CPU  por  vez  possa  operar  na  estrutura  de  dados  na  seção  crítica,  de  modo  que  nenhuma  CPU  executará  
uma  operação  na  estrutura  de  dados  quando  as  invariantes  da  estrutura  de  dados  não  forem  válidas.
 Você  pode  pensar  em  um  bloqueio  como  a  serialização  de  seções  críticas  simultâneas  para  que  sejam  
executadas  uma  de  cada  vez,  preservando  assim  as  invariantes  (assumindo  que  as  seções  críticas  estejam  corretas  
isoladamente).  Você  também  pode  pensar  em  seções  críticas  protegidas  pelo  mesmo  bloqueio  como  atômicas  entre  
si,  de  modo  que  cada  uma  veja  apenas  o  conjunto  completo  de  alterações  das  seções  críticas  anteriores  e  nunca  
veja  atualizações  parcialmente  concluídas.
 Embora  úteis  para  a  correção,  os  bloqueios  limitam  inerentemente  o  desempenho.  Por  exemplo,  se  dois  
processos  chamarem  kfree  simultaneamente,  os  bloqueios  serializarão  as  duas  seções  críticas,  de  modo  que  não  
há  benefício  em  executá-los  em  CPUs  diferentes.  Dizemos  que  múltiplos  processos  entram  em  conflito  se  desejam  
o  mesmo  bloqueio  ao  mesmo  tempo  ou  se  o  bloqueio  sofre  contenção.  Um  grande  desafio  no  projeto  do  kernel  é  
evitar  a  contenção  de  bloqueios  em  busca  do  paralelismo.  O  Xv6  faz  pouco  disso,  mas  kernels  sofisticados  
organizam  estruturas  de  dados  e  algoritmos  especificamente  para  evitar  a  contenção  de  bloqueios.
 No  exemplo  da  lista,  um  kernel  pode  manter  uma  lista  livre  separada  por  CPU  e  somente  tocar  na  lista  livre  de  outra  
CPU  se  a  lista  da  CPU  atual  estiver  vazia  e  ele  precisar  roubar  memória  de  outra  CPU.
 Outros  casos  de  uso  podem  exigir  designs  mais  complicados.
 O  posicionamento  dos  bloqueios  também  é  importante  para  o  desempenho.  Por  exemplo,  seria  correto  mover  a  
aquisição  para  o  início  do  push,  antes  da  linha  13.  Mas  isso  provavelmente  reduziria  o  desempenho,  pois  as  
chamadas  para  malloc  seriam  serializadas.  A  seção  "Usando  bloqueios"  abaixo  fornece  algumas  diretrizes  sobre  
onde  inserir  as  invocações  acquire  e  release .
 62
Machine Translated by Google
 6.2  Código:  Fechaduras
 Xv6  possui  dois  tipos  de  travas:  spinlocks  e  sleep-locks.  Começaremos  com  spinlocks.  Xv6  representa
 um  spinlock  como  uma  estrutura  spinlock  (kernel/spinlock.h:2).  O  campo  importante  na  estrutura  é
 bloqueado,  uma  palavra  que  é  zero  quando  o  bloqueio  está  disponível  e  diferente  de  zero  quando  é  mantido.  Logicamente,  xv6
 deve  adquirir  um  bloqueio  executando  um  código  como
 21
 22
 23
 24
 25
 26
 27
 28  
29  
30
 vazio
 acquire(struct  spinlock  *lk) //  não  funciona!
 {
 }
 para(;;)  {
 se(lk->bloqueado  ==  0)  {
 lk->bloqueado  =  1;
 quebrar;
 }
 }
 Infelizmente,  esta  implementação  não  garante  a  exclusão  mútua  em  um  multiprocessador.
 pode  acontecer  que  duas  CPUs  cheguem  simultaneamente  na  linha  25,  veja  que  lk->locked  é  zero,  e  então
 ambos  pegam  o  bloqueio  executando  a  linha  26.  Neste  ponto,  duas  CPUs  diferentes  seguram  o  bloqueio,  o  que
 viola  a  propriedade  de  exclusão  mútua.  O  que  precisamos  é  de  uma  maneira  de  fazer  as  linhas  25  e  26  serem  executadas  como
 uma  etapa  atômica  (ou  seja,  indivisível).
 Como  os  bloqueios  são  amplamente  utilizados,  os  processadores  multi-core  geralmente  fornecem  instruções  que  
implementam  uma  versão  atômica  das  linhas  25  e  26.  No  RISC-V,  esta  instrução  é  amoswap  r,  a.
 amoswap  lê  o  valor  no  endereço  de  memória  a,  escreve  o  conteúdo  do  registrador  r  nesse  endereço,
 e  coloca  o  valor  lido  em  r.  Ou  seja,  ele  troca  o  conteúdo  do  registrador  e  da  memória
 endereço.  Ele  executa  essa  sequência  atomicamente,  usando  hardware  especial  para  evitar  qualquer  outra  CPU
 de  usar  o  endereço  de  memória  entre  a  leitura  e  a  gravação.
 Aquisição  do  Xv6  (kernel/spinlock.c:22)  usa  a  chamada  de  biblioteca  C  portátil  __sync_lock_test_and_set,
 que  se  resume  à  instrução  amoswap ;  o  valor  de  retorno  é  o  conteúdo  antigo  (trocado)  de
 lk->bloqueado.  A  função  acquire  envolve  a  troca  em  um  loop,  tentando  novamente  (girando)  até  que  tenha
 adquiriu  o  bloqueio.  Cada  iteração  troca  um  para  lk->locked  e  verifica  o  valor  anterior;  se
 o  valor  anterior  é  zero,  então  adquirimos  o  bloqueio  e  a  troca  terá  definido  lk->locked
 para  um.  Se  o  valor  anterior  for  um,  então  alguma  outra  CPU  detém  o  bloqueio,  e  o  fato  de  que
 troquei  um  atomicamente  para  lk->locked  e  não  alterei  seu  valor.
 Depois  que  o  bloqueio  é  adquirido,  o  acquire  registra,  para  depuração,  a  CPU  que  adquiriu  o  bloqueio.
 O  campo  lk->cpu  é  protegido  pelo  bloqueio  e  só  deve  ser  alterado  enquanto  o  bloqueio  estiver  pressionado.
 A  função  release  (kernel/spinlock.c:47)  é  o  oposto  de  adquirir:  limpa  o  lk->cpu
 campo  e  então  libera  o  bloqueio.  Conceitualmente,  a  liberação  requer  apenas  atribuir  zero  a  lk->locked.
 O  padrão  C  permite  que  os  compiladores  implementem  uma  atribuição  com  várias  instruções  de  armazenamento,  portanto,  um
 A  atribuição  em  C  pode  não  ser  atômica  em  relação  ao  código  concorrente.  Em  vez  disso,  a  liberação  usa  o  C
 função  de  biblioteca  __sync_lock_release  que  realiza  uma  atribuição  atômica.  Esta  função  também
 resume-se  a  uma  instrução  RISC-V  amoswap .
 63
Machine Translated by Google
 6.3  Código:  Usando  bloqueios
 O  Xv6  usa  travas  em  muitos  lugares  para  evitar  corridas.  Como  descrito  acima,  kalloc  (kernel/kalloc.c:69)  e  kfree  (kernel/
 kalloc.c:47)  formam  um  bom  exemplo.  Experimente  os  Exercícios  1  e  2  para  ver  o  que  acontece  se  essas  funções  omitirem  
os  bloqueios.  Você  provavelmente  descobrirá  que  é  difícil  desencadear  comportamentos  incorretos,  o  que  sugere  que  é  difícil  
testar  de  forma  confiável  se  o  código  está  livre  de  erros  de  bloqueio  e  corridas.  O  Xv6  pode  muito  bem  ter  corridas  ainda  não  
descobertas.
 Uma  parte  complexa  do  uso  de  bloqueios  é  decidir  quantos  bloqueios  usar  e  quais  dados  e  invariantes  cada  bloqueio  
deve  proteger.  Existem  alguns  princípios  básicos.  Primeiro,  sempre  que  uma  variável  puder  ser  escrita  por  uma  CPU  ao  
mesmo  tempo  em  que  outra  CPU  pode  lê-la  ou  escrevê-la,  um  bloqueio  deve  ser  usado  para  evitar  que  as  duas  operações  se  
sobreponham.  Segundo,  lembre-se  de  que  os  bloqueios  protegem  invariantes:  se  uma  invariante  envolve  vários  locais  de  
memória,  normalmente  todos  eles  precisam  ser  protegidos  por  um  único  bloqueio  para  garantir  que  a  invariante  seja  mantida.
 As  regras  acima  dizem  quando  os  bloqueios  são  necessários,  mas  não  dizem  nada  sobre  quando  eles  são  
desnecessários,  e  é  importante  para  a  eficiência  não  bloquear  demais,  porque  os  bloqueios  reduzem  o  paralelismo.  Se  o  
paralelismo  não  for  importante,  pode-se  organizar  para  ter  apenas  uma  única  thread  e  não  se  preocupar  com  bloqueios.  Um  
kernel  simples  pode  fazer  isso  em  um  multiprocessador  tendo  um  único  bloqueio  que  deve  ser  adquirido  ao  entrar  no  kernel  
e  liberado  ao  sair  do  kernel  (embora  bloquear  chamadas  de  sistema,  como  leituras  de  pipe  ou  espera ,  representasse  um  
problema).  Muitos  sistemas  operacionais  uniprocessadores  foram  convertidos  para  rodar  em  multiprocessadores  usando  essa  
abordagem,  às  vezes  chamada  de  "bloqueio  de  kernel  grande",  mas  a  abordagem  sacrifica  o  paralelismo:  apenas  uma  CPU  
pode  executar  no  kernel  por  vez.  Se  o  kernel  fizer  alguma  computação  pesada,  seria  mais  eficiente  usar  um  conjunto  maior  
de  bloqueios  mais  refinados,  para  que  o  kernel  pudesse  executar  em  várias  CPUs  simultaneamente.
 Como  exemplo  de  bloqueio  de  granulação  grossa,  o  alocador  kalloc.c  do  xv6  possui  uma  única  lista  de  páginas  livres  
protegida  por  um  único  bloqueio.  Se  vários  processos  em  CPUs  diferentes  tentarem  alocar  páginas  simultaneamente,  cada  
um  terá  que  aguardar  sua  vez,  girando  em  aquisição.  Girar  desperdiça  tempo  de  CPU,  pois  não  é  um  trabalho  útil.  Se  a  
contenção  pelo  bloqueio  desperdiçasse  uma  fração  significativa  do  tempo  de  CPU,  talvez  o  desempenho  pudesse  ser  
melhorado  alterando  o  design  do  alocador  para  ter  várias  listas  de  páginas  livres,  cada  uma  com  seu  próprio  bloqueio,  para  
permitir  uma  alocação  verdadeiramente  paralela.
 Como  exemplo  de  bloqueio  de  granularidade  fina,  o  xv6  possui  um  bloqueio  separado  para  cada  arquivo,  de  modo  que  
processos  que  manipulam  arquivos  diferentes  podem  frequentemente  prosseguir  sem  esperar  pelos  bloqueios  uns  dos  outros.  
O  esquema  de  bloqueio  de  arquivos  poderia  ser  ainda  mais  granular  se  fosse  necessário  permitir  que  processos  gravassem  
simultaneamente  diferentes  áreas  do  mesmo  arquivo.  Em  última  análise,  as  decisões  sobre  a  granularidade  do  bloqueio  
precisam  ser  orientadas  por  medições  de  desempenho,  bem  como  por  considerações  de  complexidade.
 À  medida  que  os  capítulos  subsequentes  explicam  cada  parte  do  xv6,  eles  mencionarão  exemplos  do  uso  de  xv6  de
 bloqueios  para  lidar  com  a  simultaneidade.  Como  prévia,  a  Figura  6.3  lista  todos  os  bloqueios  no  xv6.
 6.4  Deadlock  e  ordenação  de  bloqueio
 Se  um  caminho  de  código  no  kernel  precisar  conter  vários  bloqueios  simultaneamente,  é  importante  que  todos  os  caminhos  
de  código  adquiram  esses  bloqueios  na  mesma  ordem.  Caso  contrário,  há  risco  de  deadlock.  Digamos  que  dois  caminhos  de  
código  no  xv6  precisem  dos  bloqueios  A  e  B,  mas  o  caminho  de  código  1  adquire  bloqueios  na  ordem  A  e  B.
 64
Machine Translated by Google
 Trancar
 bcache.lock  
cons.lock  
ftable.lock  
itable.lock  
vdisk_lock  
kmem.lock
 Descrição
 Protege  a  alocação  de  entradas  de  cache  do  buffer  de  bloco
 Serializa  o  acesso  ao  hardware  do  console,  evita  saídas  misturadas
 Serializa  a  alocação  de  um  arquivo  struct  na  tabela  de  arquivos
 Protege  a  alocação  de  entradas  de  inode  na  memória
 Serializa  o  acesso  ao  hardware  do  disco  e  à  fila  de  descritores  DMA
 Serializa  a  alocação  de  memória
 log.lock  Serializa  operações  no  log  de  transações
 pi->lock  do  pipe  Serializa  operações  em  cada  pipe
 pid_lock  Serializa  incrementos  de  next_pid
 proc's  p->lock  Serializa  as  alterações  no  estado  do  processo
 wait_lock  Ajuda  a  evitar  despertares  perdidos
 tickslock  Serializa  operações  no  contador  de  ticks
 ip->lock  do  inode  Serializa  operações  em  cada  inode  e  seu  conteúdo
 bloqueio  de  buf
 Serializa  as  operações  em  cada  buffer  de  bloco
 Figura  6.3:  Bloqueios  em  xv6
 B,  e  o  outro  caminho  os  adquire  na  ordem  B  e  depois  A.  Suponha  que  o  thread  T1  execute  o  caminho  de  código  1
 e  adquire  o  bloqueio  A,  e  o  thread  T2  executa  o  caminho  de  código  2  e  adquire  o  bloqueio  B.  O  próximo  T1  tentará
 adquirir  o  bloqueio  B,  e  T2  tentará  adquirir  o  bloqueio  A.  Ambos  os  adquirentes  bloquearão  indefinidamente,  porque  em
 em  ambos  os  casos,  a  outra  thread  mantém  o  bloqueio  necessário  e  não  o  liberará  até  que  sua  aquisição  retorne.
 Para  evitar  tais  impasses,  todos  os  caminhos  de  código  devem  adquirir  bloqueios  na  mesma  ordem.  A  necessidade  de  uma  estrutura  global
 A  ordem  de  aquisição  de  bloqueios  significa  que  os  bloqueios  são  efetivamente  parte  da  especificação  de  cada  função:  chamadores
 deve  invocar  funções  de  forma  que  os  bloqueios  sejam  adquiridos  na  ordem  acordada.
 O  Xv6  possui  muitas  cadeias  de  ordem  de  bloqueio  de  comprimento  dois  envolvendo  bloqueios  por  processo  (o  bloqueio  em  cada
 struct  proc)  devido  à  maneira  como  o  sleep  funciona  (veja  o  Capítulo  7).  Por  exemplo,  consoleintr
 (kernel/console.c:136)  é  a  rotina  de  interrupção  que  lida  com  caracteres  digitados.  Quando  uma  nova  linha  chega,  qualquer  processo  
que  esteja  aguardando  uma  entrada  do  console  deve  ser  ativado.  Para  fazer  isso,  consoleintr
 mantém  cons.lock  enquanto  chama  wakeup,  que  adquire  o  bloqueio  do  processo  de  espera  para
 Acorde-o.  Consequentemente,  a  ordem  de  bloqueio  global  para  evitar  deadlock  inclui  a  regra  de  que  cons.lock
 deve  ser  adquirido  antes  de  qualquer  bloqueio  de  processo.  O  código  do  sistema  de  arquivos  contém  as  cadeias  de  bloqueio  mais  longas  do  xv6.
 Por  exemplo,  a  criação  de  um  arquivo  requer  a  manutenção  simultânea  de  um  bloqueio  no  diretório  e  um  bloqueio  no  diretório.
 inode  do  novo  arquivo,  um  bloqueio  em  um  buffer  de  bloco  de  disco,  o  vdisk_lock  do  driver  de  disco  e  o  p->lock  do  processo  de  
chamada .  Para  evitar  deadlock,  o  código  do  sistema  de  arquivos  sempre  adquire  bloqueios  na  ordem  mencionada
 na  frase  anterior.
 Cumprir  uma  ordem  global  de  prevenção  de  impasses  pode  ser  surpreendentemente  difícil.  Às  vezes,  o  bloqueio
 conflitos  de  ordem  com  a  estrutura  lógica  do  programa,  por  exemplo,  talvez  o  módulo  de  código  M1  chame  o  módulo  M2,  mas
 a  ordem  de  bloqueio  exige  que  um  bloqueio  em  M2  seja  adquirido  antes  de  um  bloqueio  em  M1.  Às  vezes,  as  identidades
 de  fechaduras  não  são  conhecidas  com  antecedência,  talvez  porque  uma  fechadura  deva  ser  mantida  para  descobrir  a
 identidade  do  bloqueio  a  ser  adquirido  em  seguida.  Esse  tipo  de  situação  surge  no  sistema  de  arquivos  quando  ele  procura
 componentes  sucessivos  em  um  nome  de  caminho  e  no  código  para  esperar  e  sair  enquanto  pesquisam  na  tabela
 65
Machine Translated by Google
 de  processos  procurando  por  processos  filhos.  Por  fim,  o  perigo  de  deadlock  costuma  ser  uma  restrição  à  precisão  de  
um  esquema  de  bloqueio,  já  que  mais  bloqueios  geralmente  significam  mais  oportunidades  de  deadlock.  A  necessidade  
de  evitar  deadlocks  costuma  ser  um  fator  importante  na  implementação  do  kernel.
 6.5  Eclusas  reentrantes
 Pode  parecer  que  alguns  deadlocks  e  desafios  de  ordenação  de  bloqueios  poderiam  ser  evitados  com  o  uso  de  
bloqueios  reentrantes,  também  chamados  de  bloqueios  recursivos.  A  ideia  é  que,  se  o  bloqueio  for  mantido  por  um  
processo  e  esse  processo  tentar  obtê-lo  novamente,  o  kernel  poderia  simplesmente  permitir  isso  (já  que  o  processo  já  
possui  o  bloqueio),  em  vez  de  acionar  o  pânico,  como  faz  o  kernel  xv6.
 Acontece,  no  entanto,  que  bloqueios  reentrantes  dificultam  o  raciocínio  sobre  concorrência:  eles  quebram  a  
intuição  de  que  bloqueios  fazem  com  que  seções  críticas  sejam  atômicas  em  relação  a  outras  seções  críticas.  
Considere  as  duas  funções  a  seguir,  f  e  g:
 struct  spinlock  lock;  int  data  =  0; //  
protegido  por  bloqueio
 f()  
}
 g()  
}
 { adquirir(&bloquear);  
se(dados  ==  0)
 { chamar_uma  vez();  
h();  
dados  =  1;
 }  liberação(&bloqueio);
 { adquirir(&lock);  if(dados  
==  0){ chamar_uma  
vez();  dados  =  1;
 }  liberação(&bloqueio);
 Olhando  para  este  fragmento  de  código,  a  intuição  é  que  call_once  será  chamado  apenas  uma  vez:
 por  f,  ou  por  g,  mas  não  por  ambos.
 Mas  se  bloqueios  reentrantes  forem  permitidos,  e  h  chamar  g,  call_once  será  chamado  duas  vezes.
 Se  bloqueios  reentrantes  não  forem  permitidos,  então  h  chamando  g  resulta  em  um  deadlock,  o  que  também  não  
é  bom.  Mas,  supondo  que  chamar  call_once  seja  um  erro  grave ,  um  deadlock  é  preferível.  O  desenvolvedor  do  kernel  
observará  o  deadlock  (o  kernel  entra  em  pânico)  e  pode  corrigir  o  código  para  evitá-lo,  enquanto  chamar  call_once  
duas  vezes  pode  resultar  silenciosamente  em  um  erro  difícil  de  rastrear.
 Por  esse  motivo,  o  xv6  usa  bloqueios  não  reentrantes,  mais  fáceis  de  entender.  No  entanto,  desde  que  os  
programadores  mantenham  as  regras  de  bloqueio  em  mente,  qualquer  uma  das  abordagens  pode  funcionar.  Se  o  xv6  fosse...
 66
Machine Translated by Google
 Para  usar  bloqueios  reentrantes,  seria  necessário  modificar  a  aquisição  para  perceber  que  o  bloqueio  está  atualmente  sendo  
mantido  pela  thread  que  fez  a  chamada.  Também  seria  necessário  adicionar  uma  contagem  de  aquisições  aninhadas  à  struct  
spinlock,  de  forma  semelhante  a  push_off,  que  será  discutido  a  seguir.
 6.6  Bloqueios  e  manipuladores  de  interrupção
 Alguns  spinlocks  xv6  protegem  dados  usados  tanto  por  threads  quanto  por  manipuladores  de  interrupção.  Por  exemplo,  o  
manipulador  de  interrupção  do  timer  clockintr  pode  incrementar  ticks  (kernel/trap.c:164).  aproximadamente  ao  mesmo  tempo  
em  que  um  thread  do  kernel  lê  os  ticks  em  sys_sleep  (kernel/sysproc.c:59).  O  bloqueio  tickslock  serializa  os  dois  acessos.
 A  interação  de  spinlocks  e  interrupções  levanta  um  perigo  potencial.  Suponha  que  sys_sleep  esteja  com  o  tickslock  
bloqueado  e  sua  CPU  seja  interrompida  por  uma  interrupção  de  timer.  clockintr  tentaria  obter  o  tickslock,  verificar  se  ele  estava  
bloqueado  e  aguardar  sua  liberação.  Nessa  situação,  o  tickslock  nunca  será  liberado:  somente  sys_sleep  pode  liberá-lo,  mas  
sys_sleep  não  continuará  em  execução  até  que  clockintr  retorne.  Portanto,  a  CPU  entrará  em  deadlock  e  qualquer  código  que  
precise  de  qualquer  um  dos  bloqueios  também  travará.
 Para  evitar  essa  situação,  se  um  spinlock  for  usado  por  um  manipulador  de  interrupção,  uma  CPU  nunca  deve  manter  
esse  bloqueio  com  interrupções  habilitadas.  O  Xv6  é  mais  conservador:  quando  uma  CPU  adquire  qualquer  bloqueio,  o  xv6  
sempre  desabilita  as  interrupções  nessa  CPU.  Interrupções  ainda  podem  ocorrer  em  outras  CPUs,  portanto,  a  aquisição  de  
uma  interrupção  pode  esperar  que  uma  thread  libere  um  spinlock;  mas  não  na  mesma  CPU.
 O  Xv6  reativa  interrupções  quando  uma  CPU  não  possui  spinlocks;  ele  deve  fazer  uma  pequena  contabilidade  para  lidar  
com  seções  críticas  aninhadas.  acquire  chama  push_off  (kernel/spinlock.c:89)  e  libera  chamadas  pop_off  (kernel/spinlock.c:100)  
para  rastrear  o  nível  de  aninhamento  de  bloqueios  na  CPU  atual.  Quando  essa  contagem  chega  a  zero,  pop_off  restaura  o  
estado  de  ativação  de  interrupção  que  existia  no  início  da  seção  crítica  mais  externa.  As  funções  intr_off  e  intr_on  executam  
instruções  RISC-V  para  ativar  e  desativar  interrupções,  respectivamente.
 É  importante  adquirir  a  chamada  push_off  estritamente  antes  de  definir  lk->locked  (kernel/spin-  lock.c:28).  Se  as  duas  
opções  fossem  invertidas,  haveria  uma  breve  janela  de  tempo  em  que  o  bloqueio  seria  mantido  com  as  interrupções  habilitadas,  
e  uma  interrupção  com  tempo  de  execução  indesejável  causaria  um  deadlock  no  sistema.  Da  mesma  forma,  é  importante  
liberar  a  chamada  pop_off  somente  após  a  liberação  do  bloqueio  (kernel/spinlock.c:66).
 6.7  Instruções  e  ordenação  de  memória
 É  natural  pensar  em  programas  sendo  executados  na  ordem  em  que  as  instruções  do  código-fonte  aparecem.
 Este  é  um  modelo  mental  razoável  para  código  single-threaded,  mas  é  incorreto  quando  múltiplas  threads  interagem  através  
da  memória  compartilhada  [2,  4].  Um  motivo  é  que  os  compiladores  emitem  instruções  de  carregamento  e  armazenamento  em  
ordens  diferentes  daquelas  implícitas  no  código-fonte,  podendo  omiti-las  completamente  (por  exemplo,  armazenando  dados  em  
cache  em  registradores).  Outro  motivo  é  que  a  CPU  pode  executar  instruções  fora  de  ordem  para  aumentar  o  desempenho.  
Por  exemplo,  uma  CPU  pode  perceber  que,  em  uma  sequência  serial  de  instruções,  A  e  B  não  são  dependentes  uma  da  outra.  
A  CPU  pode  iniciar  a  instrução  B  primeiro,  seja  porque  suas  entradas  estão  prontas  antes  das  entradas  de  A,  seja  para  
sobrepor  a  execução  de  A  e  B.
 67
Machine Translated by Google
 Como  exemplo  do  que  poderia  dar  errado,  neste  código  para  push,  seria  um  desastre  se  o
 o  compilador  ou  a  CPU  moveu  o  armazenamento  correspondente  à  linha  4  para  um  ponto  após  o  lançamento  na  linha  6:
 2 
1
 3
 4
 5  6
 l  
=  malloc(sizeof  *l);  l->data  =  dados;  
adquirir(&listlock);  l
>next  =  lista;  lista  =  l;  
liberar(&listlock);
 Se  tal  reordenação  ocorresse,  haveria  uma  janela  durante  a  qual  outra  CPU  poderia  adquirir  o  bloqueio  e  observar  a  lista  
atualizada,  mas  veria  uma  lista  não  inicializada->próximo.
 A  boa  notícia  é  que  compiladores  e  CPUs  ajudam  programadores  simultâneos  seguindo  um  conjunto  de  regras  chamado  
modelo  de  memória  e  fornecendo  alguns  primitivos  para  ajudar  os  programadores  a  controlar  a  reordenação.
 Para  informar  ao  hardware  e  ao  compilador  para  não  reordenar,  o  xv6  usa  __sync_synchronize()  em  ambos  os  processos  
de  aquisição  (kernel/spinlock.c:22)  e  liberar  (kernel/spinlock.c:47).  __sync_synchronize()  é  uma  barreira  de  memória:  ela  
informa  ao  compilador  e  à  CPU  para  não  reordenarem  cargas  ou  armazenamentos  através  da  barreira.  As  barreiras  na  
aquisição  e  liberação  do  xv6  forçam  a  ordem  em  quase  todos  os  casos  em  que  isso  importa,  já  que  o  xv6  usa  bloqueios  em  
torno  de  acessos  a  dados  compartilhados.  O  Capítulo  9  discute  algumas  exceções.
 6.8  Bloqueios  de  sono
 Às  vezes,  o  xv6  precisa  manter  um  bloqueio  por  muito  tempo.  Por  exemplo,  o  sistema  de  arquivos  (Capítulo  8)  mantém  um  
arquivo  bloqueado  enquanto  lê  e  grava  seu  conteúdo  no  disco,  e  essas  operações  de  disco  podem  levar  dezenas  de  
milissegundos.  Manter  um  bloqueio  de  spin  por  tanto  tempo  levaria  a  desperdício  se  outro  processo  quisesse  adquiri-lo,  já  que  
o  processo  de  aquisição  consumiria  CPU  por  um  longo  tempo  enquanto  girava.  Outra  desvantagem  dos  bloqueios  de  spin  é  
que  um  processo  não  pode  ceder  a  CPU  enquanto  mantém  um  bloqueio  de  spin;  gostaríamos  de  fazer  isso  para  que  outros  
processos  possam  usar  a  CPU  enquanto  o  processo  com  o  bloqueio  aguarda  o  disco.
 Ceder  enquanto  mantém  um  spinlock  é  ilegal  porque  pode  levar  a  um  deadlock  se  uma  segunda  thread  tentar  adquirir  o  
spinlock;  como  a  aquisição  não  cede  a  CPU,  a  rotação  da  segunda  thread  pode  impedir  a  primeira  thread  de  executar  e  liberar  
o  bloqueio.  Ceder  enquanto  mantém  um  bloqueio  também  violaria  o  requisito  de  que  as  interrupções  devem  estar  desativadas  
enquanto  um  spinlock  é  mantido.  Portanto,  gostaríamos  de  um  tipo  de  bloqueio  que  ceda  a  CPU  enquanto  aguarda  para  
adquirir  e  permita  cedes  (e  interrupções)  enquanto  o  bloqueio  é  mantido.
 O  Xv6  fornece  tais  bloqueios  na  forma  de  bloqueios  de  sono.  acquiresleep  (kernel/sleeplock.c:22)  cede  a  CPU  enquanto  
aguarda,  usando  técnicas  que  serão  explicadas  no  Capítulo  7.  Em  um  nível  mais  alto,  um  sleep-lock  possui  um  campo  
bloqueado  que  é  protegido  por  um  spinlock,  e  a  chamada  de  acquiresleep  para  sleep  cede  atomicamente  a  CPU  e  libera  o  
spinlock.  O  resultado  é  que  outras  threads  podem  executar  enquanto  acquiresleep  aguarda.
 Como  os  bloqueios  de  suspensão  deixam  as  interrupções  habilitadas,  eles  não  podem  ser  usados  em  manipuladores  de  
interrupção.  Como  acquiresleep  pode  render  a  CPU,  os  bloqueios  de  suspensão  não  podem  ser  usados  dentro  de  seções  
críticas  de  spinlock  (embora  spinlocks  possam  ser  usados  dentro  de  seções  críticas  de  sleep-lock).
 68
Machine Translated by Google
 Os  spin-locks  são  mais  adequados  para  seções  críticas  curtas,  pois  esperar  por  eles  desperdiça  tempo  de  CPU;
 Os  bloqueios  de  sono  funcionam  bem  para  operações  demoradas.
 6.9  Mundo  real
 Programar  com  travas  continua  desafiador,  apesar  de  anos  de  pesquisa  sobre  primitivas  de  concorrência  e  paralelismo.  
Muitas  vezes,  é  melhor  ocultar  travas  em  construções  de  nível  superior,  como  filas  sincronizadas,  embora  o  xv6  não  faça  
isso.  Se  você  programa  com  travas,  é  aconselhável  usar  uma  ferramenta  que  tente  identificar  corridas,  pois  é  fácil  perder  
uma  invariante  que  exija  uma  trava.
 A  maioria  dos  sistemas  operacionais  suporta  threads  POSIX  (Pthreads),  que  permitem  que  um  processo  de  usuário  
tenha  várias  threads  em  execução  simultaneamente  em  diferentes  CPUs.  Pthreads  suportam  bloqueios,  barreiras,  etc.  em  
nível  de  usuário.  Pthreads  também  permitem  que  um  programador  especifique  opcionalmente  que  um  bloqueio  deve  ser  redefinido.
 participante.
 O  suporte  a  Pthreads  no  nível  do  usuário  requer  suporte  do  sistema  operacional.  Por  exemplo,  se  uma  pthread  for  
bloqueada  em  uma  chamada  de  sistema,  outra  pthread  do  mesmo  processo  poderá  ser  executada  naquela  CPU.  Outro  
exemplo:  se  uma  pthread  altera  o  espaço  de  endereço  do  seu  processo  (por  exemplo,  mapeia  ou  desmapeia  memória),  o  
kernel  deve  providenciar  para  que  outras  CPUs  que  executam  threads  do  mesmo  processo  atualizem  suas  tabelas  de  páginas  
de  hardware  para  refletir  a  alteração  no  espaço  de  endereço.
 É  possível  implementar  bloqueios  sem  instruções  atômicas  [10],  mas  é  caro  e  a  maioria
 sistemas  operacionais  usam  instruções  atômicas.
 Bloqueios  podem  ser  caros  se  muitas  CPUs  tentarem  adquirir  o  mesmo  bloqueio  ao  mesmo  tempo.  Se  uma  CPU  tiver  
um  bloqueio  armazenado  em  cache  em  seu  cache  local  e  outra  CPU  precisar  adquiri-lo,  a  instrução  atômica  para  atualizar  a  
linha  de  cache  que  contém  o  bloqueio  deve  mover  a  linha  do  cache  de  uma  CPU  para  o  cache  da  outra  CPU,  e  talvez  
invalidar  quaisquer  outras  cópias  da  linha  de  cache.  Buscar  uma  linha  de  cache  do  cache  de  outra  CPU  pode  ser  muito  mais  
caro  do  que  buscar  uma  linha  de  um  cache  local.
 Para  evitar  os  custos  associados  a  bloqueios,  muitos  sistemas  operacionais  utilizam  estruturas  de  dados  e  algoritmos  
livres  de  bloqueios  [6,  12].  Por  exemplo,  é  possível  implementar  uma  lista  encadeada  como  a  apresentada  no  início  do  
capítulo,  que  não  requer  bloqueios  durante  as  buscas  na  lista  e  uma  instrução  atômica  para  inserir  um  item  em  uma  lista.  A  
programação  livre  de  bloqueios  é  mais  complicada,  no  entanto,  do  que  programar  bloqueios;  por  exemplo,  é  preciso  se  
preocupar  com  a  reordenação  de  instruções  e  memória.  Programar  com  bloqueios  já  é  difícil,  portanto,  o  xv6  evita  a  
complexidade  adicional  da  programação  livre  de  bloqueios.

Capitulo: Traps e Chamadas de Sistema
Ha tres tipos de eventos que fazem a CPU interromper a execucao normal das instrucoes e transferir o
controle para um codigo especial que trata o evento. Um deles e uma chamada de sistema, quando um
programa em modo usuario executa a instrucao `ecall` para solicitar algo ao kernel. Outro e uma excecao:
quando uma instrucao faz algo ilegal, como dividir por zero ou acessar um endereco invalido. O terceiro tipo
e uma interrupcao de dispositivo, como quando o hardware do disco conclui uma leitura ou escrita.
Este livro usa o termo "trap" como uma referencia geral a esses eventos. A ideia e que o codigo que estava
sendo executado no momento da trap possa continuar depois, sem saber que foi interrompido. Traps devem
ser transparentes, especialmente as interrupcoes de dispositivos. A sequencia usual e: a trap transfere o
controle ao kernel; o kernel salva os registradores e o estado; executa o codigo apropriado (por exemplo,
chamada de sistema ou driver); restaura o estado salvo e retorna de onde parou.
O xv6 trata todas as traps no kernel. Isso e natural para chamadas de sistema e necessario para
interrupcoes, ja que o kernel e responsavel por controlar os dispositivos e garantir o isolamento entre
processos. Tambem e apropriado para excecoes: o xv6 termina qualquer processo que cause uma excecao.
O tratamento de traps no xv6 ocorre em quatro fases: acoes de hardware da CPU RISC-V, instrucoes em
assembly que preparam o caminho para o codigo C do kernel, uma funcao C que decide como tratar a trap,
e a rotina correspondente (chamada de sistema ou driver).
Cada CPU RISC-V tem registradores de controle (como `stvec`, `sepc`, `scause`, `sscratch`, `sstatus`) que
participam do processo de tratamento de traps. Quando ocorre uma trap, o hardware realiza uma sequencia
de etapas: desativa interrupcoes, salva o contador de programa, registra a causa, muda para o modo
supervisor e transfere o controle para o endereco em `stvec`.
No xv6, traps vindas do espaco do usuario seguem o fluxo: `uservec` `usertrap` `usertrapret` `userret`.
Uma limitacao do RISC-V e que ele nao troca a tabela de paginas automaticamente ao entrar no modo
supervisor. Para contornar isso, o xv6 usa uma pagina especial chamada *trampoline*, mapeada tanto na
tabela de paginas do usuario quanto do kernel.
O tratador `uservec` salva os 32 registradores do usuario no `trapframe` e muda a tabela de paginas para a
do kernel. Em seguida, `usertrap` identifica a causa da trap: se for chamada de sistema, chama `syscall`; se
for interrupcao, chama `devintr`; se for excecao, mata o processo.
Na volta, `usertrapret` prepara os registradores e chama `userret`, que restaura os registradores e executa
`sret` para voltar ao usuario.
Chamadas de sistema como `exec()` sao iniciadas com `ecall`. O numero da chamada vai em `a7` e os
argumentos nos registradores. O kernel busca esses valores no `trapframe` e executa a funcao
correspondente, salvando o valor de retorno em `a0`.
Para acessar ponteiros passados pelo usuario, o kernel usa funcoes como `copyin`, `copyout` e `copyinstr`,
que validam os enderecos e copiam dados entre os espacos de memoria do kernel e do usuario.
Traps originadas no kernel seguem um caminho diferente. O handler `kernelvec` salva os registradores na
pilha do kernel e chama `kerneltrap`, que trata interrupcoes e, se necessario, executa `yield`. Depois, os
registradores sao restaurados e a execucao e retomada com `sret`.
Em sistemas reais, traps e excecoes como falhas de pagina sao usadas para implementar recursos como
*copy-on-write fork*, *lazy allocation* (alocacao preguicosa), *demand paging* (paginacao sob demanda), e
*paging to disk* (paginacao em disco). Esses mecanismos aumentam a eficiencia e reduzem o uso de
memoria, permitindo que sistemas operacionais modernos escalem melhor.
Alem disso, recursos como pilhas autoexpansiveis e arquivos mapeados em memoria (`mmap`) sao
suportados atraves de excecoes de pagina. Sistemas de producao tambem fornecem chamadas como
`mmap`, `munmap`, `sigaction`, `mlock`, `madvise`, para controle avancado da memoria pelo usuario.
Por fim, embora estruturas como `trampoline` e `trapframe` possam parecer complexas, elas sao uma
consequencia do design minimalista da RISC-V, que deixa para o software grande parte da logica de
tratamento de traps, permitindo flexibilidade e desempenho.
Exercicios sugeridos incluem implementar alocacao preguicosa, fork com copia-sob-escrita, `mmap`, e
simplificar o uso de `TRAMPOLINE` e `TRAPFRAME`.

 Capítulo  3
 Tabelas  de  páginas
 Tabelas  de  páginas  são  o  mecanismo  mais  popular  pelo  qual  o  sistema  operacional  fornece  a  cada  processo  seu  próprio  
espaço  de  endereço  privado  e  memória.  Tabelas  de  páginas  determinam  o  significado  dos  endereços  de  memória  e  
quais  partes  da  memória  física  podem  ser  acessadas.  Elas  permitem  que  o  xv6  isole  os  espaços  de  endereço  de  diferentes  
processos  e  os  multiplexe  em  uma  única  memória  física.  Tabelas  de  páginas  são  um  design  popular  porque  fornecem  
um  nível  de  indireção  que  permite  aos  sistemas  operacionais  realizarem  muitos  truques.  O  xv6  realiza  alguns  truques:  
mapeia  a  mesma  memória  (uma  página  de  trampolim)  em  vários  espaços  de  endereço  e  protege  as  pilhas  do  kernel  e  do  
usuário  com  uma  página  não  mapeada.  O  restante  deste  capítulo  explica  as  tabelas  de  páginas  que  o  hardware  RISC-V  
fornece  e  como  o  xv6  as  utiliza.
 3.1  Hardware  de  paginação
 Vale  lembrar  que  as  instruções  RISC-V  (tanto  de  usuário  quanto  de  kernel)  manipulam  endereços  virtuais.  A  RAM  da  
máquina,  ou  memória  física,  é  indexada  com  endereços  físicos.  O  hardware  da  tabela  de  páginas  RISC-V  conecta  esses  
dois  tipos  de  endereços,  mapeando  cada  endereço  virtual  a  um  endereço  físico.
 27
 O  Xv6  roda  em  Sv39  RISC-V,  o  que  significa  que  apenas  os  39  bits  inferiores  de  um  endereço  virtual  de  64  bits  são  
usados;  os  25  bits  superiores  não  são  usados.  Nesta  configuração  Sv39,  uma  tabela  de  páginas  RISC-V  é  logicamente  
uma  matriz  de  2  (134.217.728)  entradas  de  tabela  de  páginas  (PTEs).  Cada  PTE  contém  um  número  de  página  física  
(PPN)  de  44  bits  e  alguns  sinalizadores.  O  hardware  de  paginação  traduz  um  endereço  virtual  usando  os  27  bits  superiores  
dos  39  bits  para  indexar  na  tabela  de  páginas  para  encontrar  um  PTE,  criando  um  endereço  físico  de  56  bits  cujos  44  bits  
superiores  vêm  do  PPN  no  PTE  e  cujos  12  bits  inferiores  são  copiados  do  endereço  virtual  original.  A  Figura  3.1  mostra  
esse  processo  com  uma  visão  lógica  da  tabela  de  páginas  como  uma  matriz  simples  de  PTEs  (veja  a  Figura  3.2  para  uma  
história  mais  completa).  Uma  tabela  de  páginas  fornece  ao  sistema  operacional  controle  sobre  as  traduções  de  endereços  
virtuais  para  físicos  na  granularidade  de  blocos  alinhados  de  4096  (212 )  bytes.  Esse  bloco  é  chamado  de  página.
 39
 No  Sv39  RISC-V,  os  25  bits  superiores  de  um  endereço  virtual  não  são  usados  para  tradução.  O  endereço  físico  
também  tem  espaço  para  crescimento:  no  formato  PTE,  há  espaço  para  que  o  número  de  páginas  físicas  aumente  em  
mais  10  bits.  Os  projetistas  do  RISC-V  escolheram  esses  números  com  base  em  previsões  tecnológicas.  2  bytes  equivalem  
a  512  GB,  o  que  deve  ser  espaço  de  endereço  suficiente  para  aplicativos  que  executam
 31
Machine Translated by Google
 Endereço  virtual
 25
 64
 EXT
 27
 Índice
 1
 0
 12
 Desvio
 44
 10
 Bandeiras  PPN
 Tabela  de  páginas
 2^27
 56
 44
 Endereço  físico
 12
 Figura  3.1:  Endereços  virtuais  e  físicos  RISC-V,  com  uma  tabela  de  páginas  lógicas  simplificada.
 56
 em  computadores  RISC-V.  2  é  espaço  de  memória  física  suficiente  para,  em  um  futuro  próximo,  acomodar  muitos  dispositivos  
de  E/S  e  chips  DRAM.  Se  mais  for  necessário,  os  projetistas  do  RISC-V  definiram  Sv48  com  endereços  virtuais  de  48  bits  [3].
 Como  mostra  a  Figura  3.2,  uma  CPU  RISC-V  traduz  um  endereço  virtual  em  um  físico  em  três  etapas.
 Uma  tabela  de  páginas  é  armazenada  na  memória  física  como  uma  árvore  de  três  níveis.  A  raiz  da  árvore  é  uma  página  de  
tabela  de  páginas  de  4.096  bytes  que  contém  512  PTEs,  que  contêm  os  endereços  físicos  das  páginas  da  tabela  de  páginas  no  
próximo  nível  da  árvore.  Cada  uma  dessas  páginas  contém  512  PTEs  para  o  nível  final  da  árvore.
 O  hardware  de  paginação  usa  os  9  bits  superiores  dos  27  bits  para  selecionar  uma  PTE  na  página  da  tabela  de  páginas  raiz,  os  
9  bits  do  meio  para  selecionar  uma  PTE  em  uma  página  da  tabela  de  páginas  no  próximo  nível  da  árvore  e  os  9  bits  inferiores  
para  selecionar  a  PTE  final.  (No  Sv48  RISC-V,  uma  tabela  de  páginas  tem  quatro  níveis,  e  os  bits  39  a  47  de  um  índice  de  
endereço  virtual  no  nível  superior.)
 Se  qualquer  um  dos  três  PTEs  necessários  para  traduzir  um  endereço  não  estiver  presente,  o  hardware  de  paginação
 gera  uma  exceção  de  falha  de  página,  deixando  para  o  kernel  lidar  com  a  exceção  (veja  Capítulo  4).
 A  estrutura  de  três  níveis  da  Figura  3.2  permite  uma  maneira  eficiente  de  gravar  PTEs  em  termos  de  memória,  em  
comparação  com  o  design  de  nível  único  da  Figura  3.1.  No  caso  comum  em  que  grandes  intervalos  de  endereços  virtuais  não  
têm  mapeamentos,  a  estrutura  de  três  níveis  pode  omitir  diretórios  de  páginas  inteiros.  Por  exemplo,  se  um  aplicativo  usa  
apenas  algumas  páginas  começando  no  endereço  zero,  as  entradas  de  1  a  511  do  diretório  de  páginas  de  nível  superior  são  
inválidas  e  o  kernel  não  precisa  alocar  páginas  para  os  511  diretórios  de  páginas  intermediárias.  Além  disso,  o  kernel  também  
não  precisa  alocar  páginas  para  os  diretórios  de  páginas  de  nível  inferior  para  esses  511  diretórios  de  páginas  intermediárias.  
Portanto,  neste  exemplo,  o  design  de  três  níveis  economiza  511  páginas  para  diretórios  de  páginas  intermediárias  e  511  ×  512  
páginas  para  diretórios  de  páginas  de  nível  inferior.
 Embora  uma  CPU  percorra  a  estrutura  de  três  níveis  no  hardware  como  parte  da  execução  de  uma  instrução  de  
carregamento  ou  armazenamento,  uma  desvantagem  potencial  dos  três  níveis  é  que  a  CPU  deve  carregar  três  PTEs  da  
memória  para  executar  a  tradução  do  endereço  virtual  na  instrução  de  carregamento/armazenamento  para  um  endereço  físico.
 Para  evitar  o  custo  de  carregamento  de  PTEs  da  memória  física,  uma  CPU  RISC-V  armazena  em  cache  as  entradas  da  tabela  
de  páginas  em  um  Translation  Look-aside  Buffer  (TLB).
 32
Machine Translated by Google
 Endereço  virtual
 9
 EXT
 44
 511
 10
 Bandeiras  PPN
 1
 511
 L2
 44
 L1
 10
 Bandeiras  PPN
 0
 satp
 63
 Reservado
 1
 Diretório  de  páginas
 53
 0
 Diretório  de  páginas
 8910  07
 Número  da  página  física
 Endereço  físico
 9
 9
 12
 Deslocamento  L0
 44
 PPN
 12
 Desvio
 44
 511
 10
 Bandeiras  PPN
 1
 0
 6  5  4  3  
Você
 RSW
 DUXR
 G
 UM
 Diretório  de  páginas
 2 1
 C
 V
 V
 R
 C
 X
 Você-  Válido-  Legível-  Gravável-  Executável-  Usuário-G  Global
 A  -  Acessado
 UM
 D-  Acessado-  Sujo  (0  no  diretório  de  páginas)
 Reservado  para  software  de  supervisão
 Figura  3.2:  Detalhes  da  tradução  de  endereços  RISC-V.
 Cada  PTE  contém  bits  de  sinalização  que  informam  ao  hardware  de  paginação  como  o  endereço  virtual  associado  
pode  ser  usado.  PTE_V  indica  se  a  PTE  está  presente:  se  não  estiver  definida,  uma  referência  à  página  causa  uma  
exceção  (ou  seja,  não  é  permitida).  PTE_R  controla  se  as  instruções  podem  ler  na  página.  PTE_W  controla  se  as  
instruções  podem  gravar  na  página.  PTE_X  controla  se  a  CPU  pode  interpretar  o  conteúdo  da  página  como  instruções  
e  executá-las.
 PTE_U  controla  se  instruções  em  modo  usuário  têm  permissão  para  acessar  a  página;  se  PTE_U  não  estiver  definido,  
o  PTE  só  poderá  ser  usado  em  modo  supervisor.  A  Figura  3.2  mostra  como  tudo  funciona.  Os  sinalizadores  e  todas  as  
outras  estruturas  relacionadas  ao  hardware  da  página  são  definidos  em  (kernel/riscv.h).
 Para  instruir  uma  CPU  a  usar  uma  tabela  de  páginas,  o  kernel  precisa  gravar  o  endereço  físico  da  página  raiz  da  
tabela  de  páginas  no  registrador  satp.  Uma  CPU  traduzirá  todos  os  endereços  gerados  por  instruções  subsequentes  
usando  a  tabela  de  páginas  apontada  por  seu  próprio  satp.  Cada  CPU  possui  seu  próprio  satp,  permitindo  que  
diferentes  CPUs  executem  processos  diferentes,  cada  um  com  um  espaço  de  endereço  privado  descrito  por  sua  
própria  tabela  de  páginas.
 Normalmente,  um  kernel  mapeia  toda  a  memória  física  em  sua  tabela  de  páginas  para  que  possa  ler  e  escrever  
em  qualquer  local  da  memória  física  usando  instruções  de  carregamento/armazenamento.  Como  os  diretórios  de  
páginas  estão  na  memória  física,  o  kernel  pode  programar  o  conteúdo  de  uma  PTE  em  um  diretório  de  páginas  
escrevendo  no  endereço  virtual  da  PTE  usando  uma  instrução  de  armazenamento  padrão.
 Algumas  notas  sobre  termos.  Memória  física  refere-se  às  células  de  armazenamento  na  DRAM.  Um  byte  de  
memória  física  possui  um  endereço,  chamado  endereço  físico.  As  instruções  usam  apenas  endereços  virtuais,  que  o  
hardware  de  paginação  traduz  em  endereços  físicos  e,  em  seguida,  envia  ao  hardware  DRAM  para  leitura.
 33
Machine Translated by Google
 ou  armazenamento  de  gravação.  Ao  contrário  da  memória  física  e  dos  endereços  virtuais,  a  memória  virtual  não  é  um  espaço  físico
 objeto,  mas  se  refere  à  coleção  de  abstrações  e  mecanismos  que  o  kernel  fornece  para  gerenciar
 memória  física  e  endereços  virtuais.
 Endereços  virtuais
 MAXVA
 PHYSTOP  
(0x88000000)
 BASE  DO  
KERN  (0x80000000)
 0x10001000
 0x10000000
 0x0C000000
 0
 Trampolim
 Página  de  guarda
 Kstack  0
 Página  de  guarda
 Kstack  1
 ...
 Memória  livre
 Dados  do  kernel
 Texto  do  kernel
 Disco  VIRTIO
 UART0
 PLIC
 Endereços  físicos
 2^56-1
 Receita  médica--
RW---
RW
RW
RW
Receita  médica
 RW
RW
RW
0x02000000
 0x1000
 0
 Não  utilizado
 Memória  física  (RAM)
 Dispositivos  
de  E/S  não  utilizados  e  outros
 Disco  VIRTIO
 UART0
 PLIC
 CLINT
 Não  utilizado
 ROM  de  inicialização
 Não  utilizado
 Figura  3.3:  À  esquerda,  o  espaço  de  endereço  do  kernel  do  xv6.  RWX  refere-se  a  leitura,  gravação  e  execução  de  PTE
 permissões.  À  direita,  o  espaço  de  endereço  físico  RISC-V  que  o  xv6  espera  ver.
 3.2  Espaço  de  endereço  do  kernel
 O  Xv6  mantém  uma  tabela  de  páginas  por  processo,  descrevendo  o  espaço  de  endereço  do  usuário  de  cada  processo,  além  de  
uma  tabela  de  página  única  que  descreve  o  espaço  de  endereço  do  kernel.  O  kernel  configura  o  layout  do  seu  espaço  de  endereço  
para  obter  acesso  à  memória  física  e  a  vários  recursos  de  hardware  em  intervalos  previsíveis.
 Endereços  virtuais.  A  Figura  3.3  mostra  como  este  layout  mapeia  endereços  virtuais  do  kernel  para  endereços  físicos.  O  arquivo  
(kernel/memlayout.h)  declara  as  constantes  para  o  layout  de  memória  do  kernel  do  xv6.
 34
Machine Translated by Google
 O  QEMU  simula  um  computador  que  inclui  RAM  (memória  física)  começando  no  endereço  físico  0x80000000  e  
continuando  até  pelo  menos  0x88000000,  que  o  xv6  chama  de  PHYSTOP.
 A  simulação  do  QEMU  também  inclui  dispositivos  de  E/S,  como  uma  interface  de  disco.  O  QEMU  expõe  as  interfaces  do  
dispositivo  ao  software  como  registradores  de  controle  mapeados  na  memória,  localizados  abaixo  de  0x80000000  no  
espaço  de  endereço  físico.  O  kernel  pode  interagir  com  os  dispositivos  lendo/escrevendo  esses  endereços  físicos  especiais;  
tais  leituras  e  gravações  se  comunicam  com  o  hardware  do  dispositivo  em  vez  da  RAM.  O  Capítulo  4  explica  como  o  xv6  
interage  com  os  dispositivos.
 O  kernel  acessa  a  RAM  e  os  registradores  de  dispositivos  mapeados  na  memória  usando  "mapeamento  direto",  ou  
seja,  mapeando  os  recursos  em  endereços  virtuais  iguais  ao  endereço  físico.  Por  exemplo,  o  próprio  kernel  está  localizado  
em  KERNBASE=0x80000000  tanto  no  espaço  de  endereço  virtual  quanto  na  memória  física.  O  mapeamento  direto  
simplifica  o  código  do  kernel  que  lê  ou  grava  na  memória  física.  Por  exemplo,  quando  o  fork  aloca  memória  do  usuário  
para  o  processo  filho,  o  alocador  retorna  o  endereço  físico  dessa  memória;  o  fork  usa  esse  endereço  diretamente  como  
um  endereço  virtual  ao  copiar  a  memória  do  usuário  do  processo  pai  para  o  filho.
 Há  alguns  endereços  virtuais  do  kernel  que  não  são  mapeados  diretamente:
 •  A  página  do  trampolim.  Ela  é  mapeada  no  topo  do  espaço  de  endereço  virtual;  as  tabelas  de  páginas  do  usuário  
têm  esse  mesmo  mapeamento.  O  Capítulo  4  discute  o  papel  da  página  do  trampolim,  mas  vemos  aqui  um  caso  de  
uso  interessante  de  tabelas  de  páginas;  uma  página  física  (contendo  o  código  do  trampolim)  é  mapeada  duas  
vezes  no  espaço  de  endereço  virtual  do  kernel:  uma  vez  no  topo  do  espaço  de  endereço  virtual  e  outra  com  um  
mapeamento  direto.
 •  As  páginas  da  pilha  do  kernel.  Cada  processo  tem  sua  própria  pilha  do  kernel,  que  é  mapeada  em  uma  posição  alta  
para  que,  abaixo  dela,  o  xv6  possa  deixar  uma  página  de  guarda  não  mapeada.  A  PTE  da  página  de  guarda  é  
inválida  (ou  seja,  PTE_V  não  está  definida),  de  modo  que,  se  o  kernel  transbordar  uma  pilha  do  kernel,  
provavelmente  causará  uma  exceção  e  o  kernel  entrará  em  pânico.  Sem  uma  página  de  guarda,  uma  pilha  
transbordando  sobrescreveria  outra  memória  do  kernel,  resultando  em  operação  incorreta.  Uma  falha  de  pânico  é  preferível.
 Embora  o  kernel  utilize  suas  pilhas  por  meio  de  mapeamentos  de  alta  memória,  elas  também  são  acessíveis  ao  
kernel  por  meio  de  um  endereço  mapeado  diretamente.  Um  projeto  alternativo  poderia  ter  apenas  o  mapeamento  direto  
e  usar  as  pilhas  no  endereço  mapeado  diretamente.  Nesse  arranjo,  no  entanto,  fornecer  páginas  de  guarda  envolveria  o  
desmapeamento  de  endereços  virtuais  que,  de  outra  forma,  se  refeririam  à  memória  física,  o  que  dificultaria  o  uso.
 O  kernel  mapeia  as  páginas  do  trampolim  e  o  texto  do  kernel  com  as  permissões  PTE_R  e  PTE_X.  O  kernel  lê  e  
executa  instruções  dessas  páginas.  O  kernel  mapeia  as  outras  páginas  com  as  permissões  PTE_R  e  PTE_W,  para  que  
possa  ler  e  gravar  na  memória  dessas  páginas.  Os  mapeamentos  para  as  páginas  de  guarda  são  inválidos.
 3.3  Código:  criando  um  espaço  de  endereço
 A  maior  parte  do  código  xv6  para  manipular  espaços  de  endereço  e  tabelas  de  páginas  reside  em  vm.c  ( kernel/vm.c:1).  
A  estrutura  de  dados  central  é  pagetable_t,  que  na  verdade  é  um  ponteiro  para  um  RISC-V
 35
Machine Translated by Google
 Tabela  de  páginas  raiz;  uma  pagetable_t  pode  ser  a  tabela  de  páginas  do  kernel  ou  uma  das  tabelas  de  páginas  por  
processo.  As  funções  centrais  são  walk,  que  encontra  o  PTE  para  um  endereço  virtual,  e  mappages,  que  instala  PTEs  
para  novos  mapeamentos.  Funções  que  começam  com  kvm  manipulam  a  tabela  de  páginas  do  kernel;  funções  que  
começam  com  uvm  manipulam  uma  tabela  de  páginas  do  usuário;  outras  funções  são  usadas  para  ambos.  copyout  e  
copyin  copiam  dados  de  e  para  endereços  virtuais  do  usuário  fornecidos  como  argumentos  de  chamada  de  sistema;  elas  
estão  em  vm.c  porque  precisam  traduzir  explicitamente  esses  endereços  para  encontrar  a  memória  física  correspondente.
 No  início  da  sequência  de  inicialização,  o  principal  chama  kvminit  (kernel/vm.c:54)  para  criar  a  tabela  de  páginas  do  
kernel  usando  kvmmake  (kernel/vm.c:20).  Esta  chamada  ocorre  antes  que  o  xv6  habilite  a  paginação  no  RISC-V,  
portanto,  os  endereços  se  referem  diretamente  à  memória  física.  O  kvmmake  primeiro  aloca  uma  página  de  memória  
física  para  armazenar  a  página  da  tabela  de  páginas  raiz.  Em  seguida,  ele  chama  o  kvmmap  para  instalar  as  traduções  
necessárias  ao  kernel.  As  traduções  incluem  as  instruções  e  os  dados  do  kernel,  a  memória  física  até  o  PHYSTOP  e  os  
intervalos  de  memória  que,  na  verdade,  são  dispositivos.  proc_mapstacks  (kernel/proc.c:33)  aloca  uma  pilha  do  kernel  
para  cada  processo.  Ele  chama  o  kvmmap  para  mapear  cada  pilha  no  endereço  virtual  gerado  pelo  KSTACK,  o  que  
deixa  espaço  para  as  páginas  de  proteção  de  pilha  inválidas.
 kvmmap  (kernel/vm.c:132)  chama  mappages  (kernel/vm.c:143),  que  instala  mapeamentos  em  uma  tabela  de  
páginas  para  um  intervalo  de  endereços  virtuais  em  um  intervalo  correspondente  de  endereços  físicos.  Ele  faz  isso  
separadamente  para  cada  endereço  virtual  no  intervalo,  em  intervalos  de  página.  Para  cada  endereço  virtual  a  ser  
mapeado,  mappages  chama  walk  para  encontrar  o  endereço  da  PTE  correspondente  a  esse  endereço.  Em  
seguida,  ele  inicializa  a  PTE  para  conter  o  número  da  página  física  relevante,  as  permissões  desejadas  (PTE_W,  
PTE_X  e/ou  PTE_R)  e  PTE_V  para  marcar  a  PTE  como  válida  (kernel/vm.c:158).
 andar  (kernel/vm.c:86)  imita  o  hardware  de  paginação  RISC-V  enquanto  procura  no  PTE  por  um  endereço  virtual  
(veja  a  Figura  3.2).  O  comando  walk  desce  a  tabela  de  páginas  de  3  níveis  9  bits  por  vez.  Ele  usa  os  9  bits  de  endereço  
virtual  de  cada  nível  para  encontrar  o  PTE  da  tabela  de  páginas  do  próximo  nível  ou  da  página  final  (kernel/vm.c:92).  Se  
o  PTE  não  for  válido,  a  página  necessária  ainda  não  foi  alocada;  se  o  argumento  alloc  estiver  definido,  walk  aloca  uma  
nova  página  da  tabela  de  páginas  e  insere  seu  endereço  físico  no  PTE.  Ele  retorna  o  endereço  do  PTE  na  camada  mais  
baixa  da  árvore  (kernel/vm.c:102).
 O  código  acima  depende  da  memória  física  ser  mapeada  diretamente  no  espaço  de  endereço  virtual  do  kernel.  Por  
exemplo,  à  medida  que  o  walk  desce  níveis  da  tabela  de  páginas,  ele  extrai  o  endereço  (físico)  da  tabela  de  páginas  do  
nível  seguinte  de  uma  PTE  (kernel/vm.c:94).  e  então  usa  esse  endereço  como  um  endereço  virtual  para  buscar  o  PTE  
no  próximo  nível  abaixo  (kernel/vm.c:92).
 chamadas  principais  kvminithart  (kernel/vm.c:62)  para  instalar  a  tabela  de  páginas  do  kernel.  Ele  grava  o  
endereço  físico  da  página  raiz  da  tabela  de  páginas  no  registrador  satp.  Depois  disso,  a  CPU  traduzirá  os  
endereços  usando  a  tabela  de  páginas  do  kernel.  Como  o  kernel  usa  um  mapeamento  de  identidade,  o  endereço  
virtual  da  próxima  instrução  será  mapeado  para  o  endereço  de  memória  física  correto.
 Cada  CPU  RISC-V  armazena  em  cache  as  entradas  da  tabela  de  páginas  em  um  Translation  Look-aside  Buffer  
(TLB)  e,  quando  o  xv6  altera  uma  tabela  de  páginas,  deve  informar  à  CPU  para  invalidar  as  entradas  TLB  em  cache  
correspondentes.  Caso  contrário,  em  algum  momento  posterior,  o  TLB  poderá  usar  um  mapeamento  em  cache  antigo,  
apontando  para  uma  página  física  que,  nesse  meio  tempo,  foi  alocada  a  outro  processo  e,  como  resultado,  um  processo  
poderá  rabiscar  na  memória  de  outro  processo.  O  RISC-V  possui  uma  instrução  sfence.vma  que  limpa  o  TLB  da  CPU  
atual.  O  xv6  executa  sfence.vma  em  kvminithart  após  recarregar  o  registrador  satp  e  no  código  trampoline  que  alterna  
para  um
 36
Machine Translated by Google
 tabela  de  página  do  usuário  antes  de  retornar  ao  espaço  do  usuário  (kernel/trampoline.S:89).
 Também  é  necessário  emitir  sfence.vma  antes  de  alterar  o  satp,  para  aguardar  a  conclusão  de  todos  os  carregamentos  e  
armazenamentos  pendentes.  Essa  espera  garante  que  as  atualizações  anteriores  na  tabela  de  páginas  tenham  sido  concluídas  e  
que  os  carregamentos  e  armazenamentos  anteriores  usem  a  tabela  de  páginas  antiga,  não  a  nova.
 um.
 Para  evitar  a  limpeza  completa  do  TLB,  as  CPUs  RISC-V  podem  suportar  identificadores  de  espaço  de  endereço  
(ASIDs)  [3].  O  kernel  pode  então  limpar  apenas  as  entradas  TLB  para  um  espaço  de  endereço  específico.  O  Xv6  não  
utiliza  esse  recurso.
 3.4  Alocação  de  memória  física
 O  kernel  deve  alocar  e  liberar  memória  física  em  tempo  de  execução  para  tabelas  de  páginas,  memória  do  usuário,  pilhas  do  kernel  
e  buffers  de  pipe.
 O  Xv6  utiliza  a  memória  física  entre  o  final  do  kernel  e  o  PHYSTOP  para  alocação  em  tempo  de  execução.  Ele  aloca  e  libera  
páginas  inteiras  de  4.096  bytes  por  vez.  Ele  rastreia  quais  páginas  estão  livres,  encadeando  uma  lista  encadeada  entre  as  próprias  
páginas.  A  alocação  consiste  em  remover  uma  página  da  lista  encadeada;  a  liberação  consiste  em  adicionar  a  página  liberada  à  
lista.
 3.5  Código:  Alocador  de  memória  física
 O  alocador  reside  em  kalloc.c  (kernel/kalloc.c:1).  A  estrutura  de  dados  do  alocador  é  uma  lista  livre  de  páginas  de  memória  física  
disponíveis  para  alocação.  Cada  elemento  da  lista  de  páginas  livres  é  uma  execução  de  struct  (kernel/kalloc.c:17).  De  onde  o  
alocador  obtém  a  memória  para  armazenar  essa  estrutura  de  dados?  Ele  armazena  a  estrutura  de  execução  de  cada  página  livre  
na  própria  página  livre,  já  que  não  há  mais  nada  armazenado  lá.  A  lista  livre  é  protegida  por  um  bloqueio  de  spin  (kernel/
 kalloc.c:21-24).  A  lista  e  o  bloqueio  são  encapsulados  em  uma  struct  para  deixar  claro  que  o  bloqueio  protege  os  campos  na  struct.  
Por  enquanto,  ignore  o  bloqueio  e  as  chamadas  para  acquire  e  release;  o  Capítulo  6  examinará  o  bloqueio  em  detalhes.
 A  função  main  chama  kinit  para  inicializar  o  alocador  (kernel/kalloc.c:27).  O  kinit  inicializa  a  lista  de  páginas  livres  para  
armazenar  todas  as  páginas  entre  o  final  do  kernel  e  o  PHYSTOP.  O  Xv6  deve  determinar  quanta  memória  física  está  disponível  
analisando  as  informações  de  configuração  fornecidas  pelo  hardware.  Em  vez  disso,  o  xv6  assume  que  a  máquina  possui  128  
megabytes  de  RAM.  O  kinit  chama  o  freerange  para  adicionar  memória  à  lista  de  páginas  livres  por  meio  de  chamadas  por  página  
ao  kfree.  Uma  PTE  só  pode  se  referir  a  um  endereço  físico  alinhado  em  um  limite  de  4096  bytes  (é  um  múltiplo  de  4096),  portanto,  
o  freerange  usa  PGROUNDUP  para  garantir  que  libere  apenas  os  endereços  físicos  alinhados.  O  alocador  inicia  sem  memória;  
essas  chamadas  ao  kfree  lhe  dão  algum  trabalho  para  gerenciar.
 O  alocador  às  vezes  trata  endereços  como  inteiros  para  realizar  operações  aritméticas  com  eles  (por  exemplo,  percorrendo  
todas  as  páginas  em  modo  livre)  e,  às  vezes,  usa  endereços  como  ponteiros  para  ler  e  escrever  na  memória  (por  exemplo,  
manipulando  a  estrutura  de  execução  armazenada  em  cada  página);  esse  uso  duplo  de  endereços  é  o  principal  motivo  pelo  qual  o  
código  do  alocador  está  cheio  de  conversões  de  tipo  em  C.  O  outro  motivo  é  que  a  liberação  e  a  alocação  alteram  inerentemente  o  
tipo  de  memória.
 37
Machine Translated by Google
 A  função  kfree  (kernel/kalloc.c:47)  começa  definindo  cada  byte  na  memória  que  está  sendo  liberada  para  o  valor  1.  Isso  fará  
com  que  o  código  que  usa  memória  após  liberá-la  (usa  “referências  pendentes”)  leia  lixo  em  vez  do  antigo  conteúdo  válido;  
esperançosamente,  isso  fará  com  que  esse  código  quebre  mais  rápido.
 Em  seguida,  kfree  anexa  a  página  à  lista  livre:  ele  converte  pa  para  um  ponteiro  para  struct  run,  registra  o  início  antigo  
da  lista  livre  em  r->next  e  define  a  lista  livre  como  igual  a  r.  kalloc  remove  e  retorna  o  primeiro  elemento  na  lista  livre.
 3.6  Espaço  de  endereço  do  processo
 Cada  processo  possui  uma  tabela  de  páginas  separada  e,  quando  o  xv6  alterna  entre  processos,  também  altera  as  
tabelas  de  páginas.  A  Figura  3.4  mostra  o  espaço  de  endereço  de  um  processo  com  mais  detalhes  do  que  a  Figura  2.3.  
A  memória  de  usuário  de  um  processo  começa  no  endereço  virtual  zero  e  pode  crescer  até  MAXVA  (kernel/riscv.h:360).  
permitindo  que  um  processo  enderece  em  princípio  256  Gigabytes  de  memória.
 O  espaço  de  endereço  de  um  processo  consiste  em  páginas  que  contêm  o  texto  do  programa  (que  o  xv6  mapeia  
com  as  permissões  PTE_R,  PTE_X  e  PTE_U),  páginas  que  contêm  os  dados  pré-inicializados  do  programa,  uma  página  
para  a  pilha  e  páginas  para  o  heap.  O  xv6  mapeia  os  dados,  a  pilha  e  o  heap  com  as  permissões  PTE_R,  PTE_W  e  
PTE_U.
 Usar  permissões  dentro  de  um  espaço  de  endereço  de  usuário  é  uma  técnica  comum  para  proteger  um  processo  de  usuário.
 Se  o  texto  fosse  mapeado  com  PTE_W,  um  processo  poderia  modificar  acidentalmente  seu  próprio  programa;  por  exemplo,  um  
erro  de  programação  pode  fazer  com  que  o  programa  escreva  em  um  ponteiro  nulo,  modificando  instruções  no  endereço  0,  e  
então  continue  em  execução,  talvez  causando  mais  confusão.  Para  detectar  esses  erros  imediatamente,  o  xv6  mapeia  o  texto  
sem  PTE_W;  se  um  programa  tentar  acidentalmente  armazenar  no  endereço  0,  o  hardware  se  recusará  a  executar  o  
armazenamento  e  gerará  uma  falha  de  página  (consulte  a  Seção  4.6).
 O  kernel  então  mata  o  processo  e  imprime  uma  mensagem  informativa  para  que  o  desenvolvedor  possa  rastrear  o  problema.
 Da  mesma  forma,  ao  mapear  dados  sem  PTE_X,  um  programa  de  usuário  não  pode  pular  acidentalmente  para  um
 endereço  nos  dados  do  programa  e  começar  a  executar  nesse  endereço.
 No  mundo  real,  fortalecer  um  processo  por  meio  da  definição  cuidadosa  de  permissões  também  auxilia  na  defesa  contra  
ataques  de  segurança.  Um  adversário  pode  fornecer  uma  entrada  cuidadosamente  construída  a  um  programa  (por  exemplo,  um  
servidor  web)  que  acione  um  bug  no  programa  na  esperança  de  transformá-lo  em  um  exploit  [14].
 Definir  permissões  cuidadosamente  e  outras  técnicas,  como  a  randomização  do  layout  do  espaço  de  endereço  do  usuário,  tornam  
esses  ataques  mais  difíceis.
 A  pilha  é  uma  única  página  e  é  exibida  com  o  conteúdo  inicial  criado  por  exec.  Strings  contendo  os  argumentos  da  linha  de  
comando,  bem  como  um  array  de  ponteiros  para  eles,  ficam  no  topo  da  pilha.  Logo  abaixo,  estão  os  valores  que  permitem  que  
um  programa  inicie  em  main  como  se  a  função  main(argc,  argv)  tivesse  acabado  de  ser  chamada.
 Para  detectar  que  uma  pilha  de  usuário  está  transbordando  a  memória  alocada,  o  xv6  coloca  uma  página  de  guarda  
inacessível  logo  abaixo  da  pilha,  limpando  o  sinalizador  PTE_U .  Se  a  pilha  de  usuário  transbordar  e  o  processo  tentar  usar  um  
endereço  abaixo  da  pilha,  o  hardware  gerará  uma  exceção  de  falha  de  página  porque  a  página  de  guarda  está  inacessível  para  
um  programa  em  execução  no  modo  de  usuário.  Um  sistema  operacional  real  pode,  em  vez  disso,  alocar  automaticamente  mais  
memória  para  a  pilha  de  usuário  quando  ela  transborda.
 Quando  um  processo  solicita  ao  xv6  mais  memória  do  usuário,  o  xv6  aumenta  o  heap  do  processo.  O  xv6  usa  primeiro
 38
Machine Translated by Google
 MAXVA
 TAMANHO  DA  PÁGINA
 Alinhamento  de  página
 0
 Receita--
estrutura  de  
armadilha  de  trampolim
 não  utilizado
 monte
 pilha
 pilha  
página  de  guarda
 dados
 não  utilizado
 texto
 RW-
R-WU
R-WU
R-WU
R-XU
 argumento  0
 ...
 argumento  N  0
 endereço  do  argumento  N
 ...
 endereço  do  argumento  0  
endereço  do  endereço  do  
argumento  0  
argc  
0xFFFFFFF
 (vazio)
 string  terminada  em  nulo  argv[argc]
 argv[0]
 argumento  argv  do  principal
 argumento  argc  do  retorno  
principal  PC  para  principal
 Figura  3.4:  Espaço  de  endereço  de  usuário  de  um  processo,  com  sua  pilha  inicial.
 kalloc  para  alocar  páginas  físicas.  Em  seguida,  ele  adiciona  PTEs  à  tabela  de  páginas  do  processo  que  apontam  para  
as  novas  páginas  físicas.  O  Xv6  define  os  sinalizadores  PTE_W,  PTE_R,  PTE_U  e  PTE_V  nesses  PTEs.  A  maioria  
dos  processos  não  utiliza  todo  o  espaço  de  endereço  do  usuário;  o  Xv6  deixa  PTE_V  livre  em  PTEs  não  utilizados.
 Vemos  aqui  alguns  bons  exemplos  do  uso  de  tabelas  de  páginas.  Primeiro,  as  tabelas  de  páginas  de  diferentes  
processos  traduzem  endereços  de  usuários  para  diferentes  páginas  de  memória  física,  de  modo  que  cada  processo  tenha  
memória  de  usuário  privada.  Segundo,  cada  processo  vê  sua  memória  como  tendo  endereços  virtuais  contíguos  começando  
em  zero,  enquanto  a  memória  física  do  processo  pode  ser  não  contígua.  Terceiro,  o  kernel  mapeia  uma  página  com  código  
trampolim  no  topo  do  espaço  de  endereço  do  usuário  (sem  PTE_U),  portanto,  uma  única  página  de  memória  física  aparece  
em  todos  os  espaços  de  endereço,  mas  pode  ser  usada  apenas  pelo  kernel.
 3.7  Código:  sbrk
 sbrk  é  a  chamada  de  sistema  para  que  um  processo  reduza  ou  aumente  sua  memória.  A  chamada  de  sistema  é  
implementada  pela  função  growproc  (kernel/proc.c:260).  growproc  chama  uvmalloc  ou  uvmdealloc,  dependendo  se  n  
é  positivo  ou  negativo.  uvmalloc  (kernel/vm.c:226)  aloca  memória  física  com  kalloc  e  adiciona  PTEs  à  tabela  de  
páginas  do  usuário  com  mappages.  uvmdealloc  chama  uvmunmap  (kernel/vm.c:171),  que  usa  walk  para  encontrar  
PTEs  e  kfree  para  liberar  a  memória  física  à  qual  se  referem.
 O  Xv6  usa  a  tabela  de  páginas  de  um  processo  não  apenas  para  informar  ao  hardware  como  mapear  os  endereços  
virtuais  do  usuário,  mas  também  como  o  único  registro  de  quais  páginas  de  memória  física  são  alocadas  para  aquele  
processo.  É  por  isso  que  a  liberação  de  memória  do  usuário  (no  uvmunmap)  requer  a  análise  da  tabela  de  páginas  do  usuário.
 39
Machine Translated by Google
 3.8  Código:  exec
 exec  é  uma  chamada  de  sistema  que  substitui  o  espaço  de  endereço  do  usuário  de  um  processo  por  dados  lidos  de  um  
arquivo,  chamado  de  arquivo  binário  ou  executável.  Um  binário  normalmente  é  a  saída  do  compilador  e  do  vinculador  e  
contém  instruções  de  máquina  e  dados  do  programa.  exec  (kernel/exec.c:23)  abre  o  caminho  binário  nomeado  usando  
namei  (kernel/exec.c:36),  que  é  explicado  no  Capítulo  8.  Em  seguida,  ele  lê  o  cabeçalho  ELF.  Os  binários  Xv6  são  
formatados  no  formato  ELF  amplamente  utilizado,  definido  em  (kernel/elf.h).  Um  binário  ELF  consiste  em  um  cabeçalho  
ELF,  struct  elfhdr  (kernel/elf.h:6),  seguido  por  uma  sequência  de  cabeçalhos  de  seção  do  programa,  struct  proghdr  (kernel/
 elf.h:25).  Cada  progvhdr  descreve  uma  seção  do  aplicativo  que  deve  ser  carregada  na  memória;  os  programas  xv6  têm  
dois  cabeçalhos  de  seção  de  programa:  um  para  instruções  e  um  para  dados.
 O  primeiro  passo  é  uma  verificação  rápida  para  verificar  se  o  arquivo  provavelmente  contém  um  binário  ELF.  Um  
binário  ELF  começa  com  o  "número  mágico"  de  quatro  bytes  0x7F,  'E',  'L',  'F'  ou  ELF_MAGIC  (kernel/elf.h:3).  Se  o  
cabeçalho  ELF  tiver  o  número  mágico  correto,  exec  assume  que  o  binário  está  bem  formado.
 exec  aloca  uma  nova  tabela  de  páginas  sem  mapeamentos  de  usuários  com  proc_pagetable  (kernel/exec.c:49),
 aloca  memória  para  cada  segmento  ELF  com  uvmalloc  (kernel/exec.c:65),  e  carrega  cada  segmento  na  memória  com  
loadseg  (kernel/exec.c:10).  loadseg  usa  walkaddr  para  encontrar  o  endereço  físico  da  memória  alocada  na  qual  gravar  
cada  página  do  segmento  ELF  e  readi  para  ler  o  arquivo.
 O  cabeçalho  da  seção  do  programa  para /init,  o  primeiro  programa  de  usuário  criado  com  exec,  se  parece  com  isto:
 #  objdump  -p  usuário/_init
 usuário/_init:
 Cabeçalho  do  programa:  
0x70000003  off
 CARREGAR  desligado
 formato  de  arquivo  elf64-little
 0x00000000000006bb0  vaddr  0x0000000000000000
 paddr  0x00000000000000000  align  2**0  filesz  0x000000000000004a  
memsz  0x0000000000000000  sinalizadores  r-
0x0000000000001000  vaddr  0x0000000000000000
 paddr  0x00000000000000000  alinhar  2**12  arquivosz  0x0000000000001000  
memsz  0x000000000001000  sinalizadores  rx
 CARREGAR  desligado
 0x0000000000002000  vaddr  0x000000000001000
 memsz  0x000000000000030  sinalizadores  rw
EMPILHAMENTO  desligado
 paddr  0x0000000000001000  alinhar  2**12  arquivosz  0x0000000000000010  
0x00000000000000000  vaddr  0x000000000000000  paddr  0x0000000000000000  
alinhar  2**4  arquivosz  0x0000000000000000  memsz  
0x0000000000000000  sinalizadores  rw
Vemos  que  o  segmento  de  texto  deve  ser  carregado  no  endereço  virtual  0x0  na  memória  (sem  permissões  de  
gravação)  a  partir  do  conteúdo  no  deslocamento  0x1000  do  arquivo.  Também  vemos  que  os  dados  devem  ser  carregados  
no  endereço  0x1000,  que  está  no  limite  de  uma  página,  e  sem  permissões  de  execução.
 O  valor  "filesz"  do  cabeçalho  de  uma  seção  de  programa  pode  ser  menor  que  o  valor  "memsz",  indicando  que  o  
espaço  entre  eles  deve  ser  preenchido  com  zeros  (para  variáveis  globais  em  C)  em  vez  de  lido  do  arquivo.  Para /init,  o  
valor  "filesz"  de  dados  é  0x10  bytes  e  o  valor  "memsz"  é  0x30  bytes,  e,  portanto,  uvmalloc  aloca  memória  física  suficiente  
para  armazenar  0x30  bytes,  mas  lê  apenas  0x10  bytes  do  arquivo /init.
 40
Machine Translated by Google
 Agora,  exec  aloca  e  inicializa  a  pilha  do  usuário.  Ele  aloca  apenas  uma  página  de  pilha.  exec  copia  as  strings  de  
argumentos  para  o  topo  da  pilha,  uma  de  cada  vez,  registrando  os  ponteiros  para  elas  em  ustack.
 Ele  coloca  um  ponteiro  nulo  no  final  do  que  será  a  lista  argv  passada  para  main.  As  três  primeiras  entradas  em  ustack  são  
o  contador  de  programa  de  retorno  falso,  argc  e  o  ponteiro  argv .
 O  comando  exec  coloca  uma  página  inacessível  logo  abaixo  da  página  da  pilha,  de  modo  que  programas  que  tentarem  
usar  mais  de  uma  página  apresentarão  falhas.  Essa  página  inacessível  também  permite  que  o  comando  exec  trate  de  
argumentos  muito  grandes;  nessa  situação,  a  cópia  (kernel/vm.c:352)  a  função  que  exec  usa  para  copiar  argumentos  para  
a  pilha  notará  que  a  página  de  destino  não  está  acessível  e  retornará  -1.
 Durante  a  preparação  da  nova  imagem  de  memória,  se  exec  detectar  um  erro,  como  um  segmento  de  programa  
inválido,  ele  salta  para  o  rótulo  "bad",  libera  a  nova  imagem  e  retorna  -1.  exec  deve  aguardar  para  liberar  a  imagem  antiga  
até  ter  certeza  de  que  a  chamada  de  sistema  será  bem-sucedida:  se  a  imagem  antiga  desaparecer,  a  chamada  de  sistema  
não  poderá  retornar  -1  para  ela.  Os  únicos  casos  de  erro  em  exec  ocorrem  durante  a  criação  da  imagem.  Assim  que  a  
imagem  estiver  concluída,  exec  pode  confirmar  a  nova  tabela  de  páginas  (kernel/exec.c:125).  e  liberar  o  antigo  (kernel/
 exec.c:129).
 exec  carrega  bytes  do  arquivo  ELF  para  a  memória  em  endereços  especificados  pelo  arquivo  ELF.  Usuários  ou  
processos  podem  colocar  quaisquer  endereços  que  desejarem  em  um  arquivo  ELF.  Portanto,  exec  é  arriscado,  pois  os  
endereços  no  arquivo  ELF  podem  se  referir  ao  kernel,  acidentalmente  ou  propositalmente.  As  consequências  para  um  
kernel  desavisado  podem  variar  de  uma  falha  a  uma  subversão  maliciosa  dos  mecanismos  de  isolamento  do  kernel  (ou  
seja,  uma  falha  de  segurança).  O  Xv6  realiza  uma  série  de  verificações  para  evitar  esses  riscos.  Por  exemplo,  if(ph.vaddr  
+  ph.memsz  <  ph.vaddr)  verifica  se  a  soma  excede  um  inteiro  de  64  bits.  O  perigo  é  que  um  usuário  possa  construir  um  
binário  ELF  com  um  ph.vaddr  que  aponte  para  um  endereço  escolhido  pelo  usuário  e  um  ph.memsz  grande  o  suficiente  
para  que  a  soma  exceda  para  0x1000,  o  que  parecerá  um  valor  válido.  Em  uma  versão  mais  antiga  do  xv6,  na  qual  o  
espaço  de  endereço  do  usuário  também  continha  o  kernel  (mas  não  era  legível/gravável  no  modo  usuário),  o  usuário  podia  
escolher  um  endereço  que  correspondesse  à  memória  do  kernel  e,  assim,  copiar  os  dados  do  binário  ELF  para  o  kernel.  
Na  versão  RISC-V  do  xv6,  isso  não  acontece,  pois  o  kernel  tem  sua  própria  tabela  de  páginas;  loadseg  carrega  na  tabela  
de  páginas  do  processo,  não  na  tabela  de  páginas  do  kernel.
 É  fácil  para  um  desenvolvedor  de  kernel  omitir  uma  verificação  crucial,  e  kernels  do  mundo  real  têm  um  longo  histórico  
de  verificações  ausentes,  cuja  ausência  pode  ser  explorada  por  programas  de  usuário  para  obter  privilégios  de  kernel.  É  
provável  que  o  xv6  não  faça  um  trabalho  completo  de  validação  dos  dados  de  nível  de  usuário  fornecidos  ao  kernel,  o  que  
um  programa  de  usuário  malicioso  pode  explorar  para  contornar  o  isolamento  do  xv6.
 3.9  Mundo  real
 Como  a  maioria  dos  sistemas  operacionais,  o  xv6  usa  o  hardware  de  paginação  para  proteção  e  mapeamento  de  memória.
 A  maioria  dos  sistemas  operacionais  faz  uso  muito  mais  sofisticado  de  paginação  do  que  o  xv6,  combinando  exceções  de  
paginação  e  falha  de  página,  que  discutiremos  no  Capítulo  4.
 O  Xv6  é  simplificado  pelo  uso,  pelo  kernel,  de  um  mapeamento  direto  entre  endereços  virtuais  e  físicos,  e  pela  
suposição  de  que  há  RAM  física  no  endereço  0x8000000,  onde  o  kernel  espera  ser  carregado.  Isso  funciona  com  o  QEMU,  
mas  em  hardware  real  acaba  sendo  uma  má  ideia;  o  hardware  real  coloca  RAM  e  dispositivos  em  endereços  físicos  
imprevisíveis,  de  modo  que  (por  exemplo)  pode  não  haver  RAM  em  0x8000000,  onde  o  xv6  espera  poder  armazenar  o  
kernel.  Um  kernel  mais  sério
 41
Machine Translated by Google
 os  projetos  exploram  a  tabela  de  páginas  para  transformar  layouts  arbitrários  de  memória  física  de  hardware  em  layouts  previsíveis  
de  endereços  virtuais  do  kernel.
 O  RISC-V  suporta  proteção  no  nível  de  endereços  físicos,  mas  o  xv6  não  usa  esse  recurso.
 Em  máquinas  com  muita  memória,  pode  fazer  sentido  usar  o  suporte  do  RISC-V  para  "superpáginas".  Páginas  pequenas  
fazem  sentido  quando  a  memória  física  é  pequena,  para  permitir  alocação  e  paginação  para  o  disco  com  granularidade  fina.  Por  
exemplo,  se  um  programa  usa  apenas  8  quilobytes  de  memória,  atribuir  a  ele  uma  superpágina  inteira  de  4  megabytes  de  memória  
física  é  um  desperdício.  Páginas  maiores  fazem  sentido  em  máquinas  com  muita  RAM  e  podem  reduzir  a  sobrecarga  na  manipulação  
da  tabela  de  páginas.
 A  falta  de  um  alocador  do  tipo  malloc  no  kernel  xv6  que  possa  fornecer  memória  para  objetos  pequenos  impede  que  o  kernel  
use  estruturas  de  dados  sofisticadas  que  exigiriam  alocação  dinâmica.
 Um  kernel  mais  elaborado  provavelmente  alocaria  muitos  tamanhos  diferentes  de  blocos  pequenos,  em  vez  de  (como  no  xv6)  
apenas  blocos  de  4.096  bytes;  um  alocador  de  kernel  real  precisaria  lidar  com  alocações  pequenas  e  também  grandes.
 A  alocação  de  memória  é  um  tópico  recorrente  e  polêmico,  cujos  problemas  básicos  são  o  uso  eficiente  de  memória  limitada  e  
a  preparação  para  solicitações  futuras  desconhecidas  [9].  Hoje  em  dia,  as  pessoas  se  preocupam  mais  com  velocidade  do  que  com  
eficiência  de  espaço.

Capítulo  4
 Armadilhas  e  chamadas  de  sistema
 Existem  três  tipos  de  eventos  que  fazem  com  que  a  CPU  interrompa  a  execução  normal  de  instruções  e  force  a  
transferência  de  controle  para  um  código  especial  que  lida  com  o  evento.  Uma  situação  é  uma  chamada  de  sistema,  
quando  um  programa  de  usuário  executa  a  instrução  ecall  para  solicitar  que  o  kernel  faça  algo  por  ele.  Outra  situação  é  
uma  exceção:  uma  instrução  (de  usuário  ou  kernel)  faz  algo  ilegal,  como  dividir  por  zero  ou  usar  um  endereço  virtual  
inválido.  A  terceira  situação  é  uma  interrupção  de  dispositivo,  quando  um  dispositivo  sinaliza  que  precisa  de  atenção,  por  
exemplo,  quando  o  hardware  do  disco  conclui  uma  solicitação  de  leitura  ou  gravação.
 Este  livro  usa  "trap"  como  um  termo  genérico  para  essas  situações.  Normalmente,  qualquer  código  que  estivesse  
em  execução  no  momento  da  trap  precisará  ser  retomado  posteriormente  e  não  precisaria  estar  ciente  de  que  algo  
especial  aconteceu.  Ou  seja,  frequentemente  queremos  que  as  traps  sejam  transparentes;  isso  é  particularmente  
importante  para  interrupções  de  dispositivo,  que  o  código  interrompido  normalmente  não  espera.  A  sequência  usual  é  
que  uma  trap  força  uma  transferência  de  controle  para  o  kernel;  o  kernel  salva  registradores  e  outros  estados  para  que  a  
execução  possa  ser  retomada;  o  kernel  executa  o  código  manipulador  apropriado  (por  exemplo,  uma  implementação  de  
chamada  de  sistema  ou  driver  de  dispositivo);  o  kernel  restaura  o  estado  salvo  e  retorna  da  trap;  e  o  código  original  
continua  de  onde  parou.
 O  Xv6  lida  com  todas  as  armadilhas  no  kernel;  as  armadilhas  não  são  entregues  ao  código  do  usuário.  Lidar  com  
armadilhas  no  kernel  é  natural  para  chamadas  de  sistema.  Faz  sentido  para  interrupções,  já  que  o  isolamento  exige  que  
apenas  o  kernel  tenha  permissão  para  usar  dispositivos  e  porque  o  kernel  é  um  mecanismo  conveniente  para  compartilhar  
dispositivos  entre  vários  processos.  Também  faz  sentido  para  exceções,  já  que  o  Xv6  responde  a  todas  as  exceções  do  
espaço  do  usuário  eliminando  o  programa  ofensivo.
 O  tratamento  de  traps  Xv6  ocorre  em  quatro  etapas:  ações  de  hardware  executadas  pela  CPU  RISC-V,  algumas  
instruções  de  montagem  que  preparam  o  caminho  para  o  código  C  do  kernel,  uma  função  C  que  decide  o  que  fazer  com  
o  trap  e  a  chamada  de  sistema  ou  rotina  de  serviço  do  driver  de  dispositivo.  Embora  a  semelhança  entre  os  três  tipos  de  
traps  sugira  que  um  kernel  poderia  lidar  com  todos  os  traps  com  um  único  caminho  de  código,  torna-se  conveniente  ter  
código  separado  para  três  casos  distintos:  traps  do  espaço  do  usuário,  traps  do  espaço  do  kernel  e  interrupções  de  
temporizador.  O  código  do  kernel  (assembler  ou  C)  que  processa  um  trap  é  frequentemente  chamado  de  manipulador;  
as  primeiras  instruções  do  manipulador  são  geralmente  escritas  em  assembler  (em  vez  de  C)  e,  às  vezes,  são  chamadas  
de  vetor.
 43
Machine Translated by Google
 4.1  Máquinas  de  armadilha  RISC-V
 Cada  CPU  RISC-V  possui  um  conjunto  de  registradores  de  controle  que  o  kernel  escreve  para  informar  
à  CPU  como  lidar  com  traps,  e  que  o  kernel  pode  ler  para  descobrir  se  um  trap  ocorreu.  Os  
documentos  RISC-V  contêm  a  história  completa  [3].  riscv.h  (kernel/riscv.h:1)  contém  definições  que  o  xv6  usa.
 Aqui  está  um  esboço  dos  registros  mais  importantes:
 •  stvec:  O  kernel  escreve  o  endereço  do  seu  manipulador  de  trap  aqui;  o  RISC-V  salta  para  o
 endereço  em  stvec  para  manipular  uma  armadilha.
 •
 sepc:  Quando  ocorre  uma  trap,  o  RISC-V  salva  o  contador  de  programa  aqui  (já  que  o  pc  é  então  sobrescrito  pelo  
valor  em  stvec).  A  instrução  sret  (retorno  da  trap)  copia  o  sepc  para  o  pc.  O  kernel  pode  escrever  sepc  para  controlar  
para  onde  o  sret  vai.
 •  causa:  o  RISC-V  coloca  um  número  aqui  que  descreve  o  motivo  da  armadilha.
 •  sscratch:  O  código  do  manipulador  de  traps  usa  sscratch  para  ajudar  a  evitar  a  substituição  de  registros  de  usuários
 antes  de  salvá-los.
 •  sstatus:  O  bit  SIE  em  sstatus  controla  se  as  interrupções  do  dispositivo  estão  habilitadas.  Se  o  kernel  limpar  o  SIE,  
o  RISC-V  adiará  as  interrupções  do  dispositivo  até  que  o  kernel  configure  o  SIE.  O  bit  SPP  indica  se  uma  
interceptação  veio  do  modo  usuário  ou  do  modo  supervisor  e  controla  para  qual  modo  sret  retorna.
 Os  registradores  acima  referem-se  a  traps  manipulados  no  modo  supervisor  e  não  podem  ser  lidos  ou  escritos  no  modo  
usuário.  Existe  um  conjunto  semelhante  de  registradores  de  controle  para  traps  manipulados  no  modo  máquina;  o  xv6  os  
utiliza  apenas  para  o  caso  especial  de  interrupções  de  temporizador.
 Cada  CPU  em  um  chip  multinúcleo  tem  seu  próprio  conjunto  desses  registradores,  e  mais  de  uma  CPU  pode  estar  
manipulando  uma  armadilha  a  qualquer  momento.
 Quando  precisa  forçar  uma  armadilha,  o  hardware  RISC-V  faz  o  seguinte  para  todos  os  tipos  de  armadilha  (exceto  
interrupções  de  temporizador):
 1.  Se  a  armadilha  for  uma  interrupção  de  dispositivo  e  o  bit  SIE  sstatus  estiver  limpo,  não  faça  nada
 seguindo.
 2.  Desabilite  as  interrupções  limpando  o  bit  SIE  em  sstatus.
 3.  Copie  o  pc  para  sepc.
 4.  Salve  o  modo  atual  (usuário  ou  supervisor)  no  bit  SPP  em  sstatus.
 5.  Defina  scause  para  refletir  a  causa  da  armadilha.
 6.  Defina  o  modo  como  supervisor.
 7.  Copie  stvec  para  o  pc.
 44
Machine Translated by Google
 8.  Comece  a  executar  no  novo  PC.
 Observe  que  a  CPU  não  alterna  para  a  tabela  de  páginas  do  kernel,  não  alterna  para  uma  pilha  no  kernel  e  não  salva  
nenhum  registrador  além  do  pc.  O  software  do  kernel  deve  executar  essas  tarefas.
 Um  motivo  pelo  qual  a  CPU  realiza  um  trabalho  mínimo  durante  uma  armadilha  é  para  fornecer  flexibilidade  ao  software;  por  
exemplo,  alguns  sistemas  operacionais  omitem  uma  troca  de  tabela  de  páginas  em  algumas  situações  para  aumentar  o  
desempenho  da  armadilha.
 Vale  a  pena  considerar  se  algum  dos  passos  listados  acima  poderia  ser  omitido,  talvez  em  busca  de  armadilhas  mais  
rápidas.  Embora  existam  situações  em  que  uma  sequência  mais  simples  possa  funcionar,  muitas  das  etapas  seriam  perigosas  
se  omitidas  em  geral.  Por  exemplo,  suponha  que  a  CPU  não  alterne  os  contadores  de  programa.  Então,  uma  armadilha  do  
espaço  do  usuário  poderia  alternar  para  o  modo  supervisor  enquanto  ainda  executa  instruções  do  usuário.  Essas  instruções  do  
usuário  poderiam  quebrar  o  isolamento  usuário/kernel,  por  exemplo,  modificando  o  registrador  satp  para  apontar  para  uma  
tabela  de  páginas  que  permitisse  o  acesso  a  toda  a  memória  física.
 Portanto,  é  importante  que  a  CPU  alterne  para  um  endereço  de  instrução  especificado  pelo  kernel,  ou  seja,  stvec.
 4.2  Armadilhas  do  espaço  do  usuário
 O  Xv6  lida  com  armadilhas  de  forma  diferente,  dependendo  se  elas  ocorrem  durante  a  execução  no  kernel  ou  no  código  
do  usuário.  Aqui  está  a  história  das  armadilhas  do  código  do  usuário;  a  Seção  4.5  descreve  as  armadilhas  do  código  do  
kernel.
 Uma  armadilha  pode  ocorrer  durante  a  execução  no  espaço  do  usuário  se  o  programa  do  usuário  fizer  uma  chamada  de  
sistema  (instrução  ecall),  ou  algo  ilegal,  ou  se  um  dispositivo  interromper.  O  caminho  de  alto  nível  de  uma  armadilha  do  espaço  
do  usuário  é  uservec  (kernel/trampoline.S:21).  então  usertrap  (kernel/trap.c:37);  e  ao  retornar,  usertrapret  (kernel/trap.c:90)  e  
então  userret  (kernel/trampoline.S:101).
 Uma  grande  limitação  no  design  do  tratamento  de  traps  do  xv6  é  o  fato  de  que  o  hardware  RISC-V  não  alterna  entre  
tabelas  de  páginas  ao  forçar  um  trap.  Isso  significa  que  o  endereço  do  manipulador  de  traps  em  stvec  deve  ter  um  mapeamento  
válido  na  tabela  de  páginas  do  usuário,  já  que  essa  é  a  tabela  de  páginas  em  vigor  quando  o  código  de  tratamento  de  traps  
começa  a  ser  executado.  Além  disso,  o  código  de  tratamento  de  traps  do  xv6  precisa  alternar  para  a  tabela  de  páginas  do  
kernel;  para  poder  continuar  a  execução  após  essa  troca,  a  tabela  de  páginas  do  kernel  também  deve  ter  um  mapeamento  para  
o  manipulador  apontado  por  stvec.
 O  Xv6  atende  a  esses  requisitos  usando  uma  página  de  trampolim.  A  página  de  trampolim  contém  uservec,
 O  código  de  manipulação  de  traps  xv6  para  o  qual  stvec  aponta.  A  página  do  trampolim  é  mapeada  na  tabela  de  páginas  de  
cada  processo  no  endereço  TRAMPOLINE,  que  está  no  topo  do  espaço  de  endereço  virtual,  de  modo  que  estará  acima  da  
memória  que  os  programas  usam  para  si  mesmos.  A  página  do  trampolim  também  é  mapeada  no  endereço  TRAMPOLINE  na  
tabela  de  páginas  do  kernel.  Veja  as  Figuras  2.3  e  3.3.  Como  a  página  do  trampolim  é  mapeada  na  tabela  de  páginas  do  
usuário,  sem  o  sinalizador  PTE_U,  as  traps  podem  começar  a  ser  executadas  lá  no  modo  supervisor.  Como  a  página  do  
trampolim  é  mapeada  no  mesmo  endereço  no  espaço  de  endereço  do  kernel,  o  manipulador  de  traps  pode  continuar  a  execução  
após  alternar  para  a  tabela  de  páginas  do  kernel.
 O  código  para  o  manipulador  de  armadilhas  uservec  está  em  trampoline.S  (kernel/trampoline.S:21).  Quando  o  uservec  
inicia,  todos  os  32  registradores  contêm  valores  pertencentes  ao  código  de  usuário  interrompido.  Esses  32  valores  precisam  
ser  salvos  em  algum  lugar  na  memória,  para  que  possam  ser  restaurados  quando  a  armadilha  retornar  ao  espaço  do  usuário.  
Armazenar  na  memória  requer  o  uso  de  um  registrador  para  armazenar  o  endereço,  mas  neste  ponto  não  há
 45
Machine Translated by Google
 Não  há  registradores  de  uso  geral  disponíveis!  Felizmente,  o  RISC-V  oferece  uma  ajuda  na  forma  do  registrador  
sscratch.  A  instrução  csrw  no  início  de  uservec  salva  a0  em  sscratch.
 Agora  o  uservec  tem  um  registrador  (a0)  para  brincar.
 A  próxima  tarefa  do  uservec  é  salvar  os  32  registradores  de  usuário.  O  kernel  aloca,  para  cada  processo,  uma  
página  de  memória  para  uma  estrutura  trapframe  que  (entre  outras  coisas)  tem  espaço  para  salvar  os  32  registradores  
de  usuário  (kernel/proc.h:43).  Como  satp  ainda  se  refere  à  tabela  de  páginas  do  usuário,  uservec  precisa  que  o  
trapframe  seja  mapeado  no  espaço  de  endereço  do  usuário.  Xv6  mapeia  o  trapframe  de  cada  processo  no  endereço  
virtual  TRAPFRAME  na  tabela  de  páginas  do  usuário  desse  processo;  TRAPFRAME  está  logo  abaixo  de  TRAMPOLINE.
 O  p->trapframe  do  processo  também  aponta  para  o  trapframe,  embora  em  seu  endereço  físico,  para  que  o  kernel  
possa  usá-lo  por  meio  da  tabela  de  páginas  do  kernel.
 Assim,  uservec  carrega  o  endereço  TRAPFRAME  em  a0  e  salva  todos  os  registros  do  usuário  lá,  incluindo  o  a0  
do  usuário,  lidos  do  zero.
 O  trapframe  contém  o  endereço  da  pilha  do  kernel  do  processo  atual,  o  hartid  da  CPU  atual,  o  endereço  da  função  
usertrap  e  o  endereço  da  tabela  de  páginas  do  kernel.  uservec  recupera  esses  valores,  alterna  satp  para  a  tabela  de  
páginas  do  kernel  e  chama  usertrap.
 O  trabalho  do  usertrap  é  determinar  a  causa  da  armadilha,  processá-la  e  retornar  (kernel/-  trap.c:37).  Primeiro,  ele  
altera  stvec  para  que  uma  trap  enquanto  estiver  no  kernel  seja  manipulada  por  kernelvec  em  vez  de  uservec.  Ele  salva  
o  registrador  sepc  (o  contador  de  programa  do  usuário  salvo),  porque  usertrap  pode  chamar  yield  para  alternar  para  a  
thread  do  kernel  de  outro  processo,  e  esse  processo  pode  retornar  ao  espaço  do  usuário,  no  processo  do  qual  ele  
modificará  sepc.  Se  a  trap  for  uma  chamada  de  sistema,  usertrap  chama  syscall  para  manipulá-la;  se  for  uma  
interrupção  de  dispositivo,  devintr;  caso  contrário,  é  uma  ex-ception,  e  o  kernel  mata  o  processo  com  falha.  O  caminho  
da  chamada  de  sistema  adiciona  quatro  ao  contador  de  programa  do  usuário  salvo  porque  o  RISC-V,  no  caso  de  uma  
chamada  de  sistema,  deixa  o  ponteiro  do  programa  apontando  para  a  instrução  ecall,  mas  o  código  do  usuário  precisa  
retomar  a  execução  na  instrução  subsequente.
 Na  saída,  o  usertrap  verifica  se  o  processo  foi  encerrado  ou  se  deve  render  a  CPU  (se  essa  armadilha  for  uma  
interrupção  de  timer).
 O  primeiro  passo  para  retornar  ao  espaço  do  usuário  é  a  chamada  para  usertrapret  (kernel/trap.c:90).  Esta  
função  configura  os  registradores  de  controle  RISC-V  para  se  prepararem  para  uma  futura  captura  a  partir  do  
espaço  do  usuário.  Isso  envolve  alterar  stvec  para  se  referir  a  uservec,  preparar  os  campos  do  trapframe  dos  
quais  uservec  depende  e  definir  sepc  para  o  contador  de  programa  do  usuário  salvo  anteriormente.  No  final,  
usertrapret  chama  userret  na  página  do  trampolim  mapeada  nas  tabelas  de  páginas  do  usuário  e  do  kernel;  o  
motivo  é  que  o  código  assembly  em  userret  alternará  as  tabelas  de  páginas.
 A  chamada  de  usertrapret  para  userret  passa  um  ponteiro  para  a  tabela  de  páginas  de  usuário  do  processo  em  
a0  (kernel/trampoline.S:101).  userret  alterna  o  satp  para  a  tabela  de  páginas  do  usuário  do  processo.  Lembre-se  de  
que  a  tabela  de  páginas  do  usuário  mapeia  tanto  a  página  do  trampolim  quanto  a  TRAPFRAME,  mas  nada  mais  do  kernel.
 O  mapeamento  da  página  do  trampolim  no  mesmo  endereço  virtual  nas  tabelas  de  páginas  do  usuário  e  do  kernel  
permite  que  o  userret  continue  executando  após  a  alteração  do  satp.  A  partir  deste  ponto,  os  únicos  dados  que  o  
userret  pode  usar  são  o  conteúdo  do  registrador  e  o  conteúdo  do  trapframe.  O  userret  carrega  o  endereço  TRAPFRAME  
em  a0,  restaura  os  registradores  de  usuário  salvos  do  trapframe  via  a0,  restaura  o  usuário  a0  salvo  e  executa  sret  
para  retornar  ao  espaço  do  usuário.
 46
Machine Translated by Google
 4.3  Código:  Chamando  chamadas  de  sistema
 O  capítulo  2  terminou  com  initcode.S  invocando  a  chamada  de  sistema  exec  (user/initcode.S:11).  Vamos  ver  como  a  chamada  
do  usuário  chega  à  implementação  da  chamada  do  sistema  exec  no  kernel.
 initcode.S  coloca  os  argumentos  para  exec  nos  registradores  a0  e  a1,  e  coloca  o  número  da  chamada  do  sistema  em  a7.  
Os  números  da  chamada  do  sistema  correspondem  às  entradas  no  array  syscalls,  uma  tabela  de  ponteiros  de  função  (kernel/
 syscall.c:107).  A  instrução  ecall  é  capturada  no  kernel  e  faz  com  que  uservec,  usertrap  e  então  syscall  sejam  executados,  
como  vimos  acima.
 chamada  de  sistema  (kernel/syscall.c:132)  recupera  o  número  de  chamada  do  sistema  do  a7  salvo  no  trapframe
 e  o  utiliza  para  indexar  chamadas  de  sistema.  Para  a  primeira  chamada  de  sistema,  a7  contém  SYS_exec  (kernel /syscall.h:8),  
resultando  em  uma  chamada  para  a  função  de  implementação  de  chamada  do  sistema  sys_exec.
 Quando  sys_exec  retorna,  syscall  registra  seu  valor  de  retorno  em  p->trapframe->a0.  Isso  fará  com  que  a  chamada  
original  do  espaço  do  usuário  para  exec()  retorne  esse  valor,  já  que  a  convenção  de  chamada  C  no  RISC-V  coloca  os  valores  
de  retorno  em  a0.  Chamadas  de  sistema  convencionalmente  retornam  números  negativos  para  indicar  erros  e  zero  ou  números  
positivos  para  sucesso.  Se  o  número  da  chamada  de  sistema  for  inválido,  syscall  imprime  um  erro  e  retorna  -1.
 4.4  Código:  Argumentos  de  chamada  de  sistema
 Implementações  de  chamadas  de  sistema  no  kernel  precisam  encontrar  os  argumentos  passados  pelo  código  do  usuário.  
Como  o  código  do  usuário  chama  funções  wrapper  de  chamadas  de  sistema,  os  argumentos  estão  inicialmente  onde  a  
convenção  de  chamada  RISC-VC  os  coloca:  em  registradores.  O  código  de  trap  do  kernel  salva  os  registradores  do  
usuário  no  quadro  de  trap  do  processo  atual,  onde  o  código  do  kernel  pode  encontrá-los.  As  funções  do  kernel  argint,  
argaddr  e  argfd  recuperam  o  n  'ésimo  argumento  de  chamada  de  sistema  do  quadro  de  trap  como  um  inteiro,  um  ponteiro  
ou  um  descritor  de  arquivo.  Todas  elas  chamam  argraw  para  recuperar  o  registrador  de  usuário  salvo  apropriado  (kernel/syscall.c:34).
 Algumas  chamadas  de  sistema  passam  ponteiros  como  argumentos,  e  o  kernel  deve  usar  esses  ponteiros  para  ler  ou  
escrever  na  memória  do  usuário.  A  chamada  de  sistema  exec,  por  exemplo,  passa  ao  kernel  um  conjunto  de  ponteiros  
referentes  a  argumentos  de  string  no  espaço  do  usuário.  Esses  ponteiros  apresentam  dois  desafios.  Primeiro,  o  programa  do  
usuário  pode  estar  com  bugs  ou  ser  malicioso,  e  pode  passar  ao  kernel  um  ponteiro  inválido  ou  um  ponteiro  com  a  intenção  
de  induzir  o  kernel  a  acessar  a  memória  do  kernel  em  vez  da  memória  do  usuário.  Segundo,  os  mapeamentos  de  tabelas  de  
páginas  do  kernel  xv6  não  são  os  mesmos  que  os  mapeamentos  de  tabelas  de  páginas  do  usuário,  portanto,  o  kernel  não  
pode  usar  instruções  comuns  para  carregar  ou  armazenar  a  partir  de  endereços  fornecidos  pelo  usuário.
 O  kernel  implementa  funções  que  transferem  dados  com  segurança  de  e  para  endereços  fornecidos  pelo  usuário.  fetchstr  
é  um  exemplo  (kernel/syscall.c:25).  Chamadas  de  sistema  de  arquivos  como  exec  usam  fetchstr  para  recuperar  argumentos  
de  nome  de  arquivo  de  string  do  espaço  do  usuário.  fetchstr  chama  copyinstr  para  fazer  o  trabalho  pesado.  copyinstr  (kernel/
 vm.c:403)  Copia  até  max  bytes  para  dst  do  endereço  virtual  srcva  na  tabela  de  páginas  do  usuário  pagetable.  Como  pagetable  
não  é  a  tabela  de  páginas  atual,  copyinstr  usa  walkaddr  (que  chama  walk)  para  procurar  srcva  na  pagetable,  retornando  o  
endereço  físico  pa0.  O  kernel  mapeia  cada  endereço  físico  de  RAM  para  o  endereço  virtual  do  kernel  correspondente,  de  
modo  que  copyinstr  pode  copiar  diretamente  bytes  de  string  de  pa0  para  dst.  walkaddr  (kernel/vm.c:109)  verifica  se  o  endereço  
virtual  fornecido  pelo  usuário  faz  parte  do  espaço  de  endereço  do  usuário  do  processo,  para  que  os  programas
 47
Machine Translated by Google
 não  consegue  enganar  o  kernel  para  que  ele  leia  outra  memória.  Uma  função  semelhante,  copyout,  copia  dados  do  kernel  para  
um  endereço  fornecido  pelo  usuário.
 4.5  Armadilhas  do  espaço  do  kernel
 O  Xv6  configura  os  registradores  de  trap  da  CPU  de  forma  um  pouco  diferente,  dependendo  se  o  código  do  usuário  ou  do  kernel  
está  em  execução.  Quando  o  kernel  está  em  execução  em  uma  CPU,  o  kernel  aponta  stvec  para  o  código  assembly  em  kernelvec  
(kernel/kernelvec.S:12).  Como  o  xv6  já  está  no  kernel,  o  kernelvec  pode  contar  com  satp  sendo  definido  na  tabela  de  páginas  do  
kernel  e  com  o  ponteiro  da  pilha  se  referindo  a  uma  pilha  do  kernel  válida.  O  kernelvec  envia  todos  os  32  registradores  para  a  
pilha,  de  onde  ele  os  restaurará  posteriormente  para  que  o  código  do  kernel  interrompido  possa  ser  retomado  sem  perturbações.
 kernelvec  salva  os  registradores  na  pilha  da  thread  do  kernel  interrompida,  o  que  faz  sentido  porque  os  valores  dos  
registradores  pertencem  a  essa  thread.  Isso  é  particularmente  importante  se  a  armadilha  causar  uma  troca  para  uma  thread  
diferente  –  nesse  caso,  a  armadilha  retornará  da  pilha  da  nova  thread,  deixando  os  registradores  salvos  da  thread  interrompida  
em  segurança  em  sua  pilha.
 kernelvec  salta  para  kerneltrap  (kernel/trap.c:135)  Após  salvar  os  registradores,  o  kerneltrap  está  
preparado  para  dois  tipos  de  armadilhas:  interrupções  de  dispositivo  e  exceções.  Ele  chama  devintr  (kernel/  -trap.c:178)  para  verificar  e  lidar  com  o  primeiro.  Se  a  armadilha  não  for  uma  interrupção  de  dispositivo,  deve  
ser  uma  exceção,  e  isso  é  sempre  um  erro  fatal  se  ocorrer  no  kernel  xv6;  o  kernel  dispara  pânico  e  interrompe  
a  execução.
 Se  o  kerneltrap  foi  chamado  devido  a  uma  interrupção  do  temporizador  e  uma  thread  do  kernel  de  um  processo  está  em  
execução  (em  oposição  a  uma  thread  do  escalonador),  o  kerneltrap  chama  o  método  yield  para  dar  a  outras  threads  a  chance  
de  executar.  Em  algum  momento,  uma  dessas  threads  irá  executar  o  método  yield,  permitindo  que  nossa  thread  e  sua  kerneltrap  
sejam  retomadas.  O  Capítulo  7  explica  o  que  acontece  no  método  yield.
 Quando  o  trabalho  do  kerneltrap  termina,  ele  precisa  retornar  ao  código  que  foi  interrompido  pela  armadilha.  Como  um  yield  
pode  ter  perturbado  o  sepc  e  o  modo  anterior  em  sstatus,  o  kerneltrap  os  salva  ao  iniciar.  Agora,  ele  restaura  esses  registradores  
de  controle  e  retorna  ao  kernelvec  (kernel/kernelvec.S:50).  kernelvec  retira  os  registradores  salvos  da  pilha  e  executa  sret,  que  
copia  sepc  para  pc  e  retoma  o  código  do  kernel  interrompido.
 Vale  a  pena  pensar  em  como  o  retorno  do  trap  acontece  se  o  kerneltrap  chamar  yield  devido  a  uma  interrupção  do  
temporizador.
 O  Xv6  define  o  stvec  de  uma  CPU  como  kernelvec  quando  essa  CPU  entra  no  kernel  pelo  espaço  do  
usuário;  você  pode  ver  isso  em  usertrap  (kernel/trap.c:29).  Há  uma  janela  de  tempo  em  que  o  kernel  inicia  a  
execução,  mas  stvec  ainda  está  definido  como  uservec,  e  é  crucial  que  nenhuma  interrupção  de  dispositivo  
ocorra  durante  essa  janela.  Felizmente,  o  RISC-V  sempre  desabilita  as  interrupções  quando  começa  a  receber  
uma  captura,  e  o  xv6  não  as  habilita  novamente  até  que  stvec  seja  definido.
 4.6  Exceções  de  falha  de  página
 A  resposta  do  Xv6  a  exceções  é  bastante  tediosa:  se  uma  exceção  acontece  no  espaço  do  usuário,  o  kernel  encerra  o  processo  
com  falha.  Se  uma  exceção  acontece  no  kernel,  o  kernel  entra  em  pânico.  Operação  real
 48
Machine Translated by Google
 os  sistemas  geralmente  respondem  de  maneiras  muito  mais  interessantes.
 Como  exemplo,  muitos  kernels  usam  falhas  de  página  para  implementar  bifurcações  de  cópia  na  escrita  (COW).  Para  explicar  a  
bifurcação  de  cópia  na  escrita,  considere  a  bifurcação  do  xv6,  descrita  no  Capítulo  3.  A  bifurcação  faz  com  que  o  conteúdo  de  memória  
inicial  do  filho  seja  o  mesmo  do  pai  no  momento  da  bifurcação.  O  Xv6  implementa  a  bifurcação  com  uvmcopy  (kernel/vm.c:306).  que  
aloca  memória  física  para  o  filho  e  copia  a  memória  do  pai  para  ela.  Seria  mais  eficiente  se  o  filho  e  o  pai  pudessem  compartilhar  a  
memória  física  do  pai.  Uma  implementação  direta  disso  não  funcionaria,  no  entanto,  pois  faria  com  que  o  pai  e  o  filho  interrompessem  
a  execução  um  do  outro  com  suas  gravações  na  pilha  e  no  heap  compartilhados.
 Pai  e  filho  podem  compartilhar  memória  física  com  segurança  por  meio  do  uso  apropriado  de  permissões  de  tabela  
de  páginas  e  falhas  de  página.  A  CPU  gera  uma  exceção  de  falha  de  página  quando  um  endereço  virtual  é  usado  sem  
mapeamento  na  tabela  de  páginas,  ou  tem  um  mapeamento  cujo  sinalizador  PTE_V  está  limpo,  ou  um  mapeamento  cujos  
bits  de  permissão  (PTE_R,  PTE_W,  PTE_X,  PTE_U)  proíbem  a  operação  que  está  sendo  tentada.  O  RISC-V  distingue  
três  tipos  de  falha  de  página:  falhas  de  página  de  carregamento  (quando  uma  instrução  de  carregamento  não  consegue  
traduzir  seu  endereço  virtual),  falhas  de  página  de  armazenamento  (quando  uma  instrução  de  armazenamento  não  
consegue  traduzir  seu  endereço  virtual)  e  falhas  de  página  de  instrução  (quando  o  endereço  no  contador  de  programa  
não  traduz).  O  registrador  scause  indica  o  tipo  de  falha  de  página  e  o  registrador  stval  contém  o  endereço  que  não  pôde  
ser  traduzido.
 O  plano  básico  em  uma  bifurcação  COW  é  que  o  pai  e  o  filho  compartilhem  inicialmente  todas  as  páginas  físicas,  mas  cada  um  
mapeie-as  como  somente  leitura  (com  o  sinalizador  PTE_W  limpo).  Pai  e  filho  podem  ler  a  partir  da  memória  física  compartilhada.  Se  
qualquer  um  deles  gravar  uma  determinada  página,  a  CPU  RISC-V  gera  uma  exceção  de  falha  de  página.  O  manipulador  de  traps  do  
kernel  responde  alocando  uma  nova  página  de  memória  física  e  copiando  para  ela  a  página  física  para  a  qual  o  endereço  com  falha  
mapeia.  O  kernel  altera  a  PTE  relevante  na  tabela  de  páginas  do  processo  com  falha  para  apontar  para  a  cópia  e  permitir  gravações,  
bem  como  leituras,  e  então  retoma  o  processo  com  falha  na  instrução  que  causou  a  falha.  Como  a  PTE  permite  gravações,  a  instrução  
reexecutada  agora  será  executada  sem  falhas.  A  cópia  na  gravação  requer  contabilidade  para  ajudar  a  decidir  quando  as  páginas  
físicas  podem  ser  liberadas,  uma  vez  que  cada  página  pode  ser  referenciada  por  um  número  variável  de  tabelas  de  páginas,  
dependendo  do  histórico  de  bifurcações,  falhas  de  página,  execuções  e  saídas.
 Essa  contabilidade  permite  uma  otimização  importante:  se  um  processo  incorrer  em  uma  falha  de  página  de  armazenamento  e  a  
página  física  for  referenciada  apenas  pela  tabela  de  páginas  desse  processo,  nenhuma  cópia  será  necessária.
 A  cópia  na  gravação  torna  a  bifurcação  mais  rápida,  já  que  a  bifurcação  não  precisa  copiar  memória.  Parte  da  memória  precisará  
ser  copiada  posteriormente,  quando  escrita,  mas  frequentemente  a  maior  parte  da  memória  nunca  precisa  ser  copiada.  Um  exemplo  
comum  é  a  bifurcação  seguida  por  um  comando  exec:  algumas  páginas  podem  ser  escritas  após  a  bifurcação,  mas  o  comando  exec  
da  filha  libera  a  maior  parte  da  memória  herdada  da  mãe.
 A  bifurcação  copy-on-write  elimina  a  necessidade  de  copiar  essa  memória.  Além  disso,  a  bifurcação  COW  é  transparente:  nenhuma  
modificação  nos  aplicativos  é  necessária  para  que  eles  se  beneficiem.
 A  combinação  de  tabelas  de  páginas  e  falhas  de  página  abre  uma  ampla  gama  de  possibilidades  interessantes,  além  da  
bifurcação  COW.  Outro  recurso  amplamente  utilizado  é  a  alocação  preguiçosa,  que  possui  duas  partes.  Primeiro,  quando  um  aplicativo  
solicita  mais  memória  chamando  sbrk,  o  kernel  percebe  o  aumento  de  tamanho,  mas  não  aloca  memória  física  e  não  cria  PTEs  para  o  
novo  intervalo  de  endereços  virtuais.  Segundo,  em  uma  falha  de  página  em  um  desses  novos  endereços,  o  kernel  aloca  uma  página  
de  memória  física  e  a  mapeia  na  tabela  de  páginas.  Assim  como  a  bifurcação  COW,  o  kernel  pode  implementar
 49
Machine Translated by Google
 alocação  preguiçosa  de  forma  transparente  para  aplicativos.
 Como  os  aplicativos  frequentemente  solicitam  mais  memória  do  que  precisam,  a  alocação  preguiçosa  é  vantajosa:  o  kernel  não  
precisa  realizar  nenhum  trabalho  para  páginas  que  o  aplicativo  nunca  usa.  Além  disso,  se  o  aplicativo  solicitar  um  aumento  significativo  
no  espaço  de  endereço,  o  sbrk  sem  alocação  preguiçosa  é  caro:  se  um  aplicativo  solicitar  um  gigabyte  de  memória,  o  kernel  precisa  
alocar  e  zerar  262.144  páginas  de  4.096  bytes.  A  alocação  preguiçosa  permite  que  esse  custo  seja  distribuído  ao  longo  do  tempo.  
Por  outro  lado,  a  alocação  preguiçosa  incorre  na  sobrecarga  extra  de  falhas  de  página,  que  envolvem  uma  transição  kernel/usuário.  
Os  sistemas  operacionais  podem  reduzir  esse  custo  alocando  um  lote  de  páginas  consecutivas  por  falha  de  página  em  vez  de  uma  
página  e  especializando  o  código  de  entrada/saída  do  kernel  para  tais  falhas  de  página.
 Outro  recurso  amplamente  utilizado  que  explora  falhas  de  página  é  a  paginação  sob  demanda.  Em  exec,  o  xv6  carrega  todo  o  
texto  e  os  dados  de  um  aplicativo  avidamente  na  memória.  Como  os  aplicativos  podem  ser  grandes  e  a  leitura  do  disco  é  cara,  esse  
custo  inicial  pode  ser  perceptível  para  os  usuários:  quando  o  usuário  inicia  um  aplicativo  grande  a  partir  do  shell,  pode  levar  muito  
tempo  até  que  ele  veja  uma  resposta.  Para  melhorar  o  tempo  de  resposta,  um  kernel  moderno  cria  a  tabela  de  páginas  para  o  espaço  
de  endereço  do  usuário,  mas  marca  as  PTEs  das  páginas  como  inválidas.  Em  uma  falha  de  página,  o  kernel  lê  o  conteúdo  da  página  
do  disco  e  o  mapeia  para  o  espaço  de  endereço  do  usuário.  Assim  como  o  COW  fork  e  a  alocação  preguiçosa,  o  kernel  pode  
implementar  esse  recurso  de  forma  transparente  para  os  aplicativos.
 Os  programas  em  execução  em  um  computador  podem  precisar  de  mais  memória  do  que  a  RAM  disponível  no  computador.
 Para  lidar  com  a  situação  de  forma  elegante,  o  sistema  operacional  pode  implementar  a  paginação  em  disco.  A  ideia  é  armazenar  
apenas  uma  fração  das  páginas  do  usuário  na  RAM  e  armazenar  o  restante  em  disco,  em  uma  área  de  paginação.  O  kernel  marca  
as  PTEs  que  correspondem  à  memória  armazenada  na  área  de  paginação  (e,  portanto,  não  na  RAM)  como  inválidas.  Se  um  aplicativo  
tentar  usar  uma  das  páginas  que  foram  paginadas  para  o  disco,  o  aplicativo  incorrerá  em  uma  falha  de  página  e  a  página  deverá  ser  
paginada:  o  manipulador  de  traps  do  kernel  alocará  uma  página  de  RAM  física,  lerá  a  página  do  disco  para  a  RAM  e  modificará  a  
PTE  relevante  para  apontar  para  a  RAM.
 O  que  acontece  se  uma  página  precisar  ser  paginada,  mas  não  houver  RAM  física  livre?  Nesse  caso,  o  kernel  deve  primeiro  
liberar  uma  página  física,  paginando-a  para  fora  ou  removendo-a  para  a  área  de  paginação  no  disco  e  marcando  as  PTEs  que  se  
referem  a  essa  página  física  como  inválidas.  A  remoção  é  dispendiosa,  portanto,  a  paginação  tem  melhor  desempenho  se  for  pouco  
frequente:  se  os  aplicativos  usarem  apenas  um  subconjunto  de  suas  páginas  de  memória  e  a  união  dos  subconjuntos  couber  na  
RAM.  Essa  propriedade  é  frequentemente  chamada  de  boa  localidade  de  referência.  Assim  como  acontece  com  muitas  técnicas  de  
memória  virtual,  os  kernels  geralmente  implementam  a  paginação  para  o  disco  de  forma  transparente  para  os  aplicativos.
 Os  computadores  frequentemente  operam  com  pouca  ou  nenhuma  memória  física  livre,  independentemente  da  quantidade  de  
RAM  fornecida  pelo  hardware.  Por  exemplo,  provedores  de  nuvem  multiplexam  muitos  clientes  em  uma  única  máquina  para  usar  
seu  hardware  de  forma  econômica.  Outro  exemplo:  usuários  executam  muitos  aplicativos  em  smartphones  com  uma  pequena  
quantidade  de  memória  física.  Em  tais  cenários,  a  alocação  de  uma  página  pode  exigir  a  remoção  prévia  de  uma  página  existente.  
Portanto,  quando  a  memória  física  livre  é  escassa,  a  alocação  é  dispendiosa.
 A  alocação  lenta  e  a  paginação  sob  demanda  são  particularmente  vantajosas  quando  a  memória  livre  é  escassa.
 A  alocação  avida  de  memória  em  sbrk  ou  exec  incorre  no  custo  extra  de  remoção  para  disponibilizar  memória.  Além  disso,  existe  o  
risco  de  que  o  trabalho  avidamente  realizado  seja  desperdiçado,  pois  antes  que  o  aplicativo  use  a  página,  o  sistema  operacional  pode  
tê-la  removido.
 Outros  recursos  que  combinam  exceções  de  paginação  e  falha  de  página  incluem  extensão  automática
 pilhas  e  arquivos  mapeados  na  memória.
 50
Machine Translated by Google
 4.7  Mundo  real
 O  trampolim  e  o  trapframe  podem  parecer  excessivamente  complexos.  Uma  força  motriz  é  que  o  RISC-V  intencionalmente  
faz  o  mínimo  possível  ao  forçar  uma  armadilha,  para  permitir  a  possibilidade  de  um  tratamento  muito  rápido  da  armadilha,  
o  que  acaba  sendo  importante.  Como  resultado,  as  primeiras  instruções  do  manipulador  de  armadilhas  do  kernel  precisam  
ser  executadas  efetivamente  no  ambiente  do  usuário:  a  tabela  de  páginas  do  usuário  e  o  conteúdo  do  registrador  do  
usuário.  E  o  manipulador  de  armadilhas  inicialmente  ignora  fatos  úteis,  como  a  identidade  do  processo  em  execução  ou  o  
endereço  da  tabela  de  páginas  do  kernel.  Uma  solução  é  possível  porque  o  RISC-V  fornece  locais  protegidos  nos  quais  o  
kernel  pode  armazenar  informações  antes  de  entrar  no  espaço  do  usuário:  o  registrador  sscratch  e  as  entradas  da  tabela  
de  páginas  do  usuário  que  apontam  para  a  memória  do  kernel,  mas  são  protegidas  pela  ausência  de  PTE_U.  O  trampolim  
e  o  trapframe  do  Xv6  exploram  esses  recursos  do  RISC-V.
 A  necessidade  de  páginas  especiais  de  trampolim  poderia  ser  eliminada  se  a  memória  do  kernel  fosse  mapeada  na  
tabela  de  páginas  de  usuário  de  cada  processo  (com  os  sinalizadores  de  permissão  PTE  apropriados).  Isso  também  
eliminaria  a  necessidade  de  uma  troca  de  tabela  de  páginas  ao  capturar  do  espaço  do  usuário  para  o  kernel.  Isso,  por  sua  
vez,  permitiria  que  implementações  de  chamadas  de  sistema  no  kernel  aproveitassem  o  mapeamento  da  memória  de  
usuário  do  processo  atual,  permitindo  que  o  código  do  kernel  desreferenciasse  diretamente  os  ponteiros  do  usuário.
 Muitos  sistemas  operacionais  têm  usado  essas  ideias  para  aumentar  a  eficiência.  O  Xv6  as  evita  para  reduzir  as  chances  
de  bugs  de  segurança  no  kernel  devido  ao  uso  inadvertido  de  ponteiros  de  usuário  e  para  reduzir  um  pouco  da  
complexidade  necessária  para  garantir  que  os  endereços  virtuais  do  usuário  e  do  kernel  não  se  sobreponham.
 Sistemas  operacionais  de  produção  implementam  bifurcação  de  cópia  na  gravação,  alocação  preguiçosa,  paginação  
sob  demanda,  paginação  em  disco,  arquivos  mapeados  na  memória,  etc.  Além  disso,  sistemas  operacionais  de  produção  
tentarão  usar  toda  a  memória  física,  seja  para  aplicativos  ou  caches  (por  exemplo,  o  cache  de  buffer  do  sistema  de  
arquivos,  que  abordaremos  mais  adiante  na  Seção  8.2).  O  Xv6  é  ingênuo  nesse  aspecto:  você  quer  que  seu  sistema  
operacional  use  a  memória  física  pela  qual  você  pagou,  mas  o  xv6  não.  Além  disso,  se  o  xv6  ficar  sem  memória,  ele  
retorna  um  erro  para  o  aplicativo  em  execução  ou  o  encerra,  em  vez  de,  por  exemplo,  remover  uma  página  de  outro  
aplicativo

 Capítulo  7
 Agendamento
 Qualquer  sistema  operacional  provavelmente  executará  com  mais  processos  do  que  a  capacidade  de  CPU  do  
computador,  portanto,  é  necessário  um  plano  para  compartilhar  o  tempo  das  CPUs  entre  os  processos.  Idealmente,  o  
compartilhamento  seria  transparente  para  os  processos  do  usuário.  Uma  abordagem  comum  é  dar  a  cada  processo  a  
ilusão  de  que  possui  sua  própria  CPU  virtual,  multiplexando  os  processos  nas  CPUs  de  hardware.  Este  capítulo  explica  
como  o  xv6  realiza  essa  multiplexação.
 7.1  Multiplexação
 O  Xv6  multiplexa  alternando  cada  CPU  de  um  processo  para  outro  em  duas  situações.  Primeiro,  o  mecanismo  de  
suspensão  e  ativação  do  xv6  alterna  quando  um  processo  aguarda  a  conclusão  da  E/S  do  dispositivo  ou  do  pipe,  ou  
aguarda  a  saída  de  um  processo  filho,  ou  aguarda  na  chamada  de  sistema  de  suspensão .  Segundo,  o  xv6  força  
periodicamente  uma  troca  para  lidar  com  processos  que  computam  por  longos  períodos  sem  suspensão.  Essa  
multiplexação  cria  a  ilusão  de  que  cada  processo  possui  sua  própria  CPU,  assim  como  o  xv6  utiliza  o  alocador  de  
memória  e  as  tabelas  de  páginas  de  hardware  para  criar  a  ilusão  de  que  cada  processo  possui  sua  própria  memória.
 A  implementação  da  multiplexação  apresenta  alguns  desafios.  Primeiro,  como  alternar  de  um  processo  para  
outro?  Embora  a  ideia  de  troca  de  contexto  seja  simples,  a  implementação  é  um  dos  códigos  mais  opacos  do  Xv6.  
Segundo,  como  forçar  trocas  de  forma  transparente  para  os  processos  do  usuário?  O  Xv6  usa  a  técnica  padrão  na  
qual  as  interrupções  de  um  temporizador  de  hardware  controlam  as  trocas  de  contexto.  Terceiro,  todas  as  CPUs  
alternam  entre  o  mesmo  conjunto  compartilhado  de  processos,  e  um  plano  de  bloqueio  é  necessário  para  evitar  
corridas.  Quarto,  a  memória  e  outros  recursos  de  um  processo  devem  ser  liberados  quando  o  processo  termina,  mas  
ele  não  pode  fazer  tudo  isso  sozinho  porque  (por  exemplo)  não  pode  liberar  sua  própria  pilha  do  kernel  enquanto  ainda  
a  utiliza.  Quinto,  cada  núcleo  de  uma  máquina  multinúcleo  deve  se  lembrar  de  qual  processo  está  executando  para  
que  as  chamadas  de  sistema  afetem  o  estado  correto  do  kernel  do  processo.  Finalmente,  a  suspensão  e  a  ativação  
permitem  que  um  processo  desista  da  CPU  e  espere  ser  ativado  por  outro  processo  ou  interrupção.  É  necessário  
cuidado  para  evitar  corridas  que  resultem  na  perda  de  notificações  de  ativação.  O  Xv6  tenta  resolver  esses  problemas  
da  forma  mais  simples  possível,  mas  mesmo  assim  o  código  resultante  é  complicado.
 71
Machine Translated by Google
 usuário
 espaço
 núcleo
 espaço
 concha
 salvar
 kstack
 concha
 gato
 interruptor
 interruptor
 agendador  
kstack
 Núcleo
 kstack
 gato
 restaurar
 Figura  7.1:  Troca  de  um  processo  de  usuário  para  outro.  Neste  exemplo,  o  xv6  é  executado  com  uma  CPU  (e,  portanto,  
uma  thread  do  escalonador).
 7.2  Código:  Troca  de  contexto
 A  Figura  7.1  descreve  as  etapas  envolvidas  na  alternância  de  um  processo  de  usuário  para  outro:  uma  transição  usuário
kernel  (chamada  de  sistema  ou  interrupção)  para  a  thread  do  kernel  do  processo  antigo,  uma  troca  de  contexto  para  a  
thread  do  escalonador  da  CPU  atual,  uma  troca  de  contexto  para  a  thread  do  kernel  de  um  novo  processo  e  um  retorno  
de  trap  para  o  processo  em  nível  de  usuário.  O  escalonador  xv6  possui  uma  thread  dedicada  (registradores  e  pilha  salvos)  
por  CPU,  pois  não  é  seguro  para  o  escalonador  executar  na  pilha  do  kernel  do  processo  antigo:  algum  outro  núcleo  pode  
despertar  o  processo  e  executá-lo,  e  seria  um  desastre  usar  a  mesma  pilha  em  dois  núcleos  diferentes.  Nesta  seção,  
examinaremos  a  mecânica  da  alternância  entre  uma  thread  do  kernel  e  uma  thread  do  escalonador.
 Alternar  de  um  thread  para  outro  envolve  salvar  os  registradores  da  CPU  do  thread  antigo  e  restaurar  os  registradores  
salvos  anteriormente  do  novo  thread;  o  fato  de  o  ponteiro  de  pilha  e  o  contador  de  programa  serem  salvos  e  restaurados  
significa  que  a  CPU  alternará  as  pilhas  e  o  código  que  está  executando.
 A  função  swtch  executa  salvamentos  e  restaurações  para  uma  troca  de  thread  do  kernel.  swtch  não  sabe  diretamente  
sobre  threads;  ele  apenas  salva  e  restaura  conjuntos  de  32  registradores  RISC-V,  chamados  contextos.
 Quando  chega  a  hora  de  um  processo  liberar  a  CPU,  a  thread  do  kernel  do  processo  chama  swtch  para  salvar  seu  próprio  
contexto  e  retornar  ao  contexto  do  escalonador.  Cada  contexto  está  contido  em  uma  estrutura  context  (kernel/proc.h:2).  
contido  na  struct  proc  de  um  processo  ou  na  struct  cpu  de  uma  CPU .  swtch  recebe  dois  argumentos:  struct  context  *old  e  
struct  context  *new.  Ele  salva  os  registradores  atuais  em  old,  carrega  os  registradores  de  new  e  retorna.
 Vamos  acompanhar  um  processo  através  de  swtch  até  o  escalonador.  Vimos  no  Capítulo  4  que  uma  
possibilidade  ao  final  de  uma  interrupção  é  que  a  usertrap  chame  yield.  O  yield,  por  sua  vez,  chama  sched,  
que  chama  swtch  para  salvar  o  contexto  atual  em  p->context  e  alternar  para  o  contexto  do  escalonador  salvo  
anteriormente  em  cpu->context  (kernel/proc.c:497).
 interruptor  (kernel/swtch.S:3)  Salva  apenas  os  registradores  salvos  pelo  chamador;  o  compilador  C  gera  código  no  
chamador  para  salvar  os  registradores  salvos  pelo  chamador  na  pilha.  O  swtch  sabe  o  deslocamento  do  campo  de  cada  
registrador  no  contexto  da  estrutura.  Ele  não  salva  o  contador  do  programa.  Em  vez  disso,  o  swtch  salva  o  registrador  ra ,  
que  contém  o  endereço  de  retorno  de  onde  o  swtch  foi  chamado.  Agora,  o  swtch  restaura  os  registradores  de
 72
Machine Translated by Google
 O  novo  contexto,  que  contém  valores  de  registradores  salvos  por  um  swtch  anterior.  Quando  swtch  retorna,  ele  retorna  
às  instruções  apontadas  pelo  registrador  ra  restaurado ,  ou  seja,  a  instrução  da  qual  a  nova  thread  chamou  swtch  
anteriormente.  Além  disso,  ele  retorna  para  a  pilha  da  nova  thread,  pois  é  para  lá  que  o  sp  restaurado  aponta.
 Em  nosso  exemplo,  o  sched  chamou  swtch  para  alternar  para  cpu->context,  o  contexto  do  escalonador  por  CPU.  
Esse  contexto  foi  salvo  no  momento  em  que  o  escalonador  chamou  swtch  (kernel /proc.c:463).  para  alternar  para  o  
processo  que  está  desistindo  da  CPU.  Quando  o  switch  que  estávamos  rastreando  retorna,  ele  não  retorna  para  sched ,  
mas  para  scheduler,  com  o  ponteiro  de  pilha  na  pilha  do  scheduler  da  CPU  atual.
 7.3  Código:  Agendamento
 A  última  seção  analisou  os  detalhes  de  baixo  nível  do  swtch;  agora,  vamos  considerar  o  swtch  como  um  dado  e  examinar  
a  comutação  de  uma  thread  do  kernel  de  um  processo  para  outro,  através  do  escalonador.  O  escalonador  existe  na  
forma  de  uma  thread  especial  por  CPU,  cada  uma  executando  a  função  de  escalonador .  Esta  função  é  responsável  
por  escolher  qual  processo  executar  em  seguida.  Um  processo  que  deseja  abrir  mão  da  CPU  deve  adquirir  seu  próprio  
bloqueio  de  processo  p->lock,  liberar  quaisquer  outros  bloqueios  que  esteja  mantendo,  atualizar  seu  próprio  estado  (p
>state)  e,  em  seguida,  chamar  sched.  Você  pode  ver  esta  sequência  em  yield  (kernel/proc.c:503).  dormir  e  sair.  sched  
verifica  novamente  alguns  desses  requisitos  (kernel/proc.c:487-492)  e  então  verifica  uma  implicação:  como  um  bloqueio  
é  mantido,  as  interrupções  devem  ser  desabilitadas.  Por  fim,  sched  chama  swtch  para  salvar  o  contexto  atual  em  p
>context  e  alternar  para  o  contexto  do  escalonador  em  cpu->context.  swtch  retorna  na  pilha  do  escalonador  como  se  o  
swtch  do  escalonador  tivesse  retornado  (kernel/proc.c:463).  O  agendador  continua  seu  loop  for ,  encontra  um  processo  
para  executar,  alterna  para  ele  e  o  ciclo  se  repete.
 Acabamos  de  ver  que  o  xv6  mantém  p->lock  em  chamadas  para  swtch:  o  chamador  de  swtch  já  deve  manter  o  
bloqueio,  e  o  controle  do  bloqueio  passa  para  o  código  comutado.  Essa  convenção  é  incomum  com  bloqueios;  
geralmente,  a  thread  que  adquire  um  bloqueio  também  é  responsável  por  liberá-lo,  o  que  facilita  o  raciocínio  sobre  a  
correção.  Para  troca  de  contexto,  é  necessário  quebrar  essa  convenção  porque  p->lock  protege  invariantes  nos  campos  
de  estado  e  contexto  do  processo  que  não  são  verdadeiros  durante  a  execução  em  swtch.  Um  exemplo  de  um  problema  
que  poderia  surgir  se  p->lock  não  fosse  mantido  durante  swtch:  uma  CPU  diferente  poderia  decidir  executar  o  processo  
após  yield  ter  definido  seu  estado  como  RUNNABLE,  mas  antes  de  swtch  fazer  com  que  ele  parasse  de  usar  sua  
própria  pilha  do  kernel.  O  resultado  seriam  duas  CPUs  rodando  na  mesma  pilha,  o  que  causaria  o  caos.
 O  único  lugar  onde  uma  thread  do  kernel  cede  sua  CPU  é  no  sched,  e  ela  sempre  alterna  para  o  mesmo  local  no  
escalonador,  que  (quase)  sempre  alterna  para  alguma  thread  do  kernel  que  anteriormente  chamava  o  sched.  Portanto,  
se  alguém  imprimisse  os  números  das  linhas  onde  o  xv6  alterna  as  threads,  observaria  o  seguinte  padrão  simples:  
(kernel/proc.c:463),  (kernel/proc.c:497),  (kernel/proc.c:463),  (kernel/proc.c:497),  e  assim  por  diante.  Procedimentos  que  
intencionalmente  transferem  o  controle  entre  si  por  meio  de  troca  de  threads  são  às  vezes  chamados  de  corrotinas;  
neste  exemplo,  sched  e  scheduler  são  corrotinas  uma  da  outra.
 Há  um  caso  em  que  a  chamada  do  planejador  para  swtch  não  termina  em  sched.  allocproc  define  o  registro  ra  de  
contexto  de  um  novo  processo  para  forkret  (kernel/proc.c:515),  para  que  sua  primeira  troca
 73
Machine Translated by Google
 “retorna”  ao  início  dessa  função.  forkret  existe  para  liberar  o  p->lock;  caso  contrário,  como  o  novo  processo  precisa  retornar  
ao  espaço  do  usuário  como  se  estivesse  retornando  de  fork,  ele  poderia  começar  em  usertrapret.  scheduler  (kernel/
 proc.c:445)  executa  
um  loop:  encontre  um  processo  para  executar,  execute-o  até  que  ele  ceda  e  repita.
 O  escalonador  percorre  a  tabela  de  processos  em  busca  de  um  processo  executável,  que  tenha  p->state  ==  
RUNNABLE.  Ao  encontrar  um  processo,  ele  define  a  variável  de  processo  atual  por  CPU  c->proc,  marca  o  
processo  como  EM  EXECUÇÃO  e,  em  seguida,  chama  swtch  para  iniciar  sua  execução  (kernel/proc.c:458-463).
 Uma  maneira  de  pensar  sobre  a  estrutura  do  código  de  escalonamento  é  que  ele  impõe  um  conjunto  de  invariantes  
sobre  cada  processo  e  mantém  p->lock  sempre  que  essas  invariantes  não  forem  verdadeiras.  Uma  invariante  é  que,  se  um  
processo  estiver  EM  EXECUÇÃO,  o  yield  de  uma  interrupção  de  temporizador  deve  ser  capaz  de  alternar  com  segurança  
para  longe  do  processo;  isso  significa  que  os  registradores  da  CPU  devem  conter  os  valores  de  registrador  do  processo  (ou  
seja,  swtch  não  os  moveu  para  um  contexto)  e  c->proc  deve  se  referir  ao  processo.  Outra  invariante  é  que,  se  um  processo  
for  EXECUTÁVEL,  deve  ser  seguro  para  o  escalonador  de  uma  CPU  ociosa  executá-lo;  isso  significa  que  p->context  deve  
conter  os  registradores  do  processo  (ou  seja,  eles  não  estão  realmente  nos  registradores  reais),  que  nenhuma  CPU  está  
executando  na  pilha  do  kernel  do  processo  e  que  nenhum  c->proc  da  CPU  se  refere  ao  processo.  Observe  que  essas  
propriedades  geralmente  não  são  verdadeiras  enquanto  p->lock  é  mantido.
 Manter  as  invariantes  acima  é  o  motivo  pelo  qual  o  xv6  frequentemente  adquire  p->lock  em  uma  thread  e  o  libera  em  
outra,  por  exemplo,  adquirindo  em  yield  e  liberando  em  scheduler.  Uma  vez  que  yield  tenha  começado  a  modificar  o  estado  
de  um  processo  em  execução  para  torná-lo  RUNNABLE,  o  bloqueio  deve  permanecer  mantido  até  que  as  invariantes  sejam  
restauradas:  o  primeiro  ponto  de  liberação  correto  é  após  o  scheduler  (executando  em  sua  própria  pilha)  limpar  c->proc.  Da  
mesma  forma,  uma  vez  que  o  scheduler  começa  a  converter  um  processo  RUNNABLE  para  RUNNING,  o  bloqueio  não  
pode  ser  liberado  até  que  a  thread  do  kernel  esteja  completamente  em  execução  (após  a  troca,  por  exemplo  em  yield).
 7.4  Código:  mycpu  e  myproc
 O  Xv6  frequentemente  precisa  de  um  ponteiro  para  a  estrutura  proc  do  processo  atual .  Em  um  processador  único,  pode-se  
ter  uma  variável  global  apontando  para  o  proc  atual.  Isso  não  funciona  em  uma  máquina  com  vários  núcleos,  já  que  cada  
núcleo  executa  um  processo  diferente.  A  maneira  de  resolver  esse  problema  é  explorar  o  fato  de  que  cada  núcleo  tem  seu  
próprio  conjunto  de  registradores;  podemos  usar  um  desses  registradores  para  ajudar  a  encontrar  informações  por  núcleo.
 O  Xv6  mantém  uma  estrutura  cpu  para  cada  CPU  (kernel/proc.h:22),  que  registra  o  processo  em  execução  naquela  
CPU  (se  houver),  os  registradores  salvos  para  a  thread  do  escalonador  da  CPU  e  a  contagem  de  spinlocks  aninhados  
necessários  para  gerenciar  a  desabilitação  de  interrupções.  A  função  mycpu  (kernel/proc.c:74)  O  Xv6  retorna  um  ponteiro  
para  a  estrutura  cpu  da  CPU  atual.  O  RISC-V  numera  suas  CPUs,  atribuindo  a  cada  uma  um  hartid.  O  Xv6  garante  que  o  
hartid  de  cada  CPU  seja  armazenado  no  registrador  tp  dessa  CPU  enquanto  estiver  no  kernel.
 Isso  permite  que  o  mycpu  use  o  tp  para  indexar  uma  matriz  de  estruturas  de  CPU  para  encontrar  a  correta.
 Garantir  que  o  tp  de  uma  CPU  sempre  mantenha  o  hartid  da  CPU  é  um  pouco  complicado.  start  define  o  
registrador  tp  no  início  da  sequência  de  inicialização  da  CPU,  enquanto  ainda  está  no  modo  de  máquina  (kernel/
 start.c:51).  usertrapret  salva  tp  na  página  do  trampolim,  pois  o  processo  do  usuário  pode  modificá  -lo.  Por  fim,  
uservec  restaura  esse  tp  salvo  ao  entrar  no  kernel  a  partir  do  espaço  do  usuário  (kernel/trampoline.S:77).
 O  compilador  garante  nunca  usar  o  registrador  tp .  Seria  mais  conveniente  se  o  xv6  pudesse
 74
Machine Translated by Google
 pergunte  ao  hardware  RISC-V  sobre  o  hartid  atual  sempre  que  necessário,  mas  o  RISC-V  permite  isso  apenas  em
 modo  máquina,  não  no  modo  supervisor.
 Os  valores  de  retorno  de  cpuid  e  mycpu  são  frágeis:  se  o  temporizador  interrompesse  e  causasse  o
 thread  para  render  e  então  mover  para  uma  CPU  diferente,  um  valor  retornado  anteriormente  não  seria  mais
 correto.  Para  evitar  esse  problema,  o  xv6  requer  que  os  chamadores  desabilitem  as  interrupções  e  as  habilitem  apenas
 depois  que  eles  terminam  de  usar  a  estrutura  retornada  cpu.
 A  função  myproc  (kernel/proc.c:83)  retorna  o  ponteiro  struct  proc  para  o  processo  que
 está  sendo  executado  na  CPU  atual.  myproc  desabilita  interrupções,  invoca  mycpu,  busca  a  CPU  atual
 ponteiro  de  processo  (c->proc)  para  fora  da  estrutura  da  CPU  e,  em  seguida,  habilita  interrupções.  O  valor  de  retorno
 de  myproc  é  seguro  de  usar  mesmo  se  as  interrupções  estiverem  habilitadas:  se  uma  interrupção  do  temporizador  move  o  processo  de  chamada
 para  uma  CPU  diferente,  seu  ponteiro  struct  proc  permanecerá  o  mesmo.
 7.5  Sono  e  despertar
 Agendamentos  e  bloqueios  ajudam  a  ocultar  as  ações  de  uma  thread  de  outra,  mas  também  precisamos  de  abstrações  que  
ajudem  as  threads  a  interagir  intencionalmente.  Por  exemplo,  o  leitor  de  um  pipe  em  xv6  pode
 precisa  esperar  que  um  processo  de  escrita  produza  dados;  a  chamada  de  um  pai  para  esperar  pode  precisar  esperar  por  um
 criança  para  sair;  e  um  processo  que  lê  o  disco  precisa  esperar  que  o  hardware  do  disco  conclua  a  leitura.
 O  kernel  xv6  usa  um  mecanismo  chamado  suspensão  e  ativação  nessas  situações  (e  em  muitas  outras).
 O  modo  de  suspensão  permite  que  uma  thread  do  kernel  aguarde  um  evento  específico;  outra  thread  pode  chamar  o  modo  de  ativação  para  indicar  
que  as  threads  que  aguardam  um  evento  devem  ser  retomadas.  O  modo  de  suspensão  e  a  ativação  são  frequentemente  chamados  de  sequências.
 mecanismos  de  coordenação  ou  sincronização  condicional.
 O  sono  e  a  vigília  fornecem  uma  interface  de  sincronização  de  nível  relativamente  baixo.  Para  motivar  o
 como  eles  funcionam  no  xv6,  nós  os  usaremos  para  construir  um  mecanismo  de  sincronização  de  nível  superior  chamado
 um  semáforo  [5]  que  coordena  produtores  e  consumidores  (xv6  não  usa  semáforos).  A
 O  semáforo  mantém  uma  contagem  e  fornece  duas  operações.  A  operação  “V”  (para  o  produtor)
 incrementa  a  contagem.  A  operação  “P”  (para  o  consumidor)  espera  até  que  a  contagem  seja  diferente  de  zero,
 e  então  decrementa  e  retorna.  Se  houvesse  apenas  um  thread  produtor  e  um  consumidor
 thread,  e  eles  foram  executados  em  CPUs  diferentes,  e  o  compilador  não  otimizou  de  forma  muito  agressiva,
 esta  implementação  estaria  correta:
 100  estrutura  semáforo  {
 101  estrutura  de  bloqueio  spinlock;
 102
 103
 104
 105  vazio
 contagem  de  int;
 };
 106  V(estrutura  semáforo  *s)
 107  {
 108
 109
 110
 111
 adquirir(&s->bloquear);
 s->contagem  +=  1;
 liberar(&s->bloquear);
 }
 75
Machine Translated by Google
 112
 113  vazio
 114  P(struct  semáforo  *s)
 115  {
 116
 117
 118  
119  
120
 121
 enquanto(s->contagem  ==  0)
 ;
 adquirir(&s->bloquear);
 s->contagem  -=  1;
 liberar(&s->bloquear);
 }
 A  implementação  acima  é  cara.  Se  o  produtor  agir  raramente,  o  consumidor  gastará
 a  maior  parte  do  tempo  girando  no  loop  while  esperando  por  uma  contagem  diferente  de  zero.  A  CPU  do  consumidor
 provavelmente  conseguiria  encontrar  trabalho  mais  produtivo  do  que  ficar  ocupado  esperando,  consultando  repetidamente  s->count.
 Evitar  a  espera  ocupada  requer  uma  maneira  para  o  consumidor  ceder  a  CPU  e  retomar  somente  após  V
 incrementa  a  contagem.
 Aqui  está  um  passo  nessa  direção,  embora,  como  veremos,  não  seja  suficiente.  Imaginemos  um  par
 de  chamadas,  sleep  e  wakeup,  que  funcionam  da  seguinte  maneira.  sleep(chan)  dorme  no  valor  arbitrário
 chan,  chamado  de  canal  de  espera.  sleep  coloca  o  processo  de  chamada  em  hibernação,  liberando  a  CPU  para  outros
 trabalho.  wakeup(chan)  desperta  todos  os  processos  que  estão  dormindo  no  chan  (se  houver),  fazendo  com  que  suas  chamadas  de  sono  sejam
 retornar.  Se  não  houver  processos  aguardando  o  canal,  o  wakeup  não  faz  nada.  Podemos  alterar  o  semáforo
 implementação  para  usar  sleep  e  wakeup  (mudanças  destacadas  em  amarelo):
 200  vazio
 201  V(struct  semáforo  *s)
 202  {
 203
 204
 205
 206  
207
 208
 209  nulo
 adquirir(&s->bloquear);
 s->contagem  +=  1;
 despertar(es);
 liberar(&s->bloquear);
 }
 210  P(estrutura  semáforo  *s)
 211  {
 212
 213
 214
 215  
216
 217
 enquanto(s->contagem  ==  0)
 sono(s);
 adquirir(&s->bloquear);
 s->contagem  -=  1;
 liberar(&s->bloquear);
 }
 P  agora  desiste  da  CPU  em  vez  de  girar,  o  que  é  bom.  No  entanto,  não  é
 fácil  de  projetar  o  sono  e  o  despertar  com  esta  interface  sem  sofrer  com  o  que  é
 conhecido  como  o  problema  do  despertar  perdido.  Suponha  que  P  encontre  s->count  ==  0  na  linha  212.  Enquanto
 P  está  entre  as  linhas  212  e  213,  V  roda  em  outra  CPU:  ele  muda  s->count  para  ser  diferente  de  zero  e
 76
Machine Translated by Google
 chama  wakeup,  que  não  encontra  nenhum  processo  dormindo  e,  portanto,  não  faz  nada.  Agora  P  continua  executando
 na  linha  213:  ele  chama  sleep  e  entra  em  modo  sleep.  Isso  causa  um  problema:  P  está  dormindo  esperando  uma  chamada  de  V
 isso  já  aconteceu.  A  menos  que  tenhamos  sorte  e  o  produtor  ligue  para  V  novamente,  o  consumidor  irá
 esperar  para  sempre,  mesmo  que  a  contagem  seja  diferente  de  zero.
 A  raiz  deste  problema  é  que  a  invariante  de  que  P  só  dorme  quando  s->count  ==  0  é  violada
 por  V  correndo  no  momento  errado.  Uma  maneira  incorreta  de  proteger  a  invariante  seria  mover
 a  aquisição  do  bloqueio  (destacada  em  amarelo  abaixo)  em  P  para  que  sua  verificação  da  contagem  e  sua  chamada  para
 o  sono  é  atômico:
 300  vazio
 301  V(struct  semáforo  *s)
 302  {
 303
 304
 305
 306
 307
 308
 309  nulo
 adquirir(&s->bloquear);
 s->contagem  +=  1;
 despertar(es);
 liberar(&s->bloquear);
 }
 310  P(estrutura  semáforo  *s)
 311  {
 312
 313  
314  
315  
316
 317
 adquirir(&s->bloquear);
 enquanto(s->contagem  ==  0)
 sono(s);
 s->contagem  -=  1;
 liberar(&s->bloquear);
 }
 Poderíamos  esperar  que  esta  versão  de  P  evitasse  o  despertar  perdido  porque  o  bloqueio  impede  V
 de  executar  entre  as  linhas  313  e  314.  Ele  faz  isso,  mas  também  gera  um  deadlock:  P  mantém  o  bloqueio
 enquanto  ele  dorme,  então  V  ficará  bloqueado  para  sempre  esperando  o  bloqueio.
 Vamos  corrigir  o  esquema  anterior  alterando  a  interface  do  sleep :  o  chamador  deve  passar  o  bloqueio  de  condição  para  sleep  para  
que  ele  possa  liberar  o  bloqueio  depois  que  o  processo  de  chamada  for  marcado  como  sleep  e
 aguardando  no  canal  de  suspensão.  O  bloqueio  forçará  um  V  concorrente  a  esperar  até  que  P  termine  de  colocar
 para  dormir,  para  que  o  despertador  encontre  o  consumidor  adormecido  e  o  acorde.  Assim  que  o  consumidor  acorda  novamente,  o  sono  
recupera  o  bloqueio  antes  de  retornar.  Nosso  novo  modo  correto  de  sono/despertar
 o  esquema  pode  ser  utilizado  da  seguinte  forma  (alteração  destacada  em  amarelo):
 400  vazio
 401  V(struct  semáforo  *s)
 402  {
 403
 404
 405
 406
 407
 adquirir(&s->bloquear);
 s->contagem  +=  1;
 despertar(es);
 liberar(&s->bloquear);
 }
 77
Machine Translated by Google
 408
 409  nulo
 410  P(struct  semáforo  *s)  411  {
 412
 413
 414
 415
 416
 417
 adquirir(&s->bloquear);  
enquanto(s->contagem  ==  0)
 dormir(s,  &s->bloqueio);  s
>contagem  -=  1;  liberar(&s
>bloqueio);
 }
 O  fato  de  P  conter  s->lock  impede  que  V  tente  acordá-lo  entre  a  verificação  de  s->count  por  P  e  sua  
chamada  para  dormir.  Observe,  no  entanto,  que  precisamos  de  dormir  para  liberar  atomicamente  s->lock  e  
colocar  o  processo  consumidor  em  sono,  a  fim  de  evitar  despertares  perdidos.
 7.6  Código:  Sono  e  despertar
 Suspensão  do  Xv6  (kernel/proc.c:536)  e  despertar  (kernel/proc.c:567)  Fornece  a  interface  mostrada  no  último  
exemplo  acima,  e  sua  implementação  (além  das  regras  de  como  usá-las)  garante  que  não  haja  despertares  
perdidos.  A  ideia  básica  é  fazer  com  que  o  sleep  marque  o  processo  atual  como  SLEEPING  e,  em  seguida,  
chamar  sched  para  liberar  a  CPU;  o  wakeup  procura  um  processo  dormindo  no  canal  de  espera  fornecido  e  o  
marca  como  RUNNABLE.  Chamadores  de  sleep  e  wakeup  podem  usar  qualquer  número  mutuamente  
conveniente  como  canal.  O  Xv6  frequentemente  usa  o  endereço  de  uma  estrutura  de  dados  do  kernel  envolvida  na  espera.
 sleep  adquire  p->lock  (kernel/proc.c:547).  Agora  o  processo  que  vai  dormir  mantém  ambos  p->lock
 e  lk.  Manter  lk  era  necessário  no  chamador  (no  exemplo,  P):  isso  garantia  que  nenhum  outro  processo  (no  
exemplo,  um  V  em  execução)  pudesse  iniciar  uma  chamada  para  wakeup(chan).  Agora  que  sleep  mantém  p
>lock,  é  seguro  liberar  lk:  algum  outro  processo  pode  iniciar  uma  chamada  para  wakeup(chan),  mas  wakeup  
aguardará  para  adquirir  p->lock  e,  portanto,  aguardará  até  que  sleep  termine  de  colocar  o  processo  para  dormir,  
evitando  que  wakeup  perca  o  sleep.
 Agora  que  o  sleep  contém  p->lock  e  nenhum  outro,  ele  pode  colocar  o  processo  para  dormir  gravando  o  
canal  de  sono,  alterando  o  estado  do  processo  para  SLEEPING  e  chamando  sched  (kernel/proc.c:551-554).
 Em  breve  ficará  claro  por  que  é  essencial  que  p->lock  não  seja  liberado  (pelo  agendador)  até  que  o  processo  
seja  marcado  como  SLEEPING.
 Em  algum  momento,  um  processo  adquirirá  o  bloqueio  de  condição,  definirá  a  condição  que  o  sleeper  está  aguardando  e  
chamará  wakeup(chan).  É  importante  que  wakeup  seja  chamado  enquanto  a  condição  lock1  estiver  sendo  mantida .  O  wakeup  
percorre  a  tabela  de  processos  (kernel/proc.c:567).  Ele  adquire  o  p->lock  de  cada  processo  que  inspeciona,  tanto  porque  pode  
manipular  o  estado  desse  processo  quanto  porque  o  p->lock  garante  que  o  modo  de  suspensão  e  o  modo  de  ativação  não  se  percam.  
Quando  o  modo  de  ativação  encontra  um  processo  no  estado  SLEEPING  com  um  canal  correspondente,  ele  altera  o  estado  desse  
processo  para  RUNNABLE.  Na  próxima  vez  que  o  agendador  for  executado,  ele  verá  que  o  processo  está  pronto  para  ser  executado.
 1A  rigor,  é  suficiente  que  o  despertar  siga  apenas  a  aquisição  (ou  seja,  pode-se  chamar  o  despertar  após  a  liberação).
 78
Machine Translated by Google
 Por  que  as  regras  de  bloqueio  para  sleep  e  wakeup  garantem  que  um  processo  em  modo  sleep  não  perca  
um  wakeup?  O  processo  em  modo  sleep  mantém  o  bloqueio  de  condição  ou  seu  próprio  p->lock ,  ou  ambos,  de  
um  ponto  antes  de  verificar  a  condição  até  um  ponto  depois  de  ser  marcado  como  SLEEPING.  O  processo  que  
chama  o  wakeup  mantém  ambos  os  bloqueios  no  loop  do  wakeup.  Assim,  o  processo  que  o  ativa  torna  a  
condição  verdadeira  antes  que  a  thread  consumidora  a  verifique;  ou  o  processo  que  o  ativa  examina  a  thread  
em  modo  sleep  estritamente  após  ela  ter  sido  marcada  como  SLEEPING.  Então,  o  processo  em  modo  wakeup  
verá  o  processo  em  modo  sleep  e  o  ativará  (a  menos  que  algo  o  ative  primeiro).
 Às  vezes,  vários  processos  estão  dormindo  no  mesmo  canal;  por  exemplo,  mais  de  um  processo  lendo  de  
um  pipe.  Uma  única  chamada  para  wakeup  despertará  todos  eles.  Um  deles  será  executado  primeiro  e  obterá  o  
bloqueio  com  o  qual  sleep  foi  chamado  e  (no  caso  de  pipes)  lerá  quaisquer  dados  que  estejam  aguardando  no  
pipe.  Os  outros  processos  descobrirão  que,  apesar  de  terem  sido  despertados,  não  há  dados  para  serem  lidos.  
Do  ponto  de  vista  deles,  o  wakeup  foi  "espúrio"  e  eles  devem  dormir  novamente.  Por  esse  motivo,  sleep  é  sempre  
chamado  dentro  de  um  loop  que  verifica  a  condição.
 Não  há  problema  algum  se  dois  usuários  de  sleep/wakeup  escolherem  acidentalmente  o  mesmo  canal:  eles  
verão  despertares  espúrios,  mas  o  looping  descrito  acima  tolerará  esse  problema.  Grande  parte  do  charme  de  
sleep/wakeup  reside  em  sua  leveza  (não  há  necessidade  de  criar  estruturas  de  dados  especiais  para  atuar  como  
canais  de  sleep)  e  em  fornecer  uma  camada  de  indireção  (os  chamadores  não  precisam  saber  com  qual  processo  
específico  estão  interagindo).
 7.7  Código:  Pipes
 Um  exemplo  mais  complexo  que  usa  sleep  e  wakeup  para  sincronizar  produtores  e  consumidores  é  a  
implementação  de  pipes  do  xv6.  Vimos  a  interface  para  pipes  no  Capítulo  1:  bytes  gravados  em  uma  extremidade  
de  um  pipe  são  copiados  para  um  buffer  no  kernel  e,  em  seguida,  podem  ser  lidos  da  outra  extremidade  do  pipe.
 Os  próximos  capítulos  examinarão  o  suporte  ao  descritor  de  arquivo  em  torno  dos  pipes,  mas  vamos  agora  dar  
uma  olhada  nas  implementações  de  pipewrite  e  piperead.
 Cada  pipe  é  representado  por  um  pipe  struct,  que  contém  um  bloqueio  e  um  buffer  de  dados .
 Os  campos  nread  e  nwrite  contam  o  número  total  de  bytes  lidos  e  gravados  no  buffer.  O  buffer  é  encapsulado:  o  
próximo  byte  escrito  após  buf[PIPESIZE-1]  é  buf[0].  As  contagens  não  são  encapsuladas.  Essa  convenção  
permite  que  a  implementação  distinga  um  buffer  cheio  (nwrite  ==  nread+PIPESIZE)  de  um  buffer  vazio  (nwrite  
==  nread),  mas  significa  que  a  indexação  no  buffer  deve  usar  buf[nread  %  PIPESIZE]  em  vez  de  apenas  
buf[nread]  (e  similarmente  para  nwrite).
 Vamos  supor  que  as  chamadas  para  piperead  e  pipewrite  aconteçam  simultaneamente  em  duas  CPUs  
diferentes.  pipewrite  (kernel/pipe.c:77)  começa  adquirindo  o  bloqueio  do  pipe,  que  protege  as  contagens,  os  
dados  e  seus  invariantes  associados.  piperead  (kernel/pipe.c:106)  Então  tenta  adquirir  o  bloqueio  também,  mas  
não  consegue.  Ele  gira  em  adquirir  (kernel/spinlock.c:22)  aguardando  o  bloqueio.  Enquanto  piperead  aguarda,  
pipewrite  percorre  os  bytes  que  estão  sendo  escritos  (addr[0..n-1]),  adicionando  cada  um  ao  pipe  por  vez  (kernel/
 pipe.c:95).  Durante  esse  loop,  pode  acontecer  que  o  buffer  fique  cheio  (kernel/pipe.c:88).  Neste  caso,  o  pipewrite  
chama  wakeup  para  alertar  os  leitores  adormecidos  sobre  o  fato  de  que  há  dados  aguardando  no  buffer  e,  então,  
entra  em  modo  sleep  em  &pi->nwrite  para  esperar  que  um  leitor  retire  alguns  bytes  do  buffer.  O  sleep  libera  pi
>lock  como  parte  do  processo  de  colocar  o  processo  do  pipewrite  em  modo  sleep.
 79
Machine Translated by Google
 Agora  que  pi->lock  está  disponível,  piperead  consegue  adquiri-lo  e  entra  em  sua  seção  crítica:  ele  descobre  que  pi
>nread !=  pi->nwrite  (kernel/pipe.c:113)  (pipewrite  entrou  em  modo  de  espera  porque  pi->nwrite  ==  pi->nread+PIPESIZE  
(kernel/pipe.c:88)),  então  ele  cai  no  loop  for ,  copia  os  dados  do  pipe  (kernel/pipe.c:120),  e  incrementa  nread  pelo  número  
de  bytes  copiados.  Essa  quantidade  de  bytes  agora  está  disponível  para  escrita,  então  piperead  chama  wakeup  ( kernel/
 pipe.c:127)  para  despertar  quaisquer  escritores  adormecidos  antes  que  ele  retorne.  O  comando  wakeup  encontra  um  
processo  adormecido  em  &pi->nwrite,  o  processo  que  estava  executando  o  pipewrite,  mas  parou  quando  o  buffer  ficou  
cheio.  Ele  marca  esse  processo  como  EXECUTÁVEL.
 O  código  do  pipe  usa  canais  de  suspensão  separados  para  leitor  e  gravador  (pi->nread  e  pi->nwrite);
 Isso  pode  tornar  o  sistema  mais  eficiente  no  caso  improvável  de  haver  muitos  leitores  e  escritores  
aguardando  o  mesmo  pipe.  O  código  do  pipe  dorme  dentro  de  um  loop  que  verifica  a  condição  de  sono;  se  
houver  vários  leitores  ou  escritores,  todos  os  processos,  exceto  o  primeiro  a  acordar,  verão  que  a  condição  
ainda  é  falsa  e  dormirão  novamente.
 7.8  Código:  Espere,  saia  e  mate
 sleep  e  wakeup  podem  ser  usados  para  muitos  tipos  de  espera.  Um  exemplo  interessante,  apresentado  no  Capítulo  1,  é  
a  interação  entre  a  saída  de  um  filho  e  a  espera  de  seu  pai .  No  momento  da  morte  do  filho,  o  pai  pode  já  estar  dormindo  
em  espera  ou  fazendo  outra  coisa;  neste  último  caso,  uma  chamada  subsequente  para  wait  deve  observar  a  morte  do  
filho,  talvez  muito  tempo  depois  de  chamar  exit.  A  maneira  como  o  xv6  registra  a  morte  do  filho  até  que  wait  a  observe  é  
para  exit  colocar  o  chamador  no  estado  ZOMBIE ,  onde  permanece  até  que  wait  do  pai  perceba,  altere  o  estado  do  filho  
para  UNUSED,  copie  o  status  de  saída  do  filho  e  retorne  o  ID  do  processo  do  filho  ao  pai.  Se  o  pai  sair  antes  do  filho,  o  
pai  entrega  o  filho  ao  processo  init ,  que  chama  wait  perpetuamente;  assim,  cada  filho  tem  um  pai  para  limpar  depois  
dele.  Um  desafio  é  evitar  disputas  e  deadlocks  entre  esperas  e  saídas  simultâneas  de  pai  e  filho,  bem  como  saídas  e  
saídas  simultâneas.
 wait  começa  adquirindo  wait_lock  (kernel/proc.c:391).  O  motivo  é  que  wait_lock  atua  como  um  bloqueio  
de  condição  que  ajuda  a  garantir  que  o  pai  não  perca  um  despertar  de  um  filho  existente.
 Em  seguida,  wait  varre  a  tabela  de  processos.  Se  encontrar  um  filho  no  estado  ZOMBIE,  libera  os  recursos  desse  filho  e  
sua  estrutura  proc ,  copia  o  status  de  saída  do  filho  para  o  endereço  fornecido  para  wait  (se  não  for  0)  e  retorna  o  ID  do  
processo  do  filho.  Se  wait  encontrar  filhos,  mas  nenhum  tiver  saído,  chama  sleep  para  aguardar  que  algum  deles  saia  
(kernel/proc.c:433).  então  verifica  novamente.  wait  geralmente  contém  dois  bloqueios,  wait_lock  e  pp->lock  de  algum  
processo ;  a  ordem  para  evitar  deadlock  é  primeiro  wait_lock  e  depois  pp->lock.  exit  (kernel/proc.c:347)  Registra  o  status  
de  saída,  libera  
alguns  recursos,  chama  reparent  para  entregar  seus  filhos  ao  processo  init ,  desperta  o  pai  caso  esteja  em  espera,  
marca  o  chamador  como  zumbi  e  libera  a  CPU  permanentemente.  exit  mantém  wait_lock  e  p->lock  durante  esta  
sequência.  Ele  mantém  wait_lock  porque  é  o  bloqueio  condicional  para  o  wakeup(p->parent),  evitando  que  um  pai  em  
espera  perca  o  wakeup.  exit  também  deve  manter  p->lock  para  esta  sequência,  para  evitar  que  um  pai  em  espera  veja  
que  o  filho  está  no  estado  ZOMBIE  antes  que  o  filho  finalmente  chame  swtch.  exit  adquire  esses  bloqueios  na  mesma  
ordem  que  wait  para  evitar  deadlock.
 80
Machine Translated by Google
 Pode  parecer  incorreto  que  exit  acorde  o  pai  antes  de  definir  seu  estado  como  ZOMBIE,  mas  isso  é  seguro:  embora  
wakeup  possa  fazer  com  que  o  pai  seja  executado,  o  loop  em  wait  não  pode  examinar  o  filho  até  que  o  p->lock  do  filho  seja  
liberado  pelo  planejador,  então  wait  não  pode  examinar  o  processo  de  saída  até  bem  depois  que  exit  tenha  definido  seu  
estado  como  ZOMBIE  (kernel/proc.c:379).
 Enquanto  exit  permite  que  um  processo  se  encerre,  kill  (kernel/proc.c:586)  permite  que  um  processo  solicite  o  término  
de  outro.  Seria  muito  complexo  para  kill  destruir  diretamente  o  processo  da  vítima,  já  que  a  vítima  pode  estar  executando  em  
outra  CPU,  talvez  no  meio  de  uma  sequência  sensível  de  atualizações  nas  estruturas  de  dados  do  kernel.  Portanto,  kill  faz  
muito  pouco:  apenas  define  o  p->killed  da  vítima  e,  se  estiver  em  modo  de  espera,  a  desperta.  Eventualmente,  a  vítima  
entrará  ou  sairá  do  kernel,  momento  em  que  o  código  em  usertrap  chamará  exit  se  p->killed  estiver  definido  (ele  verifica  
chamando  killed  (kernel/proc.c:615)).  Se  a  vítima  estiver  sendo  executada  no  espaço  do  usuário,  ela  logo  entrará  no  kernel  
fazendo  uma  chamada  de  sistema  ou  porque  o  timer  (ou  algum  outro  dispositivo)  interrompe.
 Se  o  processo  da  vítima  estiver  em  modo  de  espera,  a  chamada  de  ativação  do  comando  kill  fará  com  que  a  vítima  
retorne  do  modo  de  espera.  Isso  é  potencialmente  perigoso,  pois  a  condição  que  está  sendo  aguardada  pode  não  ser  verdadeira.
 No  entanto,  chamadas  xv6  para  sleep  são  sempre  encapsuladas  em  um  loop  while  que  testa  novamente  a  condição  após  o  
retorno  do  sleep .  Algumas  chamadas  para  sleep  também  testam  p->killed  no  loop  e  abandonam  a  atividade  atual  se  ela  
estiver  definida.  Isso  só  é  feito  quando  tal  abandono  for  correto.  Por  exemplo,  o  código  de  leitura  e  gravação  do  pipe  retorna  
se  o  sinalizador  killed  estiver  definido;  eventualmente,  o  código  retornará  para  trap,  que  verificará  novamente  p->killed  e  sairá.
 Alguns  loops  de  suspensão  xv6  não  verificam  p->killed  porque  o  código  está  no  meio  de  uma  chamada  de  
sistema  multietapas  que  deveria  ser  atômica.  O  driver  virtio  (kernel/virtio_disk.c:285)  Um  exemplo:  ele  não  verifica  
p->killed  porque  uma  operação  de  disco  pode  ser  uma  de  um  conjunto  de  gravações  necessárias  para  que  o  
sistema  de  arquivos  permaneça  em  um  estado  correto.  Um  processo  que  é  encerrado  enquanto  aguarda  E/S  de  
disco  não  será  encerrado  até  concluir  a  chamada  de  sistema  atual  e  o  usertrap  visualizar  o  sinalizador  de  encerramento.
 7.9  Bloqueio  de  Processo
 O  bloqueio  associado  a  cada  processo  (p->lock)  é  o  bloqueio  mais  complexo  no  xv6.  Uma  maneira  simples  de  
pensar  em  p->lock  é  que  ele  deve  ser  mantido  durante  a  leitura  ou  gravação  de  qualquer  um  dos  seguintes  
campos  struct  proc :  p->state,  p->chan,  p->killed,  p->xstate  e  p->pid.  Esses  campos  podem  ser  usados  por  outros  
processos  ou  por  threads  do  escalonador  em  outros  núcleos,  portanto,  é  natural  que  sejam  protegidos  por  um  
bloqueio.
 Entretanto,  a  maioria  dos  usos  de  p->lock  protegem  aspectos  de  nível  superior  da  estrutura  de  dados  do  processo  xv6.
 estruturas  e  algoritmos.  Aqui  está  o  conjunto  completo  de  coisas  que  p->lock  faz:
 •  Junto  com  p->state,  ele  evita  corridas  na  alocação  de  slots  proc[]  para  novos  processos.
 •  Ele  oculta  um  processo  da  vista  enquanto  ele  está  sendo  criado  ou  destruído.
 •  Impede  que  a  espera  de  um  pai  colete  um  processo  que  definiu  seu  estado  como  ZOMBIE,  mas  tem
 ainda  não  rendeu  a  CPU.
 •  Impede  que  o  planejador  de  outro  núcleo  decida  executar  um  processo  de  produção  após  definir  seu
 estado  para  RUNNABLE ,  mas  antes  de  terminar,  alterne.
 81
Machine Translated by Google
 •  Garante  que  apenas  o  agendador  de  um  núcleo  decida  executar  um  processo  RUNNABLE .
 •  Evita  que  uma  interrupção  do  temporizador  faça  com  que  um  processo  ceda  enquanto  estiver  em  comutação.
 •  Junto  com  o  bloqueio  de  condição,  ajuda  a  evitar  que  o  despertar  ignore  um  processo  que  está
 chamando  sleep  mas  não  terminou  de  renderizar  a  CPU.
 •  Impede  que  o  processo  de  morte  da  vítima  saia  e  talvez  seja  realocado  entre
 verificação  de  p->pid  e  configuração  de  p->killed.
 •  Torna  a  verificação  e  gravação  de  kill  de  p->state  atômicas.
 O  campo  p->parent  é  protegido  pelo  bloqueio  global  wait_lock  em  vez  de  p->lock.
 Somente  o  processo  pai  modifica  p->parent,  embora  o  campo  seja  lido  tanto  pelo  próprio  processo  quanto  por  
outros  processos  que  buscam  seus  filhos.  O  propósito  de  wait_lock  é  atuar  como  o  bloqueio  condicional  quando  
wait  está  em  modo  sleep,  aguardando  a  saída  de  qualquer  filho.  Um  filho  que  está  saindo  mantém  wait_lock  ou  p
>lock  até  que  seu  estado  seja  definido  como  ZOMBIE,  seu  pai  tenha  despertado  e  a  CPU  tenha  sido  liberada.  
wait_lock  também  serializa  saídas  simultâneas  de  um  pai  e  um  filho,  de  modo  que  o  processo  init  (que  herda  o  
filho)  tenha  a  garantia  de  ser  despertado  de  sua  espera.  wait_lock  é  um  bloqueio  global,  e  não  um  bloqueio  por  
processo  em  cada  pai,  porque,  até  que  um  processo  o  adquira,  ele  não  pode  saber  quem  é  seu  pai.
 7.10  Mundo  real
 O  escalonador  xv6  implementa  uma  política  de  escalonamento  simples,  que  executa  cada  processo  por  vez.  Essa  
política  é  chamada  de  round  robin.  Sistemas  operacionais  reais  implementam  políticas  mais  sofisticadas  que,  por  
exemplo,  permitem  que  os  processos  tenham  prioridades.  A  ideia  é  que  um  processo  executável  de  alta  prioridade  
seja  preferido  pelo  escalonador  a  um  processo  executável  de  baixa  prioridade.  Essas  políticas  podem  se  tornar  
complexas  rapidamente  porque  frequentemente  há  objetivos  conflitantes:  por  exemplo,  o  sistema  operacional  
também  pode  querer  garantir  justiça  e  alto  rendimento.  Além  disso,  políticas  complexas  podem  levar  a  interações  
não  intencionais,  como  inversão  de  prioridade  e  comboios.  A  inversão  de  prioridade  pode  ocorrer  quando  um  
processo  de  baixa  prioridade  e  um  de  alta  prioridade  usam  um  bloqueio  específico,  que,  quando  adquirido  pelo  
processo  de  baixa  prioridade,  pode  impedir  o  processo  de  alta  prioridade  de  progredir.  Um  longo  comboio  de  
processos  em  espera  pode  se  formar  quando  muitos  processos  de  alta  prioridade  aguardam  por  um  processo  de  
baixa  prioridade  que  adquire  um  bloqueio  compartilhado;  uma  vez  formado,  o  comboio  pode  persistir  por  um  longo  
tempo.  Para  evitar  esses  tipos  de  problemas,  mecanismos  adicionais  são  necessários  em  planejadores  sofisticados.
 Sleep  e  Wakeup  são  métodos  de  sincronização  simples  e  eficazes,  mas  existem  muitos  outros.  O  primeiro  
desafio  em  todos  eles  é  evitar  o  problema  de  "perdas  de  ativação"  que  vimos  no  início  do  capítulo.  A  função  sleep  
do  kernel  Unix  original  simplesmente  desabilitava  interrupções,  o  que  era  suficiente  porque  o  Unix  rodava  em  um  
sistema  com  uma  única  CPU.  Como  o  xv6  roda  em  multiprocessadores,  ele  adiciona  um  bloqueio  explícito  à  
função  sleep.  O  msleep  do  FreeBSD  adota  a  mesma  abordagem.  A  função  sleep  do  Plan  9  usa  uma  função  de  
retorno  de  chamada  que  é  executada  com  o  bloqueio  de  agendamento  mantido  pouco  antes  de  entrar  em  modo  
sleep;  a  função  serve  como  uma  verificação  de  última  hora  da  condição  de  sleep,  para  evitar  perdas  de  ativação.  O  kernel  Linux
 82
Machine Translated by Google
 sleep  usa  uma  fila  de  processos  explícita,  chamada  fila  de  espera,  em  vez  de  um  canal  de  espera;  a  fila  tem  seu  
próprio  bloqueio  interno.
 Varrer  todo  o  conjunto  de  processos  durante  o  processo  de  despertar  é  ineficiente.  Uma  solução  melhor  é  substituir  o  canal  
tanto  no  modo  de  suspensão  quanto  no  de  despertar  por  uma  estrutura  de  dados  que  contenha  uma  lista  de  processos  em  
suspensão  nessa  estrutura,  como  a  fila  de  espera  do  Linux.  O  modo  de  suspensão  e  o  modo  de  despertar  do  Plan  9  chamam  essa  
estrutura  de  ponto  de  encontro.  Muitas  bibliotecas  de  threads  se  referem  à  mesma  estrutura  como  uma  variável  de  condição;  
nesse  contexto,  as  operações  de  suspensão  e  despertar  são  chamadas  de  espera  e  sinalização.  Todos  esses  mecanismos  
compartilham  o  mesmo  princípio:  a  condição  de  suspensão  é  protegida  por  algum  tipo  de  bloqueio  que  é  descartado  atomicamente  
durante  o  processo  de  suspensão.
 A  implementação  do  wakeup  ativa  todos  os  processos  que  aguardam  em  um  canal  específico,  e  pode  ser  que  muitos  
processos  estejam  aguardando  esse  canal  específico.  O  sistema  operacional  agendará  todos  esses  processos  e  eles  verificarão  
rapidamente  a  condição  de  suspensão.
 Processos  que  se  comportam  dessa  maneira  são  às  vezes  chamados  de  "manada  estrondosa"  e  é  melhor  evitá-los.
 A  maioria  das  variáveis  de  condição  tem  duas  primitivas  para  ativação:  sinal,  que  ativa  um  processo,  e  transmissão,  
que  ativa  todos  os  processos  em  espera.
 Semáforos  são  frequentemente  usados  para  sincronização.  A  contagem  normalmente  corresponde  a  algo  como  o  número  de  
bytes  disponíveis  em  um  buffer  de  pipe  ou  o  número  de  filhos  zumbis  que  um  processo  possui.  Usar  uma  contagem  explícita  como  
parte  da  abstração  evita  o  problema  do  "despertar  perdido":  há  uma  contagem  explícita  do  número  de  despertares  ocorridos.  A  
contagem  também  evita  os  problemas  de  despertares  espúrios  e  de  manada  estrondosa.
 Terminar  processos  e  limpá-los  introduz  muita  complexidade  no  xv6.  Na  maioria  dos  sistemas  operacionais,  é  ainda  
mais  complexo,  porque,  por  exemplo,  o  processo  vítima  pode  estar  dormindo  profundamente  dentro  do  kernel,  e  desenrolar  
sua  pilha  requer  cuidado,  já  que  cada  função  na  pilha  de  chamadas  pode  precisar  fazer  alguma  limpeza.  Algumas  
linguagens  ajudam  fornecendo  um  mecanismo  de  exceção,  mas  não  C.  Além  disso,  existem  outros  eventos  que  podem  
fazer  com  que  um  processo  em  espera  seja  despertado,  mesmo  que  o  evento  que  ele  está  aguardando  ainda  não  tenha  
ocorrido.  Por  exemplo,  quando  um  processo  Unix  está  em  espera,  outro  processo  pode  enviar  um  sinal  para  ele.  Nesse  
caso,  o  processo  retornará  da  chamada  de  sistema  interrompida  com  o  valor  -1  e  com  o  código  de  erro  definido  como  
EINTR.  O  aplicativo  pode  verificar  esses  valores  e  decidir  o  que  fazer.  O  Xv6  não  suporta  sinais  e  essa  complexidade  não  
surge.
 O  suporte  do  Xv6  para  kill  não  é  totalmente  satisfatório:  existem  loops  de  sleep  que  provavelmente  deveriam  verificar  se  há  p
>killed.  Um  problema  relacionado  é  que,  mesmo  para  loops  de  sleep  que  verificam  p->killed,  há  uma  corrida  entre  sleep  e  kill;  este  
último  pode  ativar  p->killed  e  tentar  acordar  a  vítima  logo  após  o  loop  da  vítima  verificar  p->killed,  mas  antes  de  chamar  sleep.  Se  
esse  problema  ocorrer,  a  vítima  não  notará  o  p->killed  até  que  a  condição  que  está  aguardando  ocorra.  Isso  pode  ocorrer  um  
pouco  mais  tarde  ou  até  mesmo  nunca  (por  exemplo,  se  a  vítima  estiver  aguardando  uma  entrada  do  console,  mas  o  usuário  não  
digitar  nenhuma  entrada).
 Um  sistema  operacional  real  encontraria  estruturas  proc  livres  com  uma  lista  livre  explícita  em  constante
 tempo  em  vez  da  pesquisa  de  tempo  linear  em  allocproc;  xv6  usa  a  varredura  linear  para  simplificar.

Capítulo  2
 Organização  do  sistema  operacional
 Um  requisito  fundamental  para  um  sistema  operacional  é  suportar  diversas  atividades  simultaneamente.  Por  exemplo,  
usando  a  interface  de  chamada  de  sistema  descrita  no  Capítulo  1,  um  processo  pode  iniciar  novos  processos  com  fork.
 O  sistema  operacional  deve  compartilhar  o  tempo  dos  recursos  do  computador  entre  esses  processos.  Por  exemplo,  
mesmo  que  haja  mais  processos  do  que  CPUs  de  hardware,  o  sistema  operacional  deve  garantir  que  todos  os  processos  
tenham  a  chance  de  executar.  O  sistema  operacional  também  deve  providenciar  o  isolamento  entre  os  processos.  Ou  
seja,  se  um  processo  tiver  um  bug  e  apresentar  mau  funcionamento,  isso  não  deve  afetar  os  processos  que  não  dependem  
do  processo  com  bug.  O  isolamento  completo,  no  entanto,  é  muito  forte,  pois  deve  ser  possível  que  os  processos  interajam  
intencionalmente;  pipelines  são  um  exemplo.  Portanto,  um  sistema  operacional  deve  atender  a  três  requisitos:  
multiplexação,  isolamento  e  interação.
 Este  capítulo  fornece  uma  visão  geral  de  como  os  sistemas  operacionais  são  organizados  para  atender  a  esses  
três  requisitos.  Existem  muitas  maneiras  de  fazer  isso,  mas  este  texto  se  concentra  em  projetos  tradicionais  centrados  
em  um  kernel  monolítico,  usado  por  muitos  sistemas  operacionais  Unix.  Este  capítulo  também  fornece  uma  visão  geral  
de  um  processo  xv6,  que  é  a  unidade  de  isolamento  no  xv6,  e  da  criação  do  primeiro  processo  quando  o  xv6  é  iniciado.
 O  Xv6  é  executado  em  um  microprocessador  RISC-V  multi-core1 ,  e  grande  parte  de  sua  funcionalidade  de  baixo  
nível  (por  exemplo,  sua  implementação  de  processo)  é  específica  para  RISC-V.  RISC-V  é  uma  CPU  de  64  bits,  e  o  xv6  
é  escrito  em  C  "LP64",  o  que  significa  que  long  (L)  e  ponteiros  (P)  na  linguagem  de  programação  C  são  de  64  bits,  mas  
int  é  de  32  bits.  Este  livro  pressupõe  que  o  leitor  tenha  feito  um  pouco  de  programação  em  nível  de  máquina  em  
alguma  arquitetura,  e  apresentará  ideias  específicas  do  RISC-V  à  medida  que  surgirem.  Uma  referência  útil  para  RISC
V  é  "The  RISC-V  Reader:  An  Open  Architecture  Atlas"  [15].  O  ISA  em  nível  de  usuário  [2]  e  a  arquitetura  privilegiada  
[3]  são  as  especificações  oficiais.
 A  CPU  de  um  computador  completo  é  cercada  por  hardware  de  suporte,  grande  parte  dele  na  forma  de  interfaces  
de  E/S.  O  Xv6  foi  escrito  para  o  hardware  de  suporte  simulado  pela  opção  "-machine  virt"  do  qemu.  Isso  inclui  RAM,  
uma  ROM  contendo  o  código  de  inicialização,  uma  conexão  serial  com  o  teclado/tela  do  usuário  e  um  disco  para  
armazenamento.
 1Por  "multinúcleo",  este  texto  se  refere  a  múltiplas  CPUs  que  compartilham  memória,  mas  executam  em  paralelo,  cada  uma  com  seu  próprio  
conjunto  de  registradores.  Este  texto  às  vezes  usa  o  termo  multiprocessador  como  sinônimo  de  multinúcleo,  embora  multiprocessador  também  
possa  se  referir  mais  especificamente  a  um  computador  com  vários  chips  de  processador  distintos.
 21
Machine Translated by Google
 2.1  Abstraindo  recursos  físicos
 A  primeira  pergunta  que  se  pode  fazer  ao  se  deparar  com  um  sistema  operacional  é:  por  que  tê-lo?  Ou  seja,  seria  possível  
implementar  as  chamadas  de  sistema  da  Figura  1.2  como  uma  biblioteca,  com  a  qual  os  aplicativos  se  conectam.  Nesse  plano,  
cada  aplicativo  poderia  até  ter  sua  própria  biblioteca,  adaptada  às  suas  necessidades.  Os  aplicativos  poderiam  interagir  
diretamente  com  os  recursos  de  hardware  e  utilizá-los  da  melhor  maneira  possível  para  o  aplicativo  (por  exemplo,  para  atingir  um  
desempenho  alto  ou  previsível).  Alguns  sistemas  operacionais  para  dispositivos  embarcados  ou  sistemas  de  tempo  real  são  
organizados  dessa  maneira.
 A  desvantagem  dessa  abordagem  de  biblioteca  é  que,  se  houver  mais  de  um  aplicativo  em  execução,  eles  precisam  se  
comportar  bem.  Por  exemplo,  cada  aplicativo  precisa  liberar  a  CPU  periodicamente  para  que  outros  aplicativos  possam  ser  
executados.  Esse  esquema  cooperativo  de  compartilhamento  de  tempo  pode  ser  aceitável  se  todos  os  aplicativos  confiarem  uns  
nos  outros  e  não  apresentarem  bugs.  É  mais  comum  que  aplicativos  não  confiem  uns  nos  outros  e  apresentem  bugs,  portanto,  
muitas  vezes,  deseja-se  um  isolamento  mais  forte  do  que  o  proporcionado  por  um  esquema  cooperativo.
 Para  obter  um  isolamento  forte,  é  útil  proibir  que  aplicativos  acessem  diretamente  recursos  de  hardware  sensíveis  e,  
em  vez  disso,  abstrair  os  recursos  em  serviços.  Por  exemplo,  aplicativos  Unix  interagem  com  o  armazenamento  apenas  
por  meio  das  chamadas  de  sistema  de  abertura,  leitura,  gravação  e  fechamento  do  sistema  de  arquivos,  em  vez  de  ler  e  
gravar  diretamente  no  disco.  Isso  proporciona  ao  aplicativo  a  conveniência  de  nomes  de  caminho  e  permite  que  o  
sistema  operacional  (como  implementador  da  interface)  gerencie  o  disco.  Mesmo  que  o  isolamento  não  seja  uma  
preocupação,  programas  que  interagem  intencionalmente  (ou  apenas  desejam  se  manter  fora  do  caminho  uns  dos  outros)  
provavelmente  considerarão  um  sistema  de  arquivos  uma  abstração  mais  conveniente  do  que  o  uso  direto  do  disco.
 Da  mesma  forma,  o  Unix  alterna  de  forma  transparente  as  CPUs  de  hardware  entre  os  processos,  salvando  e  restaurando  
o  estado  dos  registradores  conforme  necessário,  para  que  os  aplicativos  não  precisem  se  preocupar  com  o  compartilhamento  de  
tempo.  Essa  transparência  permite  que  o  sistema  operacional  compartilhe  CPUs  mesmo  que  alguns  aplicativos  estejam  em  loops  
infinitos.
 Como  outro  exemplo,  os  processos  Unix  usam  o  comando  exec  para  construir  sua  imagem  de  memória,  em  vez  de  interagir  
diretamente  com  a  memória  física.  Isso  permite  que  o  sistema  operacional  decida  onde  colocar  um  processo  na  memória;  se  a  
memória  estiver  limitada,  o  sistema  operacional  pode  até  armazenar  alguns  dados  de  um  processo  em  disco.  O  comando  exec  
também  oferece  aos  usuários  a  conveniência  de  um  sistema  de  arquivos  para  armazenar  imagens  de  programas  executáveis.
 Muitas  formas  de  interação  entre  processos  Unix  ocorrem  por  meio  de  descritores  de  arquivo.  Os  descritores  de  arquivo  não  
apenas  abstraem  muitos  detalhes  (por  exemplo,  onde  os  dados  em  um  pipe  ou  arquivo  são  armazenados),  como  também  são  
definidos  de  forma  a  simplificar  a  interação.  Por  exemplo,  se  uma  aplicação  em  um  pipeline  falhar,  o  kernel  gera  um  sinal  de  fim  
de  arquivo  para  o  próximo  processo  no  pipeline.
 A  interface  de  chamada  de  sistema  na  Figura  1.2  foi  cuidadosamente  projetada  para  proporcionar  conveniência  ao  
programador  e  a  possibilidade  de  forte  isolamento.  A  interface  Unix  não  é  a  única  maneira  de  abstrair  recursos,  mas  provou  ser  
uma  ótima  maneira.
 2.2  Modo  de  usuário,  modo  supervisor  e  chamadas  de  sistema
 O  isolamento  forte  requer  uma  fronteira  rígida  entre  os  aplicativos  e  o  sistema  operacional.  Se  o  aplicativo  cometer  um  erro,  não  
queremos  que  o  sistema  operacional  falhe  ou  que  outros  aplicativos  sejam  afetados.
 22
Machine Translated by Google
 falha.  Em  vez  disso,  o  sistema  operacional  deve  ser  capaz  de  limpar  o  aplicativo  com  falha  e  continuar  executando  
outros  aplicativos.  Para  alcançar  um  isolamento  forte,  o  sistema  operacional  deve  garantir  que  os  aplicativos  não  
possam  modificar  (ou  mesmo  ler)  suas  estruturas  de  dados  e  instruções,  e  que  os  aplicativos  não  possam  acessar  a  
memória  de  outros  processos.
 As  CPUs  fornecem  suporte  de  hardware  para  isolamento  robusto.  Por  exemplo,  o  RISC-V  possui  três  modos  nos  
quais  a  CPU  pode  executar  instruções:  modo  máquina,  modo  supervisor  e  modo  usuário.  Instruções  de  operação  
executadas  em  modo  máquina  têm  privilégio  total;  uma  CPU  inicia  em  modo  máquina.  O  modo  máquina  destina-se  
principalmente  à  configuração  de  um  computador.  O  Xv6  executa  algumas  linhas  em  modo  máquina  e  depois  muda  
para  o  modo  supervisor.
 No  modo  supervisor,  a  CPU  tem  permissão  para  executar  instruções  privilegiadas:  por  exemplo,  habilitar  e  
desabilitar  interrupções,  ler  e  escrever  no  registrador  que  contém  o  endereço  de  uma  tabela  de  páginas,  etc.  Se  um  
aplicativo  no  modo  usuário  tentar  executar  uma  instrução  privilegiada,  a  CPU  não  executa  a  instrução,  mas  alterna  
para  o  modo  supervisor  para  que  o  código  do  modo  supervisor  possa  encerrar  o  aplicativo,  porque  ele  fez  algo  que  não  
deveria.  A  Figura  1.1  no  Capítulo  1  ilustra  essa  organização.  Um  aplicativo  pode  executar  apenas  instruções  do  modo  
usuário  (por  exemplo,  somar  números,  etc.)  e  é  dito  estar  em  execução  no  espaço  do  usuário,  enquanto  o  software  no  
modo  supervisor  também  pode  executar  instruções  privilegiadas  e  é  dito  estar  em  execução  no  espaço  do  kernel.  O  
software  em  execução  no  espaço  do  kernel  (ou  no  modo  supervisor)  é  chamado  de  kernel.
 Um  aplicativo  que  deseja  invocar  uma  função  do  kernel  (por  exemplo,  a  chamada  de  sistema  read  no  xv6)  deve  
fazer  a  transição  para  o  kernel;  um  aplicativo  não  pode  invocar  uma  função  do  kernel  diretamente.  As  CPUs  fornecem  
uma  instrução  especial  que  alterna  a  CPU  do  modo  usuário  para  o  modo  supervisor  e  entra  no  kernel  em  um  ponto  de  
entrada  especificado  pelo  kernel.  (O  RISC-V  fornece  a  instrução  ecall  para  esse  propósito.)
 Após  a  CPU  alternar  para  o  modo  supervisor,  o  kernel  pode  validar  os  argumentos  da  chamada  de  sistema  (por  
exemplo,  verificar  se  o  endereço  passado  para  a  chamada  de  sistema  faz  parte  da  memória  do  aplicativo),  decidir  se  o  
aplicativo  tem  permissão  para  executar  a  operação  solicitada  (por  exemplo,  verificar  se  o  aplicativo  tem  permissão  para  
gravar  o  arquivo  especificado)  e,  em  seguida,  negá-la  ou  executá-la.  É  importante  que  o  kernel  controle  o  ponto  de  
entrada  para  transições  para  o  modo  supervisor;  se  o  aplicativo  pudesse  decidir  o  ponto  de  entrada  do  kernel,  um  
aplicativo  malicioso  poderia,  por  exemplo,  entrar  no  kernel  em  um  ponto  onde  a  validação  dos  argumentos  é  ignorada.
 2.3  Organização  do  kernel
 Uma  questão  fundamental  de  projeto  é  qual  parte  do  sistema  operacional  deve  ser  executada  no  modo  supervisor.  
Uma  possibilidade  é  que  todo  o  sistema  operacional  resida  no  kernel,  de  modo  que  as  implementações  de  todas  as  
chamadas  de  sistema  sejam  executadas  no  modo  supervisor.  Essa  organização  é  chamada  de  kernel  monolítico.
 Nesta  organização,  todo  o  sistema  operacional  é  executado  com  privilégios  totais  de  hardware.  Essa  organização  
é  conveniente  porque  o  projetista  do  sistema  operacional  não  precisa  decidir  qual  parte  do  sistema  operacional  não  
precisa  de  privilégios  totais  de  hardware.  Além  disso,  facilita  a  cooperação  entre  diferentes  partes  do  sistema  
operacional.  Por  exemplo,  um  sistema  operacional  pode  ter  um  cache  de  buffer  que  pode  ser  compartilhado  tanto  pelo  
sistema  de  arquivos  quanto  pelo  sistema  de  memória  virtual.
 Uma  desvantagem  da  organização  monolítica  é  que  as  interfaces  entre  as  diferentes  partes  do  sistema  operacional  
são  frequentemente  complexas  (como  veremos  no  restante  deste  texto)  e,  portanto,  é
 23
Machine Translated by Google
 usuário
 espaço
 núcleo
 espaço
 concha
 Servidor  de  arquivos
 Enviar  mensagem
 Microkernel
 Figura  2.1:  Um  microkernel  com  um  servidor  de  sistema  de  arquivos
 É  fácil  para  um  desenvolvedor  de  sistema  operacional  cometer  um  erro.  Em  um  kernel  monolítico,  um  erro  é  fatal,  pois  um  erro  
no  modo  supervisor  frequentemente  causa  falha  no  kernel.  Se  o  kernel  falhar,  o  computador  para  de  funcionar  e,  portanto,  todos  
os  aplicativos  também  falham.  O  computador  precisa  ser  reiniciado  para  iniciar  novamente.
 Para  reduzir  o  risco  de  erros  no  kernel,  os  projetistas  de  sistemas  operacionais  podem  minimizar  a  quantidade  de  código  do  
sistema  operacional  executado  no  modo  supervisor  e  executar  a  maior  parte  do  sistema  operacional  no  modo  usuário.  Essa  
organização  do  kernel  é  chamada  de  microkernel.
 A  Figura  2.1  ilustra  este  projeto  de  microkernel.  Na  figura,  o  sistema  de  arquivos  é  executado  como  um  processo  em  nível  
de  usuário.  Os  serviços  do  sistema  operacional  executados  como  processos  são  chamados  de  servidores.  Para  permitir  que  os  
aplicativos  interajam  com  o  servidor  de  arquivos,  o  kernel  fornece  um  mecanismo  de  comunicação  entre  processos  para  enviar  
mensagens  de  um  processo  em  modo  de  usuário  para  outro.  Por  exemplo,  se  um  aplicativo,  como  o  shell,  deseja  ler  ou  gravar  
um  arquivo,  ele  envia  uma  mensagem  ao  servidor  de  arquivos  e  aguarda  uma  resposta.
 Em  um  microkernel,  a  interface  do  kernel  consiste  em  algumas  funções  de  baixo  nível  para  iniciar  aplicativos,  enviar  
mensagens,  acessar  hardware  do  dispositivo,  etc.  Essa  organização  permite  que  o  kernel  seja  relativamente  simples,  já  que  a  
maior  parte  do  sistema  operacional  reside  em  servidores  de  nível  de  usuário.
 No  mundo  real,  tanto  kernels  monolíticos  quanto  microkernels  são  populares.  Muitos  kernels  Unix  são  monolíticos.  Por  
exemplo,  o  Linux  possui  um  kernel  monolítico,  embora  algumas  funções  do  sistema  operacional  sejam  executadas  como  
servidores  em  nível  de  usuário  (por  exemplo,  o  sistema  de  janelas).  O  Linux  oferece  alto  desempenho  para  aplicativos  com  uso  
intensivo  de  sistema  operacional,  em  parte  porque  os  subsistemas  do  kernel  podem  ser  totalmente  integrados.
 Sistemas  operacionais  como  Minix,  L4  e  QNX  são  organizados  como  um  microkernel  com  servidores  e  têm  sido  amplamente  
implantados  em  ambientes  embarcados.  Uma  variante  do  L4,  seL4,  é  pequena  o  suficiente  para  ter  sua  segurança  de  memória  e  
outras  propriedades  de  segurança  verificadas  [8].
 Há  muito  debate  entre  desenvolvedores  de  sistemas  operacionais  sobre  qual  organização  é  melhor,  e  não  há  evidências  
conclusivas  de  um  lado  ou  de  outro.  Além  disso,  depende  muito  do  que  "melhor"  significa:  desempenho  mais  rápido,  código  
menor,  confiabilidade  do  kernel,  confiabilidade  de  todo  o  sistema  operacional  (incluindo  serviços  em  nível  de  usuário),  etc.
 Há  também  considerações  práticas  que  podem  ser  mais  importantes  do  que  a  questão  de  qual  organização.  Alguns  sistemas  
operacionais  possuem  um  microkernel,  mas  executam  alguns  dos  serviços  de  nível  de  usuário  no  espaço  do  kernel  por  questões  
de  desempenho.  Alguns  sistemas  operacionais  possuem  kernels  monolíticos  porque  foi  assim  que  começaram,  e  há  pouco  
incentivo  para  migrar  para  uma  organização  puramente  de  microkernel,  pois  novos  recursos  podem  ser  mais  importantes  do  que  
reescrever  o  sistema  operacional  existente  para  se  adequar  a  um  design  de  microkernel.
 Da  perspectiva  deste  livro,  microkernel  e  sistemas  operacionais  monolíticos  compartilham  muitas  ideias-chave.  Eles  
implementam  chamadas  de  sistema,  usam  tabelas  de  páginas,  lidam  com  interrupções  e  oferecem  suporte
 24
Machine Translated by Google
 Arquivo
 bio.c  
console.c  
entrada.S
 exec.c
 arquivo.c  fs.c
 Descrição
 Cache  de  bloco  de  disco  para  o  sistema  de  arquivos.
 Conecte-se  ao  teclado  e  à  tela  do  usuário.
 Primeiras  instruções  de  inicialização.
 chamada  de  sistema  exec().
 Suporte  ao  descritor  de  arquivo.
 Sistema  de  arquivos.
 kalloc.c  Alocador  de  páginas  físicas.
 kernelvec.S  Lida  com  armadilhas  do  kernel  e  interrupções  de  temporizador.
 log.c  Registro  do  sistema  de  arquivos  e  recuperação  de  falhas.
 main.c  Controla  a  inicialização  de  outros  módulos  durante  a  inicialização.
 pipe.c  Tubos.
 plic.c  Controlador  de  interrupção  RISC-V.
 printf.c  Saída  formatada  para  o  console.
 Processos  e  agendamento.
 proc.c  
sleeplock.c  Bloqueios  que  produzem  a  CPU.
 spinlock.c  Bloqueios  que  não  cedem  a  CPU.
 start.c  
string.c  
swtch.S  
syscall.c  
sysfile.c
 Código  de  inicialização  inicial  do  modo  máquina.
 Biblioteca  de  strings  e  matrizes  de  bytes  em  C.
 Troca  de  thread.
 Despachar  chamadas  do  sistema  para  a  função  de  manuseio.
 Chamadas  de  sistema  relacionadas  a  arquivos.
 Chamadas  de  sistema  relacionadas  ao  processo.
 sysproc.c  
trampoline.S  Código  de  montagem  para  alternar  entre  usuário  e  kernel.
 trap.c  
uart.c
 Código  C  para  manipular  e  retornar  armadilhas  e  interrupções.
 Driver  de  dispositivo  de  console  de  porta  serial.
 virtio_disk.c  Driver  de  dispositivo  de  disco.
 vm.c
 Gerenciar  tabelas  de  páginas  e  espaços  de  endereço.
 Figura  2.2:  Arquivos  de  origem  do  kernel  Xv6.
 processos,  eles  usam  bloqueios  para  controle  de  simultaneidade,  eles  implementam  um  sistema  de  arquivos,  etc.  Este  livro
 concentra-se  nessas  ideias  principais.
 O  Xv6  é  implementado  como  um  kernel  monolítico,  como  a  maioria  dos  sistemas  operacionais  Unix.  Assim,  o  xv6
 A  interface  do  kernel  corresponde  à  interface  do  sistema  operacional,  e  o  kernel  implementa  o  sistema  operacional  completo.  Como  o  xv6  
não  fornece  muitos  serviços,  seu  kernel  é  menor  do  que  alguns
 microkernels,  mas  conceitualmente  o  xv6  é  monolítico.
 2.4  Código:  organização  xv6
 O  código-fonte  do  kernel  xv6  está  no  subdiretório  kernel/.  O  código-fonte  é  dividido  em  arquivos,  a  seguir
 uma  noção  aproximada  de  modularidade;  a  Figura  2.2  lista  os  arquivos.  As  interfaces  entre  módulos  são  definidas  em
 25
Machine Translated by Google
 defs.h  (kernel/defs.h).
 2.5  Visão  geral  do  processo
 MAXVA
 0
 estrutura  de  
armadilha  de  trampolim
 monte
 pilha  de  usuário
 texto  do  usuário
 e  dados
 Figura  2.3:  Layout  do  espaço  de  endereço  virtual  de  um  processo
 A  unidade  de  isolamento  no  xv6  (assim  como  em  outros  sistemas  operacionais  Unix)  é  um  processo.  A  abstração  de  
processo  impede  que  um  processo  destrua  ou  espione  a  memória,  a  CPU,  os  descritores  de  arquivo,  etc.  de  outro  processo.  
Também  impede  que  um  processo  destrua  o  próprio  kernel,  impedindo  que  um  processo  subverta  os  mecanismos  de  
isolamento  do  kernel.  O  kernel  deve  implementar  a  abstração  de  processo  com  cuidado,  pois  um  aplicativo  malicioso  ou  
com  bugs  pode  induzir  o  kernel  ou  o  hardware  a  fazer  algo  errado  (por  exemplo,  burlar  o  isolamento).  Os  mecanismos  
usados  pelo  kernel  para  implementar  processos  incluem  o  sinalizador  de  modo  usuário/supervisor,  espaços  de  endereço  e  
fatiamento  de  tempo  de  threads.
 Para  ajudar  a  reforçar  o  isolamento,  a  abstração  do  processo  dá  a  um  programa  a  ilusão  de  que  ele  possui  sua  própria  
máquina  privada.  Um  processo  fornece  ao  programa  o  que  parece  ser  um  sistema  de  memória  privado,  ou  espaço  de  
endereço,  que  outros  processos  não  conseguem  ler  ou  escrever.  Um  processo  também  fornece  ao  programa  o  que  parece  
ser  sua  própria  CPU  para  executar  as  instruções  do  programa.
 O  Xv6  utiliza  tabelas  de  páginas  (implementadas  por  hardware)  para  dar  a  cada  processo  seu  próprio  espaço  de  
endereço.  A  tabela  de  páginas  RISC-V  traduz  (ou  "mapeia")  um  endereço  virtual  (o  endereço  que  uma  instrução  RISC-V  
manipula)  para  um  endereço  físico  (um  endereço  que  o  chip  da  CPU  envia  para  a  memória  principal).
 O  Xv6  mantém  uma  tabela  de  páginas  separada  para  cada  processo  que  define  o  espaço  de  endereço  desse  processo.
 Conforme  ilustrado  na  Figura  2.3,  um  espaço  de  endereço  inclui  a  memória  do  usuário  do  processo,  começando  no  endereço  
virtual  zero.  As  instruções  vêm  primeiro,  seguidas  pelas  variáveis  globais,  depois  pela  pilha  e,  finalmente,  por  uma  área  de  
"heap"  (para  malloc)  que  o  processo  pode  expandir  conforme  necessário.  Há  uma  série  de  fatores  que  limitam  o  tamanho  
máximo  do  espaço  de  endereço  de  um  processo:  ponteiros  no  RISC-V  têm  64  bits  de  largura;  o  hardware  usa  apenas  os  39  
bits  mais  baixos  ao  procurar  endereços  virtuais  em  tabelas  de  páginas;  e  o  xv6  usa  apenas  38  desses  39  bits.  Portanto,  o  
endereço  máximo  é  2  ÿ  1  =  0x3ffffffffff,  que  é
 38
 26
Machine Translated by Google
 MAXVA  (kernel/riscv.h:363).  No  topo  do  espaço  de  endereço,  o  xv6  reserva  uma  página  para  um  trampolim  e  uma  
página  que  mapeia  o  trapframe  do  processo.  O  xv6  usa  essas  duas  páginas  para  fazer  a  transição  de  entrada  e  
saída  do  kernel;  a  página  do  trampolim  contém  o  código  para  a  transição  de  entrada  e  saída  do  kernel,  e  o  
mapeamento  do  trapframe  é  necessário  para  salvar/restaurar  o  estado  do  processo  do  usuário,  como  explicaremos  no  Capítulo  4.
 O  kernel  xv6  mantém  muitas  partes  de  estado  para  cada  processo,  que  ele  reúne  em  uma  estrutura  proc
 (kernel/proc.h:85).  As  partes  mais  importantes  do  estado  do  kernel  de  um  processo  são  sua  tabela  de  páginas,  sua  pilha  
do  kernel  e  seu  estado  de  execução.  Usaremos  a  notação  p->xxx  para  nos  referirmos  aos  elementos  da  estrutura  proc ;  
por  exemplo,  p->pagetable  é  um  ponteiro  para  a  tabela  de  páginas  do  processo.
 Cada  processo  possui  uma  thread  de  execução  (ou  thread,  para  abreviar)  que  executa  as  instruções  do  processo.  
Uma  thread  pode  ser  suspensa  e  posteriormente  retomada.  Para  alternar  entre  processos  de  forma  transparente,  o  kernel  
suspende  a  thread  em  execução  e  retoma  a  thread  de  outro  processo.  Grande  parte  do  estado  de  uma  thread  (variáveis  
locais,  endereços  de  retorno  de  chamadas  de  função)  é  armazenada  nas  pilhas  da  thread.
 Cada  processo  possui  duas  pilhas:  uma  pilha  do  usuário  e  uma  pilha  do  kernel  (p->kstack).  Quando  o  processo  está  
executando  instruções  do  usuário,  apenas  sua  pilha  do  usuário  está  em  uso,  e  sua  pilha  do  kernel  está  vazia.  Quando  o  
processo  entra  no  kernel  (para  uma  chamada  de  sistema  ou  interrupção),  o  código  do  kernel  é  executado  na  pilha  do  kernel  
do  processo;  enquanto  um  processo  está  no  kernel,  sua  pilha  do  usuário  ainda  contém  dados  salvos,  mas  não  é  usada  
ativamente.  A  thread  de  um  processo  alterna  entre  usar  ativamente  sua  pilha  do  usuário  e  sua  pilha  do  kernel.
 A  pilha  do  kernel  é  separada  (e  protegida  do  código  do  usuário)  para  que  o  kernel  possa  ser  executado  mesmo  que  um  
processo  tenha  destruído  sua  pilha  do  usuário.
 Um  processo  pode  fazer  uma  chamada  de  sistema  executando  a  instrução  ecall  do  RISC-V .  Essa  instrução  aumenta  
o  nível  de  privilégio  de  hardware  e  altera  o  contador  de  programa  para  um  ponto  de  entrada  definido  pelo  kernel.
 O  código  no  ponto  de  entrada  alterna  para  uma  pilha  do  kernel  e  executa  as  instruções  do  kernel  que  implementam  a  
chamada  do  sistema.  Quando  a  chamada  do  sistema  é  concluída,  o  kernel  retorna  para  a  pilha  do  usuário  e  retorna  ao  
espaço  do  usuário  chamando  a  instrução  sret ,  que  reduz  o  nível  de  privilégio  de  hardware  e  retoma  a  execução  das  
instruções  do  usuário  logo  após  a  instrução  de  chamada  do  sistema.  A  thread  de  um  processo  pode  "bloquear"  no  kernel  
para  aguardar  a  E/S  e  retomar  de  onde  parou  quando  a  E/S  terminar.
 p->state  indica  se  o  processo  está  alocado,  pronto  para  execução,  em  execução,  aguardando  E/S  ou  encerrando.  
p
>pagetable  contém  a  tabela  de  páginas  do  processo,  no  formato  esperado  pelo  hardware  RISC-V.  O  Xv6  faz  com  
que  o  hardware  de  paginação  utilize  a  p->pagetable  de  um  processo  ao  executá-lo  no  espaço  do  usuário.  A  tabela  
de  páginas  de  um  processo  também  serve  como  registro  dos  endereços  das  páginas  físicas  alocadas  para  
armazenar  a  memória  do  processo.
 Em  resumo,  um  processo  agrupa  duas  ideias  de  design:  um  espaço  de  endereço  para  dar  ao  processo  a  ilusão  de  ter  
sua  própria  memória  e  uma  thread  para  dar  ao  processo  a  ilusão  de  ter  sua  própria  CPU.  No  xv6,  um  processo  consiste  
em  um  espaço  de  endereço  e  uma  thread.  Em  sistemas  operacionais  reais,  um  processo  pode  ter  mais  de  uma  thread  
para  aproveitar  múltiplas  CPUs.
 2.6  Código:  iniciando  o  xv6,  o  primeiro  processo  e  chamada  do  sistema
 Para  tornar  o  xv6  mais  concreto,  descreveremos  como  o  kernel  inicia  e  executa  o  primeiro  processo.  Os  capítulos  
subsequentes  descreverão  os  mecanismos  que  aparecem  nesta  visão  geral  com  mais  detalhes.
 27
Machine Translated by Google
 Quando  o  computador  RISC-V  é  ligado,  ele  se  inicializa  e  executa  um  carregador  de  boot  armazenado  na  memória  
somente  leitura.  O  carregador  de  boot  carrega  o  kernel  xv6  na  memória.  Em  seguida,  no  modo  máquina,  a  CPU  executa  
o  xv6  a  partir  de  _entry  (kernel/entry.S:7).  O  RISC-V  começa  com  o  hardware  de  paginação  desabilitado:  endereços  
virtuais  são  mapeados  diretamente  para  endereços  físicos.
 O  carregador  carrega  o  kernel  xv6  na  memória  no  endereço  físico  0x80000000.  O  motivo  pelo  qual  ele  coloca  o  
kernel  em  0x80000000  em  vez  de  0x0  é  porque  o  intervalo  de  endereços  0x0:0x80000000  contém  dispositivos  de  E/S.
 As  instruções  em  _entry  configuram  uma  pilha  para  que  o  xv6  possa  executar  código  C.  O  xv6  declara  espaço  para  
uma  pilha  inicial,  stack0,  no  arquivo  start.c  (kernel/start.c:11).  O  código  em  _entry  carrega  o  registrador  de  ponteiro  de  
pilha  sp  com  o  endereço  stack0+4096,  o  topo  da  pilha,  porque  a  pilha  no  RISC-V  cresce  para  baixo.  Agora  que  o  kernel  
tem  uma  pilha,  _entry  chama  o  código  C  no  início  (kernel/start.c:21).
 A  função  start  realiza  algumas  configurações  permitidas  apenas  no  modo  máquina  e,  em  seguida,  alterna  para  o  
modo  supervisor.  Para  entrar  no  modo  supervisor,  o  RISC-V  fornece  a  instrução  mret.  Essa  instrução  é  mais  
frequentemente  usada  para  retornar  de  uma  chamada  anterior  do  modo  supervisor  para  o  modo  máquina.  start  não  
retorna  de  tal  chamada  e,  em  vez  disso,  configura  tudo  como  se  houvesse  uma:  define  o  modo  de  privilégio  anterior  como  
supervisor  no  registrador  mstatus,  define  o  endereço  de  retorno  como  main  escrevendo  o  endereço  de  main  no  registrador  
mepc,  desabilita  a  tradução  de  endereços  virtuais  no  modo  supervisor  escrevendo  0  no  registrador  de  tabela  de  páginas  
satp  e  delega  todas  as  interrupções  e  exceções  ao  modo  supervisor.
 Antes  de  entrar  no  modo  supervisor,  start  executa  mais  uma  tarefa:  programa  o  chip  do  relógio  
para  gerar  interrupções  de  tempo.  Com  essa  tarefa  concluída,  start  "retorna"  ao  modo  supervisor  
chamando  mret.  Isso  faz  com  que  o  contador  de  programa  mude  para  o  modo  principal  (kernel/main.c:11).
 Após  o  principal  (kernel/main.c:11)  inicializa  vários  dispositivos  e  subsistemas,  ele  cria  o  primeiro  processo  
chamando  userinit  (kernel/proc.c:233).  O  primeiro  processo  executa  um  pequeno  programa  escrito  em  assembly  RISC-V,  
que  faz  a  primeira  chamada  de  sistema  no  xv6.  initcode.S  (user/initcode.S:3)  carrega  o  número  para  a  chamada  do  
sistema  exec ,  SYS_EXEC  (kernel/syscall.h:8),  no  registro  a7  e,  em  seguida,  chama  ecall  para  entrar  novamente  no  
kernel.
 O  kernel  usa  o  número  no  registro  a7  em  syscall  (kernel/syscall.c:132)  para  chamar  a  chamada  de  sistema  desejada.  
A  tabela  de  chamadas  de  sistema  (kernel/syscall.c:107)  mapeia  SYS_EXEC  para  sys_exec,  que  o  kernel  invoca.  Como  
vimos  no  Capítulo  1,  exec  substitui  a  memória  e  os  registradores  do  processo  atual  por  um  novo  programa  (neste  caso, /
 init).
 Depois  que  o  kernel  conclui  o  exec,  ele  retorna  ao  espaço  do  usuário  no  processo /init .  init  (user/init.c:15)  Cria  um  
novo  arquivo  de  dispositivo  de  console,  se  necessário,  e  o  abre  como  descritores  de  arquivo  0,  1  e  2.  Em  seguida,  inicia  
um  shell  no  console.  O  sistema  está  funcionando.
 2.7  Modelo  de  Segurança
 Você  pode  se  perguntar  como  o  sistema  operacional  lida  com  códigos  maliciosos  ou  com  bugs.  Como  lidar  com  códigos  
maliciosos  é  mais  difícil  do  que  lidar  com  bugs  acidentais,  é  razoável  considerar  este  tópico  como  relacionado  à  
segurança.  Aqui  está  uma  visão  geral  das  premissas  e  objetivos  típicos  de  segurança  no  design  de  sistemas  operacionais.
 28
Machine Translated by Google
 O  sistema  operacional  deve  assumir  que  o  código  de  nível  de  usuário  de  um  processo  fará  o  possível  para  destruir  o  
kernel  ou  outros  processos.  O  código  do  usuário  pode  tentar  desreferenciar  ponteiros  fora  de  seu  espaço  de  endereço  
permitido;  pode  tentar  executar  quaisquer  instruções  RISC-V,  mesmo  aquelas  não  destinadas  ao  código  do  usuário;  pode  
tentar  ler  e  escrever  qualquer  registrador  de  controle  RISC-V;  pode  tentar  acessar  diretamente  o  hardware  do  dispositivo;  e  
pode  passar  valores  inteligentes  para  chamadas  de  sistema  na  tentativa  de  enganar  o  kernel  para  travar  ou  fazer  algo  
estúpido.  O  objetivo  do  kernel  é  restringir  cada  processo  do  usuário  para  que  tudo  o  que  ele  possa  fazer  seja  ler/escrever/
 executar  sua  própria  memória  de  usuário,  usar  os  32  registradores  RISC-V  de  uso  geral  e  afetar  o  kernel  e  outros  processos  
das  maneiras  que  as  chamadas  de  sistema  devem  permitir.  O  kernel  deve  impedir  quaisquer  outras  ações.  Isso  normalmente  
é  um  requisito  absoluto  no  projeto  do  kernel.
 As  expectativas  para  o  código  do  kernel  são  bem  diferentes.  O  código  do  kernel  é  assumido  como  escrito  por  
programadores  bem-intencionados  e  cuidadosos.  Espera-se  que  o  código  do  kernel  esteja  livre  de  bugs  e  certamente  não  
contenha  nada  malicioso.  Essa  suposição  afeta  a  forma  como  analisamos  o  código  do  kernel.  Por  exemplo,  existem  muitas  
funções  internas  do  kernel  (por  exemplo,  os  bloqueios  de  spin)  que  causariam  sérios  problemas  se  o  código  do  kernel  as  
utilizasse  incorretamente.  Ao  examinar  qualquer  trecho  específico  do  código  do  kernel,  queremos  nos  convencer  de  que  ele  
se  comporta  corretamente.  Assumimos,  no  entanto,  que  o  código  do  kernel  em  geral  está  escrito  corretamente  e  segue  
todas  as  regras  sobre  o  uso  das  funções  e  estruturas  de  dados  do  kernel.  No  nível  de  hardware,  presume-se  que  a  CPU,  a  
RAM,  o  disco  etc.  do  RISC-V  operem  conforme  anunciado  na  documentação,  sem  bugs  de  hardware.
 É  claro  que  na  vida  real  as  coisas  não  são  tão  simples.  É  difícil  evitar  que  um  código  de  usuário  inteligente  torne  um  
sistema  inutilizável  (ou  o  faça  entrar  em  pânico)  consumindo  recursos  protegidos  pelo  kernel  –  espaço  em  disco,  tempo  de  
CPU,  slots  na  tabela  de  processos,  etc.  Geralmente  é  impossível  escrever  código  livre  de  bugs  ou  projetar  hardware  livre  de  
bugs;  se  os  escritores  de  código  de  usuário  malicioso  estiverem  cientes  de  bugs  no  kernel  ou  no  hardware,  eles  os  
explorarão.  Mesmo  em  kernels  maduros  e  amplamente  utilizados,  como  o  Linux,  as  pessoas  descobrem  novas  
vulnerabilidades  continuamente  [1].  Vale  a  pena  projetar  salvaguardas  no  kernel  contra  a  possibilidade  de  ele  ter  bugs:  
asserções,  verificação  de  tipos,  páginas  de  proteção  de  pilha,  etc.  Finalmente,  a  distinção  entre  código  de  usuário  e  kernel  
às  vezes  é  confusa:  alguns  processos  privilegiados  em  nível  de  usuário  podem  fornecer  serviços  essenciais  e  efetivamente  
fazer  parte  do  sistema  operacional  e,  em  alguns  sistemas  operacionais,  o  código  de  usuário  privilegiado  pode  inserir  novo  
código  no  kernel  (como  com  os  módulos  carregáveis  do  kernel  do  Linux).
 2.8  Mundo  real
 A  maioria  dos  sistemas  operacionais  adotou  o  conceito  de  processo,  e  a  maioria  dos  processos  se  assemelha  aos  do  xv6.  
Os  sistemas  operacionais  modernos,  no  entanto,  suportam  várias  threads  dentro  de  um  processo,  permitindo  que  um  único  
processo  explore  múltiplas  CPUs.  Suportar  múltiplas  threads  em  um  processo  envolve  uma  série  de  mecanismos  que  o  xv6  
não  possui,  incluindo  possíveis  mudanças  na  interface  (por  exemplo,  o  clone  do  Linux,  uma  variante  do  fork)  para  controlar  
quais  aspectos  de  um  processo  as  threads  compartilham.

 Capítulo  5
 Interrupções  e  drivers  de  dispositivo
 Um  driver  é  o  código  em  um  sistema  operacional  que  gerencia  um  dispositivo  específico:  ele  configura  o  hardware  do  
dispositivo,  instrui  o  dispositivo  a  executar  operações,  lida  com  as  interrupções  resultantes  e  interage  com  processos  
que  podem  estar  aguardando  E/S  do  dispositivo.  O  código  do  driver  pode  ser  complicado,  pois  ele  é  executado  
simultaneamente  com  o  dispositivo  que  gerencia.  Além  disso,  o  driver  precisa  entender  a  interface  de  hardware  do  
dispositivo,  que  pode  ser  complexa  e  mal  documentada.
 Dispositivos  que  precisam  da  atenção  do  sistema  operacional  geralmente  podem  ser  configurados  para  gerar  
interrupções,  que  são  um  tipo  de  armadilha.  O  código  de  tratamento  de  armadilhas  do  kernel  reconhece  quando  
um  dispositivo  gera  uma  interrupção  e  chama  o  manipulador  de  interrupções  do  driver;  no  xv6,  esse  despacho  
ocorre  em  devintr  (kernel/trap.c:178).
 Muitos  drivers  de  dispositivo  executam  código  em  dois  contextos:  uma  metade  superior,  que  é  executada  na  thread  
do  kernel  de  um  processo,  e  uma  metade  inferior,  que  é  executada  em  tempo  de  interrupção.  A  metade  superior  é  
chamada  por  meio  de  chamadas  de  sistema,  como  leitura  e  gravação,  que  exigem  que  o  dispositivo  execute  E/S.  Esse  
código  pode  solicitar  ao  hardware  que  inicie  uma  operação  (por  exemplo,  solicitar  ao  disco  que  leia  um  bloco);  em  
seguida,  o  código  aguarda  a  conclusão  da  operação.  Por  fim,  o  dispositivo  conclui  a  operação  e  gera  uma  interrupção.  
O  manipulador  de  interrupção  do  driver,  atuando  como  a  metade  inferior,  descobre  qual  operação  foi  concluída,  ativa  
um  processo  em  espera,  se  apropriado,  e  informa  ao  hardware  para  iniciar  a  execução  de  qualquer  próxima  operação  em  espera.
 5.1  Código:  Entrada  do  console
 O  driver  do  console  (kernel/console.c)  é  uma  ilustração  simples  da  estrutura  do  driver.  O  driver  do  console  aceita  
caracteres  digitados  por  um  humano,  por  meio  do  hardware  de  porta  serial  UART  conectado  ao  RISC-V.
 O  driver  do  console  acumula  uma  linha  de  entrada  por  vez,  processando  caracteres  de  entrada  especiais,  como  
backspace  e  control-u.  Processos  do  usuário,  como  o  shell,  usam  a  chamada  de  sistema  read  para  buscar  linhas  de  
entrada  do  console.  Quando  você  digita  uma  entrada  para  o  xv6  no  QEMU,  suas  teclas  digitadas  são  enviadas  ao  xv6  
por  meio  do  hardware  UART  simulado  do  QEMU.
 O  hardware  UART  com  o  qual  o  driver  se  comunica  é  um  chip  16550  [13]  emulado  pelo  QEMU.  Em  um  computador  
real,  um  16550  gerenciaria  um  link  serial  RS232  conectado  a  um  terminal  ou  outro  computador.
 Ao  executar  o  QEMU,  ele  é  conectado  ao  seu  teclado  e  monitor.
 O  hardware  UART  parece,  para  o  software,  um  conjunto  de  registradores  de  controle  mapeados  na  memória.
 53
Machine Translated by Google
 Ou  seja,  existem  alguns  endereços  físicos  que  o  hardware  RISC-V  conecta  ao  dispositivo  UART,  de  modo  que  
cargas  e  armazenamentos  interagem  com  o  hardware  do  dispositivo  em  vez  da  RAM.  Os  endereços  mapeados  
na  memória  para  a  UART  começam  em  0x10000000,  ou  UART0  (kernel/memlayout.h:21).  Existem  vários  
registradores  de  controle  UART,  cada  um  com  a  largura  de  um  byte.  Seus  deslocamentos  em  relação  à  UART0  
são  definidos  em  (kernel/uart.c:22).  Por  exemplo,  o  registrador  LSR  contém  bits  que  indicam  se  os  caracteres  
de  entrada  estão  aguardando  para  serem  lidos  pelo  software.  Esses  caracteres  (se  houver)  estão  disponíveis  
para  leitura  no  registrador  RHR.  Cada  vez  que  um  caractere  é  lido,  o  hardware  da  UART  o  exclui  de  um  FIFO  
interno  de  caracteres  em  espera  e  limpa  o  bit  "pronto"  no  LSR  quando  o  FIFO  está  vazio.  O  hardware  de  
transmissão  da  UART  é  amplamente  independente  do  hardware  de  recepção;  se  o  software  escreve  um  byte  no  
THR,  a  UART  transmite  esse  byte.
 As  principais  chamadas  do  Xv6  consoleinit  (kernel/console.c:182)  para  inicializar  o  hardware  da  UART.  Este  
código  configura  a  UART  para  gerar  uma  interrupção  de  recebimento  quando  a  UART  recebe  cada  byte  de  
entrada  e  uma  interrupção  de  transmissão  completa  sempre  que  a  UART  termina  de  enviar  um  byte  de  saída  
(kernel/uart.c:53).
 O  shell  xv6  lê  o  console  por  meio  de  um  descritor  de  arquivo  aberto  por  init.c  (usuário/init.c:19).
 As  chamadas  para  o  sistema  read  passam  pelo  kernel  até  o  consoleread  (kernel/  console.c:80).  consoleread  
aguarda  a  chegada  da  entrada  (por  meio  de  interrupções)  e  seu  armazenamento  em  buffer  em  cons.buf,  copia  
a  entrada  para  o  espaço  do  usuário  e  (após  a  chegada  de  uma  linha  inteira)  retorna  ao  processo  do  usuário.  Se  
o  usuário  ainda  não  tiver  digitado  uma  linha  inteira,  qualquer  processo  de  leitura  aguardará  na  chamada  sleep  
(kernel/  console.c:96).  (O  Capítulo  7  explica  os  detalhes  do  sono).
 Quando  o  usuário  digita  um  caractere,  o  hardware  UART  solicita  ao  RISC-V  que  lance  uma  interrupção,  que  
ativa  o  manipulador  de  traps  do  xv6.  O  manipulador  de  traps  chama  devintr  (kernel/trap.c:178).  que  analisa  o  
registrador  RISC-V  Scause  para  descobrir  se  a  interrupção  é  de  um  dispositivo  externo.  Em  seguida,  solicita  a  
uma  unidade  de  hardware  chamada  PLIC  [3]  que  informe  qual  dispositivo  interrompeu  (kernel/trap.c:187).  Se  
fosse  o  UART,  devintr  chamaria  uartintr.
 uartintr  (kernel/uart.c:176)  lê  quaisquer  caracteres  de  entrada  em  espera  do  hardware  UART  e  os  entrega  
ao  consoleintr  (kernel/console.c:136);  Ele  não  espera  por  caracteres,  pois  entradas  futuras  gerarão  uma  nova  
interrupção.  A  função  do  consoleintr  é  acumular  caracteres  de  entrada  em  cons.buf  até  que  uma  linha  inteira  
chegue.  O  consoleintr  trata  o  backspace  e  alguns  outros  caracteres  de  forma  especial.  Quando  uma  nova  linha  
chega,  o  consoleintr  desperta  um  consoleread  em  espera  (se  houver).
 Uma  vez  ativado,  o  consoleread  observará  uma  linha  completa  em  cons.buf,  a  copiará  para  o  espaço  do  usuário  e  retornará  (por  
meio  do  mecanismo  de  chamada  do  sistema)  para  o  espaço  do  usuário.
 5.2  Código:  Saída  do  console
 Uma  chamada  de  sistema  de  gravação  em  um  descritor  de  arquivo  conectado  ao  console  eventualmente  
chega  em  uartputc  (kernel/uart.c:87).  O  driver  do  dispositivo  mantém  um  buffer  de  saída  (uart_tx_buf)  
para  que  os  processos  de  escrita  não  precisem  esperar  o  UART  terminar  de  enviar;  em  vez  disso,  o  
uartputc  anexa  cada  caractere  ao  buffer,  chama  o  uartstart  para  iniciar  a  transmissão  do  dispositivo  (se  
ainda  não  estiver)  e  retorna.  A  única  situação  em  que  o  uartputc  aguarda  é  se  o  buffer  já  estiver  cheio.
 Cada  vez  que  a  UART  termina  de  enviar  um  byte,  ela  gera  uma  interrupção.  uartintr  chama  uartstart,
 54
Machine Translated by Google
 que  verifica  se  o  dispositivo  realmente  concluiu  o  envio  e  entrega  ao  dispositivo  o  próximo  caractere  de  saída  armazenado  em  
buffer.  Assim,  se  um  processo  gravar  vários  bytes  no  console,  normalmente  o  primeiro  byte  será  enviado  pela  chamada  de  
uartputc  para  uartstart,  e  os  bytes  restantes  armazenados  em  buffer  serão  enviados  por  chamadas  de  uartstart  de  uartintr  à  
medida  que  as  interrupções  de  transmissão  completa  chegam.
 Um  padrão  geral  a  ser  observado  é  o  desacoplamento  da  atividade  do  dispositivo  da  atividade  do  processo  por  meio  de  
buffer  e  interrupções.  O  driver  do  console  pode  processar  a  entrada  mesmo  quando  nenhum  processo  está  aguardando  para  
lê-la;  uma  leitura  subsequente  verá  a  entrada.  Da  mesma  forma,  os  processos  podem  enviar  a  saída  sem  precisar  esperar  
pelo  dispositivo.  Esse  desacoplamento  pode  aumentar  o  desempenho,  permitindo  que  os  processos  sejam  executados  
simultaneamente  com  a  E/S  do  dispositivo,  e  é  particularmente  importante  quando  o  dispositivo  está  lento  (como  no  caso  da  
UART)  ou  precisa  de  atenção  imediata  (como  no  caso  de  caracteres  digitados  ecoando).  Essa  ideia  às  vezes  é  chamada  de  E/S.
 simultaneidade.
 5.3  Concorrência  em  drivers
 Você  deve  ter  notado  chamadas  para  acquire  em  consoleread  e  em  consoleintr.  Essas  chamadas  adquirem  um  bloqueio,  que  
protege  as  estruturas  de  dados  do  driver  de  console  contra  acesso  concorrente.  Há  três  perigos  de  simultaneidade  aqui:  dois  
processos  em  CPUs  diferentes  podem  chamar  consoleread  ao  mesmo  tempo;  o  hardware  pode  solicitar  que  uma  CPU  envie  
uma  interrupção  de  console  (na  verdade,  UART)  enquanto  essa  CPU  já  estiver  em  execução  dentro  de  consoleread;  e  o  
hardware  pode  enviar  uma  interrupção  de  console  em  uma  CPU  diferente  enquanto  consoleread  estiver  em  execução.  Esses  
perigos  podem  resultar  em  corridas  ou  deadlocks.  O  Capítulo  6  explora  esses  problemas  e  como  os  bloqueios  podem  resolvê
los.
 Outra  maneira  pela  qual  a  concorrência  requer  cuidado  nos  drivers  é  que  um  processo  pode  estar  aguardando  a  entrada  
de  um  dispositivo,  mas  a  chegada  da  sinalização  de  interrupção  da  entrada  pode  ocorrer  quando  um  processo  diferente  (ou  
nenhum  processo)  estiver  em  execução.  Assim,  os  manipuladores  de  interrupção  não  podem  pensar  sobre  o  processo  ou  
código  que  interromperam.  Por  exemplo,  um  manipulador  de  interrupção  não  pode  chamar  copyout  com  segurança  com  a  
tabela  de  páginas  do  processo  atual.  Os  manipuladores  de  interrupção  normalmente  realizam  relativamente  pouco  trabalho  
(por  exemplo,  apenas  copiam  os  dados  de  entrada  para  um  buffer)  e  ativam  o  código  da  metade  superior  para  fazer  o  resto.
 5.4  Interrupções  do  temporizador
 O  Xv6  utiliza  interrupções  de  temporizador  para  manter  seu  relógio  e  permitir  a  alternância  entre  processos  computacionais;  
as  chamadas  yield  em  usertrap  e  kerneltrap  causam  essa  alternância.  As  interrupções  de  temporizador  vêm  do  hardware  de  
relógio  conectado  a  cada  CPU  RISC-V.  O  Xv6  programa  esse  hardware  de  relógio  para  interromper  cada  CPU  periodicamente.
 O  RISC-V  exige  que  as  interrupções  de  timer  sejam  executadas  no  modo  máquina,  não  no  modo  supervisor.  O  modo  
máquina  RISC-V  executa  sem  paginação  e  com  um  conjunto  separado  de  registradores  de  controle,  portanto,  não  é  prático  
executar  código  kernel  xv6  comum  no  modo  máquina.  Como  resultado,  o  xv6  lida  com  interrupções  de  timer  completamente  
separadamente  do  mecanismo  de  captura  descrito  acima.
 O  código  executado  no  modo  de  máquina  em  start.c,  antes  do  main,  é  configurado  para  receber  interrupções  do  
temporizador  (kernel/start.c:63).  Parte  do  trabalho  consiste  em  programar  o  hardware  CLINT  (interruptor  local  do  núcleo)  para  
gerar  uma  interrupção  após  um  determinado  atraso.  Outra  parte  consiste  em  configurar  uma  área  de  risco,  análoga  à
 55
Machine Translated by Google
 trapframe,  para  ajudar  o  manipulador  de  interrupção  do  temporizador  a  salvar  registradores  e  o  endereço  dos  registradores  CLINT.
 Por  fim,  start  define  mtvec  como  timervec  e  habilita  interrupções  de  timer.
 Uma  interrupção  de  temporizador  pode  ocorrer  a  qualquer  momento  durante  a  execução  de  código  do  usuário  ou  do  kernel;  
não  há  como  o  kernel  desabilitar  interrupções  de  temporizador  durante  operações  críticas.  Portanto,  o  manipulador  de  interrupções  
de  temporizador  deve  executar  sua  função  de  forma  a  garantir  que  não  perturbe  o  código  do  kernel  interrompido.  A  estratégia  básica  
é  que  o  manipulador  solicite  ao  RISC-V  que  gere  uma  "interrupção  de  software"  e  retorne  imediatamente.  O  RISC-V  entrega  
interrupções  de  software  ao  kernel  com  o  mecanismo  de  captura  comum  e  permite  que  o  kernel  as  desabilite.  O  código  para  lidar  
com  a  interrupção  de  software  gerada  por  uma  interrupção  de  temporizador  pode  ser  visto  em  devintr  (kernel/trap.c:205).
 O  manipulador  de  interrupção  do  temporizador  no  modo  máquina  é  timervec  (kernel/kernelvec.S:95).  Ele  salva  alguns  
registradores  na  área  de  rascunho  preparada  por  start,  informa  ao  CLINT  quando  gerar  a  próxima  interrupção  do  
temporizador,  solicita  ao  RISC-V  que  lance  uma  interrupção  de  software,  restaura  os  registradores  e  retorna.  Não  há  
código  C  no  manipulador  de  interrupção  do  temporizador.
 5.5  Mundo  real
 O  Xv6  permite  interrupções  de  dispositivo  e  temporizador  durante  a  execução  no  kernel,  bem  como  durante  a  execução  de  
programas  do  usuário.  As  interrupções  de  temporizador  forçam  uma  troca  de  thread  (uma  chamada  para  yield)  a  partir  do  
manipulador  de  interrupções  de  temporizador,  mesmo  durante  a  execução  no  kernel.  A  capacidade  de  dividir  o  tempo  da  CPU  de  
forma  justa  entre  as  threads  do  kernel  é  útil  se  as  threads  do  kernel  às  vezes  gastam  muito  tempo  computando,  sem  retornar  ao  
espaço  do  usuário.  No  entanto,  a  necessidade  de  o  código  do  kernel  estar  ciente  de  que  pode  ser  suspenso  (devido  a  uma  
interrupção  de  temporizador)  e  posteriormente  retomado  em  uma  CPU  diferente  é  a  fonte  de  alguma  complexidade  no  xv6  (consulte  
a  Seção  6.6).  O  kernel  poderia  ser  um  pouco  mais  simples  se  as  interrupções  de  dispositivo  e  temporizador  ocorressem  apenas  
durante  a  execução  do  código  do  usuário.
 Oferecer  suporte  a  todos  os  dispositivos  em  um  computador  típico  em  todo  o  seu  esplendor  dá  muito  trabalho,  pois  
há  muitos  dispositivos,  os  dispositivos  têm  muitos  recursos  e  o  protocolo  entre  o  dispositivo  e  o  driver  pode  ser  complexo  
e  mal  documentado.  Em  muitos  sistemas  operacionais,  os  drivers  ocupam  mais  código  do  que  o  kernel  principal.
 O  driver  UART  recupera  dados  um  byte  por  vez,  lendo  os  registradores  de  controle  UART;  esse  padrão  é  chamado  de  E/S  
programada,  pois  o  software  controla  a  movimentação  dos  dados.  A  E/S  programada  é  simples,  mas  lenta  demais  para  ser  usada  
em  altas  taxas  de  dados.  Dispositivos  que  precisam  mover  grandes  quantidades  de  dados  em  alta  velocidade  normalmente  usam  
acesso  direto  à  memória  (DMA).  O  hardware  do  dispositivo  DMA  grava  os  dados  de  entrada  diretamente  na  RAM  e  lê  os  dados  de  
saída  da  RAM.  Dispositivos  de  disco  e  rede  modernos  usam  DMA.
 Um  driver  para  um  dispositivo  DMA  prepararia  dados  na  RAM  e,  então,  usaria  uma  única  gravação  em  um  registro  de  controle  para  
dizer  ao  dispositivo  para  processar  os  dados  preparados.
 Interrupções  fazem  sentido  quando  um  dispositivo  precisa  de  atenção  em  momentos  imprevisíveis  e  não  com  muita  frequência.
 Mas  as  interrupções  têm  alta  sobrecarga  de  CPU.  Assim,  dispositivos  de  alta  velocidade,  como  controladores  de  rede  e  de  disco,  
usam  truques  que  reduzem  a  necessidade  de  interrupções.  Um  truque  é  gerar  uma  única  interrupção  para  um  lote  inteiro  de  
solicitações  de  entrada  ou  saída.  Outro  truque  é  o  driver  desabilitar  completamente  as  interrupções  e  verificar  o  dispositivo  
periodicamente  para  ver  se  ele  precisa  de  atenção.  Essa  técnica  é  chamada  de  polling.  O  polling  faz  sentido  se  o  dispositivo  executa  
operações  muito  rapidamente,  mas  desperdiça  tempo  de  CPU  se  o  dispositivo  estiver  ocioso  na  maior  parte  do  tempo.  Alguns  
drivers  alternam  dinamicamente  entre  polling  e  interrupções.
 56
Machine Translated by Google
 dependendo  da  carga  atual  do  dispositivo.
 O  driver  UART  copia  os  dados  recebidos  primeiro  para  um  buffer  no  kernel  e  depois  para  o  espaço  do  usuário.
 Isso  faz  sentido  em  baixas  taxas  de  dados,  mas  essa  cópia  dupla  pode  reduzir  significativamente  o  desempenho  de  
dispositivos  que  geram  ou  consomem  dados  muito  rapidamente.  Alguns  sistemas  operacionais  conseguem  mover  dados  
diretamente  entre  os  buffers  do  espaço  do  usuário  e  o  hardware  do  dispositivo,  geralmente  com  DMA.
 Conforme  mencionado  no  Capítulo  1,  o  console  aparece  para  os  aplicativos  como  um  arquivo  comum,  e  os  aplicativos  
leem  a  entrada  e  escrevem  a  saída  usando  as  chamadas  de  sistema  de  leitura  e  escrita .  Os  aplicativos  podem  querer  
controlar  aspectos  de  um  dispositivo  que  não  podem  ser  expressos  pelas  chamadas  de  sistema  de  arquivos  padrão  (por  
exemplo,  habilitar/desabilitar  o  buffer  de  linha  no  driver  do  console).  Os  sistemas  operacionais  Unix  suportam  a  chamada  de  
sistema  ioctl  para  esses  casos.
 Alguns  usos  de  computadores  exigem  que  o  sistema  responda  em  um  tempo  limitado.  Por  exemplo,  em  sistemas  
críticos  de  segurança,  perder  um  prazo  pode  levar  a  desastres.  O  Xv6  não  é  adequado  para  configurações  de  tempo  real  
rígido.  Sistemas  operacionais  para  tempo  real  rígido  tendem  a  usar  bibliotecas  que  se  conectam  ao  aplicativo  de  forma  a  
permitir  uma  análise  para  determinar  o  pior  tempo  de  resposta.  O  Xv6  também  não  é  adequado  para  aplicativos  de  tempo  
real  flexível,  quando  perder  um  prazo  ocasionalmente  é  aceitável,  porque  o  escalonador  do  Xv6  é  muito  simplista  e  possui  
um  caminho  de  código  do  kernel  onde  as  interrupções  são  desabilitadas  por  um  longo  período.

\chapter{Summary}
\label{CH:SUM}

Este texto apresentou os conceitos fundamentais de sistemas operacionais por meio de um estudo detalhado, linha por linha, de um único sistema: o xv6. Algumas linhas de código representam a essência de conceitos fundamentais de sistemas operacionais—como troca de contexto, a fronteira entre o espaço do usuário e o do kernel, e mecanismos de bloqueio—e são cruciais para entender como um sistema operacional funciona. Outras linhas servem como exemplos ilustrativos de como implementar certas ideias, e poderiam ser desenvolvidas de maneiras diferentes—por exemplo, com algoritmos de escalonamento mais eficientes, estruturas de dados em disco mais eficazes, ou mecanismos de registro (logging) aprimorados para permitir transações concorrentes.

Todos os conceitos foram explorados sob a perspectiva de uma interface de chamadas de sistema específica e amplamente consagrada: a interface Unix. Embora a implementação tenha sido específica ao xv6, as ideias fundamentais se aplicam amplamente ao projeto e à compreensão de diversos sistemas operacionais modernos.


{
% The following prevents latex from splitting a bibliography entry with a page
% break
\interlinepenalty=10000
% Since we're using natbib in numbers mode, we don't need plainnat,
% which exists to feed authors and years back in to natbib.  As a
% result, it complains about entries without years, which we don't
% care about.
%\bibliographystyle{plainnat}
\bibliographystyle{plain}
\bibliography{book}
}

\printindex

\end{document}
